\chapter{RTL、Verilog 与验证}

\begin{keyidea}
RTL 的目标不是描述一串顺序动作，而是描述寄存器、组合网络及其时钟边界。本章训练把 Verilog 读回硬件，并区分仿真、综合和真实时序。
\end{keyidea}

\begin{longtable}{L{0.22\linewidth}L{0.4\linewidth}L{0.28\linewidth}}
\toprule
Verilog 写法 & 对应的硬件 & 误读为程序时的危险 \\
\midrule
\code{assign} / \code{always\_comb} & 组合逻辑网络 & 以为有执行顺序 \\\tablerowrule
\code{always\_ff @(posedge clk)} & 一组触发器 & 以为“每一步都执行一遍” \\\tablerowrule
非阻塞赋值 \code{<=} & 触发器输入端 & 以为立即生效 \\\tablerowrule
\code{case} 不写满 & 可能推断出锁存器 & 以为“没写就不变”是免费的 \\
\bottomrule
\end{longtable}

\begin{keyidea}
读 Verilog 的三问：这一行产生了什么硬件？这个信号是组合路径还是触发器？综合器在哪些写法下会推断出我不想要的东西？能回答这三问，就能读懂绝大多数工程 RTL。
\end{keyidea}

\section{一、从电路图到 RTL：两种描述同一个硬件}

\topic{module 是带端口的电路框图}先给结论：\code{module} 不是程序里的函数或类，它是一张有名字的电路图。\code{module} 与 \code{endmodule} 之间的全部文字，描述的是这张图里有哪些元件、哪些线连到图的边缘。图的边缘——这张电路图与外界交换信号的位置——由端口表给出：\code{input} 是进图的引脚，只能被外界驱动；\code{output} 是出图的引脚，必须由图内某个元件驱动。方向声明的本质是约定驱动权：一根线在同一时刻只能有一个驱动者，这和“数字硬件基础”部分总线规则是同一条纪律。

图内部的中间节点用 \code{wire}（线网）声明。它对应面包板上两个芯片引脚之间那根独立的跳线：没有自己的值，值永远来自驱动它的那一端。另一类声明 \code{reg}（SystemVerilog 中建议写成 \code{logic}）表示“可以被 always 块驱动的信号”——这个名字是历史遗留的误导。它叫“寄存器”，但不一定是寄存器：信号是不是触发器，取决于驱动它的 always 块是时序块还是组合块。组合块驱动的 \code{logic} 综合出来依然是一团门，一个触发器都没有。“时序逻辑的 Verilog 写法”一节写时序块时会把这件事讲透，这里先记住：声明的是“驱动资格”，不是“存储行为”。

逐项对应关系可以列成一张表。这张表是本节的索引，值得记熟：

\begin{longtable}{L{0.22\linewidth}L{0.32\linewidth}L{0.32\linewidth}}
\toprule
Verilog 元素 & 电路图上的对象 & 关键问题 \\
\midrule
\code{module}/\code{endmodule} & 一张电路图的图名与边界 & 这张图对外暴露哪些引脚 \\\tablerowrule
\code{input} & 进图的引脚，只能被外界驱动 & 谁驱动它；断线时默认值是什么 \\\tablerowrule
\code{output} & 出图的引脚，默认是线网型 & 图内哪个元件驱动它 \\\tablerowrule
\code{wire} & 图内一根连线（或一束线） & 两端各连什么；是否有且仅有一个驱动者 \\\tablerowrule
\code{reg}/\code{logic} & 可被 always 块驱动的节点 & 驱动它的是时序块还是组合块 \\\tablerowrule
\code{parameter} & 图纸上的标注常数，不参与连接 & 改标注会不会改变接线方式 \\
\bottomrule
\end{longtable}

最小的完整例子是一个两输入 NAND，对应 74HC00 里四路门中的一路：

\begin{lstlisting}[language=Verilog]
module nand2 (
    input  wire a,   // 引脚：进图
    input  wire b,   // 引脚：进图
    output wire y    // 引脚：出图
);
    assign y = ~(a & b);  // 一个 NAND 门，输出焊在 y 上
endmodule
\end{lstlisting}

逐行读回硬件：第一行给出图名 \code{nand2}；端口表画出三根连到图边缘的线，方向标得一清二楚；块内唯一的 \code{assign} 放一个 NAND 门，两个输入端分别焊在 \code{a}、\code{b} 上，输出端焊在 \code{y} 上；\code{endmodule} 收拢图纸边界。这几行文字就是一张完整电路图，没有一句“会被执行的代码”。

上面这段逐行解读可以直接画成对照图：左边是 \code{nand2} 的这几行代码，右边是它描述的电路图，虚线把每行代码连到它对应的图元。

\begin{pdfdiagram}
  % 左列：代码行
  \node[anchor=west, font=\scriptsize\ttfamily, text=BookInk] (c1) at (-5.4,3.4) {module nand2 (};
  \node[anchor=west, font=\scriptsize\ttfamily, text=BookInk] (c2) at (-5.4,2.6) {input wire a, b,};
  \node[anchor=west, font=\scriptsize\ttfamily, text=BookInk] (c3) at (-5.4,1.8) {output wire y );};
  \node[anchor=west, font=\scriptsize\ttfamily, text=BookInk] (c4) at (-5.4,1.0) {assign y = \textasciitilde(a \& b);};
  \node[anchor=west, font=\scriptsize\ttfamily, text=BookInk] (c5) at (-5.4,0.2) {endmodule};
  % 右：module 图纸边界与内部元件
  \draw[draw=BookNavy, dashed, rounded corners=3pt, line width=0.5pt] (-0.7,-0.5) rectangle (4.8,3.1);
  \node[hw label] at (3.55,2.82) {电路图 nand2};
  \node[hw compute, minimum width=1.7cm, minimum height=1.1cm] (nand) at (1.9,1.4) {NAND 门};
  \draw[hw bus] (-0.7,2.6) -- (0.5,2.6) |- (1.05,1.7);
  \draw[hw bus] (-0.7,1.0) -- (0.65,1.0) |- (1.05,1.1);
  \draw[hw bus] (nand.east) -- (4.8,1.4);
  \fill[BookInk] (-0.7,2.6) circle (0.045);
  \fill[BookInk] (-0.7,1.0) circle (0.045);
  \fill[BookInk] (4.8,1.4) circle (0.045);
  \node[hw label] at (-0.42,2.42) {a};
  \node[hw label] at (-0.42,0.82) {b};
  \node[hw label] at (4.35,1.58) {y};
  \node[hw label] at (0.45,2.85) {wire = 跳线};
  % 对应关系（虚线，不是电路连接）
  \draw[hw return, dashed] (c1.east) -- (-2.0,3.4) -- (-2.0,3.55) -- (1.2,3.55) -- (1.2,3.1);
  \node[hw label, above] at (-0.4,3.55) {module = 图名与边界};
  \draw[hw return, dashed] (c2.east) -- (-0.78,2.6);
  \node[hw label, above] at (-2.1,2.6) {端口 = 引脚};
  \draw[hw return, dashed] (c3.east) -- (-2.4,1.8) -- (-2.4,-0.85) -- (5.15,-0.85) -- (5.15,1.4) -- (4.88,1.4);
  \node[hw label, anchor=west] at (-2.25,-1.1) {output = 出图引脚};
  \draw[hw return, dashed] (c4.east) -- (-1.4,1.0) -- (-1.4,0.3) -- (1.9,0.3) -- (1.9,0.82);
  \node[hw label, above] at (0.6,0.3) {assign = 门与连线};
  \draw[hw return, dashed] (c5.east) -- (-3.0,0.2) -- (-3.0,-1.35) -- (2.4,-1.35) -- (2.4,-0.5);
  \node[hw label, below] at (-0.3,-1.35) {endmodule = 收拢图纸边界};
\end{pdfdiagram}
\diagramnote{这几行代码与电路图元的逐项对应（虚线为对应关系，不是电路连接）：\code{module}/\code{endmodule} 圈出图纸边界，端口表给出进出图的三根引脚，\code{wire} 是引脚到门之间的跳线，\code{assign} 把一只 NAND 门接在 a、b 与 y 之间。读任何 module 都用这个顺序：先看引脚，再数跳线，最后对每个门对号入座。}

端口与驱动的关系可以浓缩成一张速查表，读任何 module 时拿来对照：

\begin{longtable}{L{0.2\linewidth}L{0.32\linewidth}L{0.32\linewidth}}
\toprule
对象 & 谁来驱动 & 误用的后果 \\
\midrule
\code{input} & 图外，经端口送入 & 图内给它赋值：方向冲突 \\\tablerowrule
\code{output}（线网型） & 图内唯一一个驱动者 & 两处驱动：输出对短路 \\\tablerowrule
\code{output reg}/\code{logic} & 图内某一个 always 块 & 多个块驱动同一信号：多驱动错误 \\\tablerowrule
\code{inout} & 双向，靠三态结构让出 & 让出不及时：总线争用，见存储与设备事务相关章节 \\\tablerowrule
内部 \code{wire} & 图内一处 \code{assign} 或实例输出 & 无驱动：悬空，电平不确定 \\
\bottomrule
\end{longtable}

由此得到一个读 module 的固定顺序：先看端口表，弄清这张图与外界交换哪些信号、各进各出；再看内部声明，数出图里有几根中间跳线、几个可被块驱动的节点；最后逐行看驱动关系，把每个门、每根线对号入座。端口表是图纸的图例，先读图例再看图——这个顺序到“从 Verilog 读回硬件：综合边界”一节读大设计时会省大量回头路。

边界上说三件事。第一，module 之间不能嵌套定义，一个 module 是一张独立的图；图与图发生关系的唯一方式是实例化，见本节末尾。第二，端口方向一旦声明就固定了驱动关系，试图在图内驱动 \code{input}，或在两个地方同时驱动同一根 \code{output}，都是连线冲突，与面包板上把两个输出脚对短路同罪。第三，SystemVerilog 的 \code{logic} 可以同时接受 \code{assign} 驱动或单个 always 块驱动，能替换绝大多数 \code{wire}/\code{reg}；本书沿用两种写法并存的工程习惯，读到时只需问“谁驱动它”，不必纠结声明用的是哪个关键字。

顺带认识一种老写法。早期的 Verilog 把端口表只写成一串名字，方向和宽度在块内另行声明：

\begin{lstlisting}[language=Verilog]
module nand2_old (a, b, y);   // 端口表只列名字
    input  a;               // 方向与宽度在块内声明
    input  b;
    output y;

    assign y = ~(a & b);
endmodule
\end{lstlisting}

两种写法描述的是同一张图，硬件没有任何差别；老写法只是把“图边缘有哪些引脚”拆成两处说。新代码一律用前面那种在端口表内直接声明的 ANSI 式写法，读旧代码时认出这是同一件事的两种记法即可。

最后看 \code{parameter}。它不是信号，不占任何一根线，是图纸上的标注常数——像原理图空白处写的“本板按 8 位总线装配”。标注参与描述，改动标注就换了一张不同尺寸的图纸：

\begin{lstlisting}[language=Verilog]
module and_n #(
    parameter WIDTH = 8        // 图纸标注：默认 8 位宽
) (
    input  wire [WIDTH-1:0] a,
    input  wire [WIDTH-1:0] b,
    output wire [WIDTH-1:0] y
);
    assign y = a & b;          // WIDTH 个并排的 AND 门
endmodule
\end{lstlisting}

\code{WIDTH} 决定了这张图里画几个门、每束线有几根；实例化时可以用 \code{\#(.WIDTH(16))} 改标注，得到一张 16 位的图。参数化让一张图纸服务多种规格，省的是画图遍数——每一组参数值仍然综合出一份独立的硬件，这一点与实例化的边界一致。

\topic{assign 是连线，不是赋值}先讲机制。\code{assign y = ~(a \& b);} 在程序语言里读作“把某个值算出来存进 y”，在 RTL 里这句话的含义完全不同：综合器看到它，就把右边表达式对应的组合门网络画出来，把网络的输出端直接焊到左边那根线上。焊完之后没有“执行”这回事——\code{a} 或 \code{b} 一变化，经过门的传播延迟，\code{y} 跟着变。所以它的正式名字叫连续赋值（continuous assignment）：“连续”指的是连接不间断，不是“不停地计算”。

连续赋值有一个程序语言里没有的性质：顺序无关。看下面三行：

\begin{lstlisting}[language=Verilog]
assign carry    = a & b;
assign sum      = a ^ b;
assign overflow = sum & carry;
\end{lstlisting}

三行任意调换次序，综合出来的电路一模一样。这不是特例，而是连续赋值的普遍性质：每一行各自描述一条独立的焊接关系，焊接的先后不影响焊完的板子。仿真器也维持这个性质——任何源信号变化都会触发与之相关的 assign 重新求值（事件驱动），而不是按行号轮流执行。所以读 RTL 时，遇到一组 assign 的正确读法是“这一片有哪些门、各自连到哪”，而不是“先算什么、后算什么”。

这是 Verilog 与程序的第一处分歧，也是把 Verilog 当程序读的读者最容易产生误解的地方。程序里“后面的语句看到前面语句的结果”；RTL 里所有 assign 同时成立，像一张图上所有门同时通电。“数字硬件基础”部分的安全启动控制板通电瞬间，所有门一起开始工作，没有哪个门“先运行”——assign 描述的正是这种并行存在。

再看一个略大的网络，把算术与数据通路相关章节的完整加法器写成门级 assign：

\begin{lstlisting}[language=Verilog]
// 一位全加器：五行 assign 对应六个门，书写次序任意
assign sum  = a ^ b ^ cin;
assign t1   = a & b;
assign t2   = a & cin;
assign t3   = b & cin;
assign cout = t1 | t2 | t3;
\end{lstlisting}

把这五行画成门级电路，每一行对应一到两只门，\code{t1}--\code{t3} 三根中间跳线也有了实体：

\begin{pdfdiagram}
  % 输入分发干线（竖线，圆点为连接点）
  \node[hw label, anchor=east] at (-2.0,3.9) {a};
  \node[hw label, anchor=east] at (-2.0,3.55) {b};
  \node[hw label, anchor=east] at (-2.0,-1.5) {cin};
  \draw[BookMuted, line width=0.62pt] (-1.9,3.9) -- (-0.75,3.9) -- (-0.75,0.85);
  \draw[BookMuted, line width=0.62pt] (-1.9,3.55) -- (-1.05,3.55) -- (-1.05,-0.35);
  \draw[BookMuted, line width=0.62pt] (-1.9,-1.5) -- (-1.35,-1.5) -- (-1.35,2.45);
  % 门（行号标在门内）
  \node[hw compute, minimum width=1.15cm, minimum height=1.0cm] (xor1) at (1.2,3.3) {XOR 行 1};
  \node[hw compute, minimum width=1.15cm, minimum height=1.0cm] (and1) at (1.2,1.7) {AND 行 2};
  \node[hw compute, minimum width=1.15cm, minimum height=1.0cm] (and2) at (1.2,0.5) {AND 行 3};
  \node[hw compute, minimum width=1.15cm, minimum height=1.0cm] (and3) at (1.2,-0.7) {AND 行 4};
  \node[hw compute, minimum width=1.15cm, minimum height=1.0cm] (xor2) at (4.0,3.3) {XOR 行 1};
  \node[hw compute, minimum width=1.15cm, minimum height=1.0cm] (or1) at (4.0,0.5) {OR 行 5};
  % 输入分支（箭头水平进门边）
  \draw[hw bus] (-0.75,3.65) -- (0.6,3.65);
  \draw[hw bus] (-0.75,2.05) -- (0.6,2.05);
  \draw[hw bus] (-0.75,0.85) -- (0.6,0.85);
  \draw[hw bus] (-1.05,2.95) -- (0.6,2.95);
  \draw[hw bus] (-1.05,1.35) -- (0.6,1.35);
  \draw[hw bus] (-1.05,-0.35) -- (0.6,-0.35);
  \draw[hw bus] (-1.35,2.45) -- (3.1,2.45) -- (3.1,2.95) -- (3.4,2.95);
  \draw[hw bus] (-1.35,0.15) -- (0.6,0.15);
  \draw[hw bus] (-1.35,-1.05) -- (0.6,-1.05);
  \fill[BookMuted] (-0.75,3.65) circle (0.04);
  \fill[BookMuted] (-0.75,2.05) circle (0.04);
  \fill[BookMuted] (-0.75,0.85) circle (0.04);
  \fill[BookMuted] (-1.05,2.95) circle (0.04);
  \fill[BookMuted] (-1.05,1.35) circle (0.04);
  \fill[BookMuted] (-1.05,-0.35) circle (0.04);
  \fill[BookMuted] (-1.35,2.45) circle (0.04);
  \fill[BookMuted] (-1.35,0.15) circle (0.04);
  \fill[BookMuted] (-1.35,-1.05) circle (0.04);
  % 门间连线
  \draw[hw bus] (xor1.east) -- (3.4,3.3);
  \draw[hw bus] (and1.east) -- (2.6,1.7) -- (2.6,0.85) -- (3.4,0.85);
  \draw[hw bus] (and2.east) -- (3.4,0.5);
  \draw[hw bus] (and3.east) -- (2.3,-0.7) -- (2.3,0.15) -- (3.4,0.15);
  \draw[hw bus] (xor2.east) -- (5.8,3.3);
  \draw[hw bus] (or1.east) -- (5.8,0.5);
  \node[hw label, anchor=west] at (5.9,3.3) {sum};
  \node[hw label, anchor=west] at (5.9,0.5) {cout};
  \node[hw label, above] at (2.2,1.7) {t1};
  \node[hw label, above] at (2.05,0.5) {t2};
  \node[hw label, below] at (2.0,-0.7) {t3};
\end{pdfdiagram}
\diagramnote{一位全加器的门级电路，行号标在门内：行 1 的两只 XOR 产生 \code{sum}，行 2--4 的三只 AND 与行 5 的 OR 产生 \code{cout}。左侧竖线是输入分发干线，圆点为连接点，十字交叉处不相连。五行任意调换书写次序，综合出的仍是这张图。}

\code{t1}、\code{t2}、\code{t3} 是图内三根中间跳线：不给名字、把 \code{cout} 写成一个大表达式，综合器画出的网络相同；命名的价值在于可读性和调试测点——这与“数字硬件基础”部分给 \code{safe\_chain\_n} 这类中间节点起名的理由完全一样。把 \code{cout} 那行挪到最前面，板子不变；网络上每一点的电平只由此刻的 \code{a}、\code{b}、\code{cin} 决定。两种读法的分野可以归纳成一张表：

\begin{longtable}{L{0.4\linewidth}L{0.44\linewidth}}
\toprule
程序式误读 & 硬件正读 \\
\midrule
“先算 \code{sum}，再算 \code{t1}……” & 六个门同时通电，没有先后 \\\tablerowrule
“\code{t1} 里存着一次运算的结果” & \code{t1} 是一根线，电平随输入即时变化 \\\tablerowrule
“调换两行，结果会不一样” & 焊接次序不影响焊完的板子 \\\tablerowrule
“\code{cout} 用到 \code{t1}，所以它后执行” & \code{cout} 的门接在 \code{t1} 线上，只有传播延迟 \\
\bottomrule
\end{longtable}

\begin{warningbox}
常见误区：“assign 把结果存下来了，后面直接用。”没有任何东西被存下来。assign 的左边是一根线，线上此刻的电平就是右边门网络此刻的输出。想“存”，就得画出存储元件——“状态、时钟与时序逻辑”一章的锁存器、触发器——那对应 always 块和时钟，是后面两节的内容。
\end{warningbox}

边界上只有一条：assign 只能驱动线网型信号，且每根线网只能有一处 assign 驱动（多驱动是另一回事，需要三态总线那种特殊结构，见存储与设备事务相关章节）。想在一个地方根据条件写多条“赋值”，用 always\_comb，那是“组合逻辑的 Verilog 写法”一节的话题。

\topic{位宽、拼接与切片}机制层面，每根信号都有固定宽度，对应一束并排的导线。\code{wire [7:0] data;} 声明八根线，编号从 7 到 0，最高位写在左边——与算术与数据通路相关章节的 bit 编号约定一致，\code{data[7]} 是最高位。切片 \code{data[3:0]} 是从线束里抽出低四根；拼接用一对花括号把几束线并成一束更宽的；常数写成 \code{8'hFF} 这种“位宽-进制-值”三段式，位宽写在最前面。

\begin{lstlisting}[language=Verilog]
wire [7:0] data;
wire [3:0] lo;
wire [7:0] swapped;

assign lo      = data[3:0];                 // 抽出低 4 根线
assign swapped = {data[3:0], data[7:4]};    // 高低半字节换位
\end{lstlisting}

拼接和切片不消耗任何门——它们只是排线方式。这一点必须和运算分清：\code{swapped} 在硅片上的成本是零个门、八根重新排列的连线，而 \code{data + 1} 的成本是一整条加法器进位链。读到拼接时，脑子里应该出现“几束线换了排法”，而不是“发生了什么计算”。花括号里写在前面的信号占高位，就像把几扎电线从左到右排好：\code{data[7:4]} 写在后面，所以换到低位。次序写反，线束就接反——和本书末章的教学计算机项目总线位序接反是同一类事故，不冒烟、不发热，只表现为数值诡异。

把这两行画成位图就更清楚了：切片是抽线，拼接是重排，两处都没有门的影子。

\begin{pdfdiagram}
  % ===== 左：切片 lo = data[3:0] =====
  \node[hw label] at (2.0,4.25) {切片：\code{lo = data[3:0]}};
  \foreach \x/\b in {0/7, 0.5/6, 1/5, 1.5/4, 2/3, 2.5/2, 3/1, 3.5/0} {
    \draw[BookTeal, line width=0.7pt, fill=BookCodeBg] (\x,2.7) rectangle ({\x+0.5},3.5);
    \node at ({\x+0.25},3.1) {\b};
  }
  \node[hw label] at (1.0,2.15) {高 4 位未取};
  \foreach \x/\b in {2/3, 2.5/2, 3/1, 3.5/0} {
    \draw[BookTeal, line width=0.7pt, fill=BookCodeBg] (\x,0.8) rectangle ({\x+0.5},1.6);
    \node at ({\x+0.25},1.2) {\b};
  }
  \node[hw label, anchor=east] at (1.85,1.2) {lo[3:0]};
  \draw[hw bus] (2.25,2.7) -- (2.25,1.6);
  \draw[hw bus] (2.75,2.7) -- (2.75,1.6);
  \draw[hw bus] (3.25,2.7) -- (3.25,1.6);
  \draw[hw bus] (3.75,2.7) -- (3.75,1.6);
  % ===== 右：拼接 swapped = {data[3:0], data[7:4]} =====
  \node[hw label] at (7.9,4.25) {拼接：\code{\{data[3:0], data[7:4]\}}};
  \foreach \x/\b in {5.9/7, 6.4/6, 6.9/5, 7.4/4, 7.9/3, 8.4/2, 8.9/1, 9.4/0} {
    \draw[BookTeal, line width=0.7pt, fill=BookCodeBg] (\x,2.7) rectangle ({\x+0.5},3.5);
    \node at ({\x+0.25},3.1) {\b};
  }
  \node[hw label, anchor=south] at (6.9,3.58) {data[7:4]};
  \node[hw label, anchor=south] at (8.9,3.58) {data[3:0]};
  \foreach \x/\b in {5.9/7, 6.4/6, 6.9/5, 7.4/4, 7.9/3, 8.4/2, 8.9/1, 9.4/0} {
    \draw[BookTeal, line width=0.7pt, fill=white] (\x,0.8) rectangle ({\x+0.5},1.6);
    \node at ({\x+0.25},1.2) {\b};
  }
  \node[hw label, anchor=north] at (6.9,0.62) {= data[3:0]};
  \node[hw label, anchor=north] at (8.9,0.62) {= data[7:4]};
  % 高 4 位换到低位、低 4 位换到高位（规范十字交叉一处）
  \draw[hw bus] (6.6,2.7) -- (6.6,2.25) -- (9.6,2.25) -- (9.6,1.6);
  \node[hw label, above, fill=white, inner sep=1pt] at (8.1,2.27) {高 4 位到低位};
  \draw[hw bus] (9.3,2.7) -- (9.3,1.85) -- (6.95,1.85) -- (6.95,1.6);
  \node[hw label, below, fill=white, inner sep=1pt] at (8.1,1.83) {低 4 位到高位};
\end{pdfdiagram}
\diagramnote{左：切片 \code{data[3:0]} 从八根线里抽出低四根，高四位不参与。右：拼接 \code{\{data[3:0], data[7:4]\}} 把低半字节换到高位、高半字节换到低位，\code{swapped} 的每一位仍来自同一束线。两种操作都只是重新排线——图上没有一只门，硬件成本为零。}

位宽不匹配时，Verilog 的处理规则是机械的：表达式先按参与运算的最大位宽对齐，短的一侧高位补 0；结果赋给左侧时，左侧更宽就继续补 0，左侧更窄就直接丢弃高位。这套规则没有任何“智能”，只有约定。它导致的经典事故有两种：两个 8 位数相加，和需要 9 位，左边只给 8 位，进位被静默丢弃；1 位标志与 8 位数据比较，前者先被补成 8 位再比，比较对象与直觉不符。统计工程实践中 RTL 低级 bug 的来源，位宽类问题常年排在前列，远比“逻辑写反”常见——因为每一行单独看都像是写对了。

两段有病的代码，病都在位宽：

\begin{lstlisting}[language=Verilog]
wire [7:0] a, b;
wire [7:0] sum_bad;
wire [8:0] sum_good;
wire       flag;
wire [7:0] threshold;
wire       too_low;

assign sum_bad  = a + b;            // 病：第 9 位进位被静默丢弃
assign sum_good = {1'b0, a} + {1'b0, b};  // 先扩到 9 位再加，进位有处可去
assign too_low  = (threshold < flag);     // 病：flag 补成 8 位再比，阈值形同虚设
\end{lstlisting}

\code{sum\_bad} 仿真时大部分输入都对，只在和溢出时出错——最折磨人的一类 bug。\code{too\_low} 更隐蔽：\code{flag} 只有 0 和 1 两种取值，比较退化成“\code{threshold} 是否为 0”，阈值设定完全失效。两种病的共同点是编译不报错、每行都“像是对的”，只有逐位推算才能发现。

\begin{longtable}{L{0.24\linewidth}L{0.34\linewidth}L{0.28\linewidth}}
\toprule
写法 & 含义 & 硬件成本 \\
\midrule
\code{wire [7:0] d} & 8 根并排的线 & 8 根连线 \\\tablerowrule
\code{d[3:0]} & 抽出低 4 根 & 0 个门 \\\tablerowrule
\code{d[3:0]} 与 \code{d[7:4]} 拼接 & 两束并成一束，次序即排位 & 0 个门 \\\tablerowrule
\code{4'b0} 与 \code{d[3:0]} 拼接 & 高位补 0 扩展到 8 位 & 0 个门，补位接固定电平 \\\tablerowrule
\code{d + 8'd1} & 8 位加法 & 一条 8 位进位链 \\
\bottomrule
\end{longtable}

写 RTL 的两条纪律由此而来。第一，常数永远写全位宽，\code{8'd0} 而不是裸 \code{0}，让每一行的宽度意图可见。第二，加法结果先想进位去向：要么左边多给一位，要么把左边写成 \code{carry} 与 \code{sum} 的拼接，让进位进入有名信号。算术与数据通路相关章节给 ALU 算标志时就是这么处理的，这里只是把同一件事写成文字。

读别人的代码时，位宽是头号排查对象，按三个检查点过一遍：每个运算符左右两侧的宽度各是多少；每个常数有没有写明位宽；每处拼接的次序与预期的高位排位是否一致。三问都答得出，位宽类 bug 基本无处藏身；答不出，就先停下来把宽度逐个标注在草稿上，再继续往下读。

\topic{案例：安全启动控制板的并排对照}“数字硬件基础”部分用一片 74HC04、一片 74HC00、一片 74HC32 搭过三级报警优先权电路：\code{fire} 优先级最高，\code{fault} 次之，\code{info} 最低，任一时刻只点亮当前最高优先级那盏灯，高优先级请求以反相形式出现在低优先级的乘积项里，从逻辑上把低优先级输出强制为 0。现在把同一张电路图写成 Verilog：

\begin{lstlisting}[language=Verilog]
module alarm_priority (
    input  wire fire,
    input  wire fault,
    input  wire info,
    output wire fire_lamp,
    output wire fault_lamp,
    output wire info_lamp,
    output wire any_alarm
);
    wire fire_n;    // fire 的反相，U1 门 1
    wire any_hi;    // fire OR fault，U3 门 1
    wire any_hi_n;  // any_hi 的反相，U1 门 3

    assign fire_n     = ~fire;
    assign fire_lamp  = fire;               // 直连，不经门
    assign fault_lamp = fault & fire_n;     // U2 门 1 加 U1 门 2
    assign any_hi     = fire | fault;
    assign any_hi_n   = ~any_hi;
    assign info_lamp  = info & any_hi_n;    // U2 门 2 加 U1 门 4
    assign any_alarm  = any_hi | info;      // 校验输出，U3 门 2
endmodule
\end{lstlisting}

代码行与芯片门的对应是一一可数的：

\begin{longtable}{L{0.3\linewidth}L{0.3\linewidth}L{0.26\linewidth}}
\toprule
代码行 & “数字硬件基础”部分对应资源 & 检查点 \\
\midrule
\code{fire\_lamp = fire} & 直连，不经门 & 最高优先级路径最短 \\\tablerowrule
\code{fire\_n = ~fire} & U1（74HC04）门 1 & 反相信号参与屏蔽项 \\\tablerowrule
\code{fault\_lamp = fault \& fire\_n} & U2（74HC00）门 1 加 U1 门 2 & NAND 加反相实现 AND \\\tablerowrule
\code{any\_hi = fire | fault} & U3（74HC32）门 1 & 高优先级汇总 \\\tablerowrule
\code{any\_hi\_n = ~any\_hi} & U1 门 3 & 供 \code{info} 屏蔽项使用 \\\tablerowrule
\code{info\_lamp = info \& any\_hi\_n} & U2 门 2 加 U1 门 4 & 德摩根改写省去三输入 AND \\\tablerowrule
\code{any\_alarm = any\_hi | info} & U3 门 2 & 两条独立路径算同一事实 \\
\bottomrule
\end{longtable}

并排看有两点值得停留。第一，结构完全对应：三片芯片、八个门、若干跳线，变成二十行文字，每个门都能在代码里点到名，\code{any\_alarm} 这条刻意多算的校验路径也原样保留——两条独立路径算同一事实的调试思想，在 RTL 里同样成立。第二，控制精度发生了变化：“数字硬件基础”部分能精确说出“\code{fire} 到 \code{fault\_lamp} 三级门”，因为门是亲手挑的；RTL 把“用哪几个门”的决定权交给了综合器和工艺库，文字里只剩逻辑关系。放弃对门数的精确控制，换取写法的紧凑和可维护——这就是 RTL 相对原理图的定位。想保住精确门数也有办法：门级实例化写法，下一个专题就讲。

还有一些原理图上的要素在 RTL 里根本没有对应行：电源脚、去耦电容、未用门的输入处理、电平余量。它们并没有消失，而是移交给了工艺库、后端工具和板级设计——模块实例化成芯片之后，这些约定仍然成立，只是不再由 RTL 文字表达。读 RTL 时心里要留着这张移交清单，尤其在把 FPGA 引脚接到真实报警灯的时候：驱动电流、上拉下拉这类电气约定，永远要去原理图和数据手册里查，代码里找不到。

反方向的练习同样重要：拿到一段 RTL，把它读回电路图。方法是对每一行 assign 问三个问题——右边表达式是什么门网络，左边线网连到哪里，线网还有没有其他驱动者。三个问题答完，\code{alarm\_priority} 就重新变回八个门、三根内部跳线、一条直连。再用“数字硬件基础”部分的真值表复核：\code{fire}=1 的四行里 \code{fire\_lamp} 恒为 1 且另两灯恒为 0，\code{fault} 只在 \code{fire}=0 时才能点亮自己的灯——代码与真值表逐行相符，这段 RTL 才算真正读完。

\topic{层次化实例化}机制：实例化是把一张画好的电路图整个放进另一张电路图。写 \code{alarm\_priority u\_alarm (...)} 时，\code{alarm\_priority} 是图名，\code{u\_alarm} 是这一次放置的元件编号——就是原理图上的 U1、U2。同一张图可以放很多次，每一次放置都是一份独立的硬件拷贝，像同一型号的芯片买了好几片，各自通电、互不干扰。

端口连接推荐点名法 \code{.端口(信号)}：

\begin{lstlisting}[language=Verilog]
module panel_top (
    input  wire       fire,
    input  wire       fault,
    input  wire       info,
    output wire [2:0] lamps,
    output wire       buzzer
);
    wire any_alarm_w;   // 子图输出到蜂鸣器之间的跳线

    alarm_priority u_alarm (
        .fire       (fire),       // 子图 fire 脚接本图 fire 线
        .fault      (fault),
        .info       (info),
        .fire_lamp  (lamps[2]),   // 三盏灯捆成一束线
        .fault_lamp (lamps[1]),
        .info_lamp  (lamps[0]),
        .any_alarm  (any_alarm_w)
    );

    assign buzzer = any_alarm_w;  // 蜂鸣器由校验输出驱动
endmodule
\end{lstlisting}

\code{.fire(fire)} 读作“把子图的 \code{fire} 引脚焊到本图的 \code{fire} 线上”。点名法里端口次序随意，名字对应即可；三盏灯拼成 \code{lamps[2:0]} 一束线，子图的三个单位输出分别焊到束中各根——拼接在端口连接处同样只是排线。还有一种按端口表顺序连的位置法，端口一增删就错位，工程代码一律用点名法。漏连的端口默认悬空，和真实芯片引脚悬空一样危险：“数字硬件基础”部分讲过输入悬空的电平漂移，综合器对悬空端口通常只给警告，警告必须逐条看完。

大设计按“小图进中图、中图进整图”自底向上组织。处理器核心相关章节的最小 CPU 就是这样分图的：ALU 一张、寄存器一张、控制器一张、存储器一张，顶层图只画它们之间的总线和控制线。Verilog 的 module 层次就是这套图纸体系的文字版。读别人写的大设计，正确顺序也是先找顶层 module 的实例化列表，弄清有哪几张子图、子图之间连了哪些线，再逐张进去看细节——和拿到一台新设备先看整机框图、再拆电路板是同一个动作。

把这套“小图进中图、中图进整图”画出来，module 层次就是图套图：顶层图纸的边界圈住若干张已完成的子图，顶层自己只画子图之间的连线。

\begin{pdfdiagram}
  % 顶层图纸边界
  \draw[draw=BookNavy, dashed, rounded corners=3pt, line width=0.5pt] (-0.9,-2.1) rectangle (8.3,2.5);
  \node[hw label] at (5.55,2.22) {顶层图 cpu\_top：只画总线与控制线};
  % 两张子图（寄存器与存储器子图从略）
  \node[hw compute, minimum width=2.6cm, minimum height=1.3cm] (ualu) at (1.5,0.3) {u\_alu\\子图 alu8};
  \node[hw control, minimum width=2.6cm, minimum height=1.3cm] (uctrl) at (5.6,0.3) {u\_ctrl\\子图 controller};
  % 子图间的数据通路（实线）
  \draw[hw bus] (ualu.east) -- (uctrl.west);
  \node[hw label, above] at (3.55,0.32) {result / flags};
  % 控制线（虚线）：u_ctrl 底部通道折入 u_alu 底部
  \draw[hw return, dashed] (uctrl.south) -- (5.6,-1.35) -- (1.5,-1.35) -- (ualu.south);
  \node[hw label, below] at (3.55,-1.37) {控制线 op};
  % clk 从图外进入，沿顶部通道分发到两子图
  \node[hw label, anchor=east] at (-1.75,1.5) {clk};
  \draw[hw return, dashed] (-1.65,1.5) -- (5.6,1.5);
  \draw[hw return, dashed] (1.5,1.5) -- (ualu.north);
  \draw[hw return, dashed] (5.6,1.5) -- (uctrl.north);
  \fill[BookMuted] (1.5,1.5) circle (0.04);
  \fill[BookMuted] (5.6,1.5) circle (0.04);
  \fill[BookInk] (-0.9,1.5) circle (0.045);
  % 输出端口：穿出图纸边界
  \draw[hw bus] (uctrl.east) -- (8.9,0.3);
  \fill[BookInk] (8.3,0.3) circle (0.045);
  \node[hw label, anchor=west] at (9.0,0.3) {输出};
  \node[hw label] at (6.5,-1.75) {寄存器与存储器子图从略};
\end{pdfdiagram}
\diagramnote{module 层次就是图纸层次：顶层图 \code{cpu\_top} 的虚线边界内放着两张子图——\code{u\_alu} 是 \code{alu8} 的一次放置，\code{u\_ctrl} 是控制器图的一次放置；同一张图可放置多次，每次都是一份独立硬件。子图之间只有两类线：实线的数据通路（result/flags）与虚线的控制线（op、clk），穿过边界的圆点是顶层图的端口。寄存器与存储器两子图从略。}

边界上有一句必须说的话：实例化产生的是真实的第二份硬件，不是“调用”。三个实例就是三倍面积、三倍功耗。软件里“抽个函数复用代码”的直觉在这里要改：RTL 复用模块，省的是画图和验证的工夫，不是硅片。想要“一份硬件分时复用”，那是状态机加数据通路的课题，处理器核心相关章节已经见过——控制器让同一台 ALU 在不同周期干不同的活，那才是硬件意义上的“子程序”。

实例化还有一层用途，正好兑现上一个专题的承诺：门级实例化。把基本门各自写成 module，再全部用实例化拼装，就能让门数精确到个位，恢复“数字硬件基础”部分原理图的控制精度：

\begin{lstlisting}[language=Verilog]
module not1 (input wire a, output wire y);
    assign y = ~a;
endmodule

module or2 (input wire a, input wire b, output wire y);
    assign y = a | b;
endmodule

// 三级报警的门级版本：与“数字硬件基础”部分芯片分配逐门对应
module alarm_priority_gates (
    input  wire fire,
    input  wire fault,
    input  wire info,
    output wire fire_lamp,
    output wire fault_lamp,
    output wire info_lamp,
    output wire any_alarm
);
    wire fire_n, m1, any_hi, any_hi_n, m2;

    not1 u_inv1 (.a(fire),          .y(fire_n));      // U1 门 1
    nand2 u_nand1 (.a(fault),       .b(fire_n),  .y(m1));      // U2 门 1
    not1 u_inv2 (.a(m1),            .y(fault_lamp));  // U1 门 2
    or2  u_or1  (.a(fire),          .b(fault),   .y(any_hi));  // U3 门 1
    not1 u_inv3 (.a(any_hi),        .y(any_hi_n));    // U1 门 3
    nand2 u_nand2 (.a(info),        .b(any_hi_n), .y(m2));     // U2 门 2
    not1 u_inv4 (.a(m2),            .y(info_lamp));   // U1 门 4
    or2  u_or2  (.a(any_hi),        .b(info),    .y(any_alarm)); // U3 门 2

    assign fire_lamp = fire;                          // 直连
endmodule
\end{lstlisting}

每个实例对应“数字硬件基础”部分芯片分配表里的一个门，实例名 \code{u\_nand1} 之类就是元件编号，\code{m1}、\code{m2} 这些中间线名也原样保留。这种写法叫结构级描述：文字里没有任何逻辑表达式，只有元件和连线，读起来像一份网表（门与连线的清单）。它的用途有二：一是教学上证明“module 层次就是图纸层次”到底到了什么程度；二是某些对实现有精确要求的场合（手工摆放的关键路径、复用经过验证的硬核），工程师真的会写到这一层。代价同样明显：逻辑意图被拆散到几十个实例里，“这段代码想干什么”要拼完整张图才能回答。所以日常 RTL 停在 assign 和 always 块这一层，把门级留给需要精确控制的地方——抽象层次的选择本身就是一种工程判断。

收拢本节，五份对应已经齐备：\code{module} 是带端口的电路框图，端口表是图边缘的引脚；\code{assign} 是连线不是赋值，次序无关是它区别于程序的第一处；位宽、拼接、切片说的都是线束的排法，大多数低级 bug 藏在这里；实例化是把一张图放进另一张图，实例名就是元件编号。这套对应不依赖任何工具，纸笔即可验证——这正是“用文字画电路”的含义：每一行都能读回具体硬件，读不回去的行，多半也没写对。

“组合逻辑的 Verilog 写法”一节进入组合逻辑的第二种写法。当分支和中间量多到 \code{assign} 一行写不下时，\code{always\_comb} 登场——但请记住，它只是换了一种描述手段，焊出来的仍然是没有记忆的门网络。

\section{二、组合逻辑的 Verilog 写法}

组合逻辑有两种文字写法：“从电路图到 RTL：两种描述同一个硬件”一节的 \code{assign}，和本节的过程块 \code{always\_comb}。

\topic{assign 与 always\_comb 的分工}机制上，两者完全等价：都描述组合网络，输出都直接连到对应信号上，输入一变输出就跟着变。区别只在描述手段。\code{assign} 一行写一条连接，右边是一个表达式；\code{always\_comb} 的一个块里可以写多行、引入临时变量、使用 \code{if}/\code{case} 分支，综合器读完整个块，按数据依赖把最终结果展开成门网络。写法不同，焊出来的东西是同一类。

同一个 2 选 1 选择器，两种写法并排看：

\begin{lstlisting}[language=Verilog]
assign y = sel ? b : a;

always_comb begin
    if (sel) y = b;
    else     y = a;
end
\end{lstlisting}

综合结果完全相同：每比特一个二选一选择器，\code{sel} 接选择端。既然如此，分工按可读性来：一行能写完的表达式用 \code{assign}；分支多、中间量多、有临时代号的逻辑用 \code{always\_comb}。算术与数据通路相关章节 ALU 那种“候选结果并排算好、再按 \code{op} 挑一个”的结构，用 \code{always\_comb} 写远比嵌套三元运算符清楚，本节末尾会给出完整例子。

\code{always\_comb} 里最容易误读的还是顺序。块内语句从上到下书写，这只是描述顺序：对同一个信号，后面的赋值覆盖前面的，最后一笔生效；对不同信号，赋值之间没有先后，展开成并行存在的门。仿真器按书写顺序执行块内语句（阻塞赋值），配合 \code{always\_comb} 自带的完整敏感表——块内读到的任何信号一变，整块立即重算——最终稳态结果与综合出的门网络一致。两套视角殊途同归，但硬件图景只有一个：块里没有“步骤”，只有连接。把 \code{always\_comb} 读成“每当输入变化就执行一遍的小程序”，就又把 RTL 读回程序了。

把“描述顺序”与“数据流”并排画出来，两种读法的差别就成了图上的差别：左边是给人看的书写次序，右边才是综合器真正关心的依赖关系。

\begin{pdfdiagram}[x=0.85cm]
  % 左：块内书写顺序（只是描述顺序）
  \node[hw label] at (-4.3,2.7) {块内书写顺序：只是描述顺序};
  \node[anchor=west, font=\scriptsize\ttfamily, text=BookInk] at (-6.2,1.9) {t1 = a \& b;};
  \node[anchor=west, font=\scriptsize\ttfamily, text=BookInk] at (-6.2,1.2) {t2 = t1 | c;};
  \node[anchor=west, font=\scriptsize\ttfamily, text=BookInk] at (-6.2,0.5) {y  = t2 \textasciicircum d;};
  \node[hw label, align=center] at (-4.3,-0.5) {三行任意调换次序\\综合结果不变};
  % 右：综合按数据流展开
  \node[hw label] at (2.6,2.7) {综合按数据流展开：只有依赖，没有步骤};
  \node[hw compute, minimum width=1.3cm, minimum height=0.95cm, align=center] (g1) at (0.3,1.2) {AND\\行 1};
  \node[hw compute, minimum width=1.3cm, minimum height=0.95cm, align=center] (g2) at (2.6,1.2) {OR\\行 2};
  \node[hw compute, minimum width=1.3cm, minimum height=0.95cm, align=center] (g3) at (4.9,1.2) {XOR\\行 3};
  \node[hw label, anchor=east] at (-1.15,1.42) {a};
  \node[hw label, anchor=east] at (-1.15,0.98) {b};
  \draw[hw bus] (-1.0,1.42) -- ($(g1.west)+(0,0.22)$);
  \draw[hw bus] (-1.0,0.98) -- ($(g1.west)+(0,-0.22)$);
  \node[hw label] at (2.6,2.2) {c};
  \draw[hw bus] (2.6,2.0) -- (g2.north);
  \node[hw label] at (4.9,2.2) {d};
  \draw[hw bus] (4.9,2.0) -- (g3.north);
  \draw[hw bus] (g1.east) -- (g2.west);
  \node[hw label, above] at (1.45,1.22) {t1};
  \draw[hw bus] (g2.east) -- (g3.west);
  \node[hw label, above] at (3.75,1.22) {t2};
  \draw[hw bus] (g3.east) -- (6.5,1.2);
  \node[hw label, anchor=west] at (6.6,1.2) {y};
  \node[hw label, anchor=north] at (2.6,0.25) {y 只由 a、b、c、d 的当前值决定；t1、t2 是中间跳线};
\end{pdfdiagram}
\diagramnote{同一个 \code{always\_comb} 块的两种看法。左：块内三行有书写先后，那只是描述顺序，任意调换，综合结果不变。右：综合器按数据依赖把三行展开成 AND、OR、XOR 的级联——行号标在门内，\code{t1}、\code{t2} 是门之间的中间跳线。块里没有“步骤”，只有连接；输出只由输入的当前值经传播延迟决定。}

临时变量是 \code{always\_comb} 相对 \code{assign} 多出来的主要表达手段。把“从电路图到 RTL：两种描述同一个硬件”一节的三级报警改写成一个块：

\begin{lstlisting}[language=Verilog]
always_comb begin
    // 临时变量是描述中的代号，综合后是中间跳线
    fire_n     = ~fire;
    any_hi     = fire | fault;
    any_hi_n   = ~any_hi;

    fire_lamp  = fire;
    fault_lamp = fault & fire_n;
    info_lamp  = info & any_hi_n;
    any_alarm  = any_hi | info;
end
\end{lstlisting}

\code{fire\_n}、\code{any\_hi}、\code{any\_hi\_n} 在块内先写后用，读起来有先后，但这只是给人看的书写次序——综合器按数据依赖展开后，与“从电路图到 RTL：两种描述同一个硬件”一节七行 \code{assign} 的版本是同一批门、同一批跳线。块内赋值用阻塞赋值 \code{=}，它的语义是“代号立即更新，后面的语句看到新值”；综合之后没有任何元件对应这个“立即”，它只是描述规则。正是为了守住“组合块用阻塞、时序块用非阻塞”这条纪律，“从 Verilog 读回硬件：综合边界”一节的移位寄存器反例才有判据。

分工可以浓缩成一张速查表：

\begin{longtable}{L{0.4\linewidth}L{0.42\linewidth}}
\toprule
场景 & 推荐写法 \\
\midrule
单行表达式：一根线、一个门网络 & \code{assign} \\\tablerowrule
多分支选择、需要中间代号 & \code{always\_comb} \\\tablerowrule
模块之间穿线、端口直连 & \code{assign} \\\tablerowrule
需要 if/case 表达译码或优先权 & \code{always\_comb} \\
\bottomrule
\end{longtable}

\begin{keyidea}
\code{assign} 与 \code{always\_comb} 的选择标准是表达力，不是功能：表达式一行写得清，用 \code{assign}；需要分支和中间量，用 \code{always\_comb}。两者综合出的都是无记忆的组合网络——判断依据永远是“输入变，输出经传播延迟后跟着变”，没有时钟、没有保持。
\end{keyidea}

\topic{运算符与硬件的对应表}每个运算符背后都是一类具体电路，面积和延迟的量级各不相同。写 RTL 的人心里要有这张表，写完一行就能报出代价：

\begin{longtable}{L{0.16\linewidth}L{0.32\linewidth}L{0.38\linewidth}}
\toprule
运算符 & 综合出的硬件 & 面积与延迟量级（n 位宽） \\
\midrule
\code{a + b} & 串行进位加法器 & n 个全加器；延迟沿进位链约 n 级门 \\\tablerowrule
\code{a - b} & 加补器：加法器加 B 反相、Cin=1 & 与加法器同级，多 n 个可控反相 \\\tablerowrule
\code{a == b} & 逐位 XNOR 加与归约树 & n 个 XNOR，约 $\log n$ 级门延迟 \\\tablerowrule
\code{a < b} & 复用减法路径，看借位或符号位 & 与减法器同级 \\\tablerowrule
\code{sel ? x : y} & 每比特一个 2 选 1 mux & n 个 mux，1 级门延迟 \\\tablerowrule
逐位 \code{\&}、\code{|} & n 个并行 AND/OR 门 & n 个门，1 级门延迟 \\\tablerowrule
归约 \code{\&a} & 与归约树：n 位收成 1 位 & n-1 个 AND，约 $\log n$ 级 \\\tablerowrule
归约异或 & XOR 归约树，即奇偶校验 & n-1 个 XOR，约 $\log n$ 级 \\\tablerowrule
\code{~a} & n 个反相器 & n 个门，1 级门延迟 \\\tablerowrule
\code{a << 1}（常数移位） & 重新排线，空位接固定电平 & 0 个门 \\
\bottomrule
\end{longtable}

三组区分值得单独说。第一组是逐位与归约：\code{a \& b} 中两个 n 位输入对应位各自过门，输出仍是 n 位；归约 \code{\&a} 只有一个操作数，把 n 位收成 1 位，硬件是一棵树。写法只差一个操作数的位置，硬件从“n 个门并排”变成“一棵 n 输入的树”，读代码时先看操作数个数再下结论。第二组是算术与比较：减法不单独造电路，沿用算术与数据通路相关章节的补码方案，B 反相加最低位进位 1，与加法共用进位链；比较 \code{a < b} 通常就复用这条减法路径，看最高位借位或符号位——所以比较器不便宜，它背后藏着一整条加法器。第三组是移位：移常数位只是排线，零成本；移变量位才是桶形移位器，面积随位宽平方增长，不可混为一谈。

链与树的形状差别画出来，比量级公式更直观。

\begin{pdfdiagram}
  % ===== 上：8 级串行进位链 =====
  \node[hw label, anchor=south] at (4.6,3.85) {串行进位链：进位逐级穿过 8 个全加器};
  \node[hw compute, minimum width=0.85cm] (fa0) at (0.575,2.2) {FA};
  \node[hw compute, minimum width=0.85cm] (fa1) at (1.725,2.2) {FA};
  \node[hw compute, minimum width=0.85cm] (fa2) at (2.875,2.2) {FA};
  \node[hw compute, minimum width=0.85cm] (fa3) at (4.025,2.2) {FA};
  \node[hw compute, minimum width=0.85cm] (fa4) at (5.175,2.2) {FA};
  \node[hw compute, minimum width=0.85cm] (fa5) at (6.325,2.2) {FA};
  \node[hw compute, minimum width=0.85cm] (fa6) at (7.475,2.2) {FA};
  \node[hw compute, minimum width=0.85cm] (fa7) at (8.625,2.2) {FA};
  \node[hw label, anchor=east] at (-0.2,2.2) {$C_0$};
  \draw[hw bus] (0.05,2.2) -- (fa0.west);
  \draw[hw bus] (fa0.east) -- (fa1.west);
  \draw[hw bus] (fa1.east) -- (fa2.west);
  \draw[hw bus] (fa2.east) -- (fa3.west);
  \draw[hw bus] (fa3.east) -- (fa4.west);
  \draw[hw bus] (fa4.east) -- (fa5.west);
  \draw[hw bus] (fa5.east) -- (fa6.west);
  \draw[hw bus] (fa6.east) -- (fa7.west);
  \draw[hw bus] (fa7.east) -- (9.3,2.2);
  \node[hw label, anchor=west] at (9.4,2.2) {$C_8$};
  % 每级从上方进一位操作数
  \draw[hw bus] (fa0.north) -- ++(0,0.55);
  \draw[hw bus] (fa1.north) -- ++(0,0.55);
  \draw[hw bus] (fa2.north) -- ++(0,0.55);
  \draw[hw bus] (fa3.north) -- ++(0,0.55);
  \draw[hw bus] (fa4.north) -- ++(0,0.55);
  \draw[hw bus] (fa5.north) -- ++(0,0.55);
  \draw[hw bus] (fa6.north) -- ++(0,0.55);
  \draw[hw bus] (fa7.north) -- ++(0,0.55);
  \node[hw label, anchor=south] at (0.575,3.2) {$a_0,b_0$};
  \node[hw label, anchor=south] at (4.6,3.2) {$\cdots$};
  \node[hw label, anchor=south] at (8.625,3.2) {$a_7,b_7$};
  % 延迟标尺
  \draw[{Stealth[length=1.8mm]}-{Stealth[length=1.8mm]}, BookRed, line width=0.6pt] (0.575,1.5) -- (8.625,1.5);
  \node[hw label, anchor=north] at (4.6,1.38) {8 级门延迟（随 $n$ 线性）};
  % ===== 下：归约树 =====
  \node[hw label, anchor=north] at (4.6,-4.25) {归约树（\&a）：每级减半，8 位收成 1 位};
  % 8 个输入
  \fill[BookMuted] (0.575,0.2) circle (0.045);
  \fill[BookMuted] (1.725,0.2) circle (0.045);
  \fill[BookMuted] (2.875,0.2) circle (0.045);
  \fill[BookMuted] (4.025,0.2) circle (0.045);
  \fill[BookMuted] (5.175,0.2) circle (0.045);
  \fill[BookMuted] (6.325,0.2) circle (0.045);
  \fill[BookMuted] (7.475,0.2) circle (0.045);
  \fill[BookMuted] (8.625,0.2) circle (0.045);
  \node[hw label, anchor=south] at (0.575,0.28) {$x_0$};
  \node[hw label, anchor=south] at (1.725,0.28) {$x_1$};
  \node[hw label, anchor=south] at (2.875,0.28) {$x_2$};
  \node[hw label, anchor=south] at (4.025,0.28) {$x_3$};
  \node[hw label, anchor=south] at (5.175,0.28) {$x_4$};
  \node[hw label, anchor=south] at (6.325,0.28) {$x_5$};
  \node[hw label, anchor=south] at (7.475,0.28) {$x_6$};
  \node[hw label, anchor=south] at (8.625,0.28) {$x_7$};
  % 第一级：4 个与门
  \node[hw compute, minimum width=0.75cm] (g1) at (1.15,-0.9) {\&};
  \node[hw compute, minimum width=0.75cm] (g2) at (3.45,-0.9) {\&};
  \node[hw compute, minimum width=0.75cm] (g3) at (5.75,-0.9) {\&};
  \node[hw compute, minimum width=0.75cm] (g4) at (8.05,-0.9) {\&};
  \draw[hw bus] (0.575,0.2) -- (0.575,-0.9) -- (g1.west);
  \draw[hw bus] (1.725,0.2) -- (1.725,-0.9) -- (g1.east);
  \draw[hw bus] (2.875,0.2) -- (2.875,-0.9) -- (g2.west);
  \draw[hw bus] (4.025,0.2) -- (4.025,-0.9) -- (g2.east);
  \draw[hw bus] (5.175,0.2) -- (5.175,-0.9) -- (g3.west);
  \draw[hw bus] (6.325,0.2) -- (6.325,-0.9) -- (g3.east);
  \draw[hw bus] (7.475,0.2) -- (7.475,-0.9) -- (g4.west);
  \draw[hw bus] (8.625,0.2) -- (8.625,-0.9) -- (g4.east);
  % 第二级：2 个与门
  \node[hw compute, minimum width=0.75cm] (G1) at (2.3,-2.1) {\&};
  \node[hw compute, minimum width=0.75cm] (G2) at (6.9,-2.1) {\&};
  \draw[hw bus] (g1.south) -- (1.15,-2.1) -- (G1.west);
  \draw[hw bus] (g2.south) -- (3.45,-2.1) -- (G1.east);
  \draw[hw bus] (g3.south) -- (5.75,-2.1) -- (G2.west);
  \draw[hw bus] (g4.south) -- (8.05,-2.1) -- (G2.east);
  % 第三级：根
  \node[hw compute, minimum width=0.75cm] (rt) at (4.6,-3.1) {\&};
  \draw[hw bus] (G1.south) -- (2.3,-3.1) -- (rt.west);
  \draw[hw bus] (G2.south) -- (6.9,-3.1) -- (rt.east);
  \draw[hw bus] (rt.south) -- (4.6,-3.8);
  \node[hw label, anchor=west] at (4.8,-3.8) {1 位结果};
  % 深度标尺
  \draw[{Stealth[length=1.8mm]}-{Stealth[length=1.8mm]}, BookRed, line width=0.6pt] (9.5,0.2) -- (9.5,-3.1);
  \node[hw label, anchor=west] at (9.65,-1.45) {3 级};
\end{pdfdiagram}
\diagramnote{同样 8 位，两种形状的延迟量级一眼可分。上：串行进位链，进位必须逐级穿过 8 个全加器，延迟随位宽 $n$ 线性增长，是数据通路上最贵的一级。下：归约树，每级把输入减半，8 位收成 1 位只要 3 级，延迟随 $\log_2 n$ 增长。比较与归约便宜、算术加法贵，差别就在这两种形状上。}

量级的意义在于预算。一条 32 位加法链是几十级门的延迟，一个 32 位相等比较只要六级左右；写 \code{cnt == 32'd50000} 比写 \code{cnt >= 32'd50000} 便宜得多。“电源、时钟与信号完整性”一章算时序预算时会再回到这张表；现在先养成习惯：每写一个运算符，报得出它是什么电路、什么量级。

把预算练成直觉，最快的办法是给熟悉的模块估一遍价。以本书末章的教学计算机项目教学计算机的 8 位数据通路为例：

\begin{longtable}{L{0.3\linewidth}L{0.26\linewidth}L{0.28\linewidth}}
\toprule
数据通路上的操作 & 对应硬件 & 量级估计 \\
\midrule
\code{A + B}（加法指令） & 8 位串行进位链 & 约 8 级进位延迟 \\\tablerowrule
\code{A == 8'h00}（零标志） & 8 个 XNOR 加 3 级与树 & 约 4 级门延迟 \\\tablerowrule
\code{sel ? alu\_out : mem\_out} & 8 个 2 选 1 mux & 1 级门延迟 \\\tablerowrule
操作码译码选控制信号 & 4 位译码加 mux 阵列 & 2 到 3 级门延迟 \\
\bottomrule
\end{longtable}

对照“电源、时钟与信号完整性”一章的时钟周期核算：这一条链路上最贵的是加法器的进位链，所以它往往决定最小时钟周期；译码和 mux 加起来不过几级门，通常不是瓶颈。能在写代码之前说出“这条路径上最贵的是哪个运算符”，就已经具备了一半时序直觉；剩下的一半在触发器和时钟偏斜里，第三、四节接着讲。

\topic{case 语句与优先结构}\code{case} 和 if-else 链都可能综合成多路选择结构，但形状不同，表达的意图也不同。\code{case} 表达“按选择信号的值挑一个”：综合器对选择信号译码，译码结果控制一组并行的选择路径，各分支地位平等，数据到输出的路径深度相同。if-else 链表达“先看第一个条件，不成立再看下一个”：综合成一串级联的 2 选 1 mux，越靠前的条件控制越靠近输出的 mux——路径最短；最后一个条件的数据要穿过整条链。四个分支，两种写法：

\begin{lstlisting}[language=Verilog]
// 写法一：并行译码，四路地位平等
always_comb begin
    case (sel)
        2'b00:   y = d0;
        2'b01:   y = d1;
        2'b10:   y = d2;
        default: y = d3;
    endcase
end

// 写法二：优先级链，req0 最优先
always_comb begin
    if      (req0) y = d0;
    else if (req1) y = d1;
    else if (req2) y = d2;
    else           y = d3;
end
\end{lstlisting}

硬件差别对照如下：

\begin{longtable}{L{0.2\linewidth}L{0.32\linewidth}L{0.32\linewidth}}
\toprule
 & 并行译码（case） & 优先级链（if-else） \\
\midrule
选择结构 & 译码器加并行 mux & 级联 2 选 1 mux \\\tablerowrule
各分支路径 & 等深 & 逐级加深，链尾最深 \\\tablerowrule
分支间关系 & 互斥，由译码保证 & 靠前者屏蔽靠后者 \\\tablerowrule
天然适用 & 按编码选路：操作码、地址段 & 按优先权裁决：中断、报警 \\
\bottomrule
\end{longtable}

两种综合结果画成结构图，“等深”与“逐级加深”的差别一眼可分：

\begin{pdfdiagram}[x=0.8cm]
  % ===== 左：并行译码 =====
  \node[hw label] at (0.9,3.4) {并行译码：四路等深};
  \node[hw compute, minimum width=1.5cm, minimum height=2.6cm] (mux4) at (0.9,0) {4 选 1\\mux};
  \node[hw label, anchor=east] at (-1.9,0.98) {d0};
  \node[hw label, anchor=east] at (-1.9,0.33) {d1};
  \node[hw label, anchor=east] at (-1.9,-0.33) {d2};
  \node[hw label, anchor=east] at (-1.9,-0.98) {d3};
  \draw[hw bus] (-1.75,0.98) -- ($(mux4.west)+(0,0.98)$);
  \draw[hw bus] (-1.75,0.33) -- ($(mux4.west)+(0,0.33)$);
  \draw[hw bus] (-1.75,-0.33) -- ($(mux4.west)+(0,-0.33)$);
  \draw[hw bus] (-1.75,-0.98) -- ($(mux4.west)+(0,-0.98)$);
  \node[hw control, minimum width=1.5cm, minimum height=0.9cm] (dec) at (0.9,2.5) {译码器};
  \node[hw label, anchor=east] at (-1.9,2.5) {sel[1:0]};
  \draw[hw return, dashed] (-1.75,2.5) -- (dec.west);
  \draw[hw return, dashed] (dec.south) -- (mux4.north);
  \node[hw label, anchor=west] at (1.1,1.75) {4 条互斥选择线};
  \draw[hw bus] (mux4.east) -- (3.0,0);
  \node[hw label, anchor=west] at (3.1,0) {y};
  \node[hw label, anchor=north] at (0.9,-1.7) {每条数据路径同为 1 级};
  % ===== 右：优先级链 =====
  \node[hw label] at (7.3,2.2) {优先级链：req0 最优先};
  \node[hw control, minimum width=1.1cm, minimum height=0.95cm] (m2) at (5.6,0) {2 选 1};
  \node[hw control, minimum width=1.1cm, minimum height=0.95cm] (m1) at (7.3,0) {2 选 1};
  \node[hw control, minimum width=1.1cm, minimum height=0.95cm] (m0) at (9.0,0) {2 选 1};
  \node[hw label, above] at (4.65,0.02) {d3};
  \draw[hw bus] (4.3,0) -- (m2.west);
  \draw[hw bus] (m2.east) -- (m1.west);
  \draw[hw bus] (m1.east) -- (m0.west);
  \draw[hw bus] (m0.east) -- (10.3,0);
  \node[hw label, anchor=west] at (10.4,0) {y};
  \node[hw label] at (5.6,1.6) {d2};
  \draw[hw bus] (5.6,1.35) -- (m2.north);
  \node[hw label] at (7.3,1.6) {d1};
  \draw[hw bus] (7.3,1.35) -- (m1.north);
  \node[hw label] at (9.0,1.6) {d0};
  \draw[hw bus] (9.0,1.35) -- (m0.north);
  \node[hw label] at (5.6,-1.15) {req2};
  \draw[hw return, dashed] (5.6,-0.9) -- (m2.south);
  \node[hw label] at (7.3,-1.15) {req1};
  \draw[hw return, dashed] (7.3,-0.9) -- (m1.south);
  \node[hw label] at (9.0,-1.15) {req0};
  \draw[hw return, dashed] (9.0,-0.9) -- (m0.south);
  \node[hw label, anchor=north] at (7.3,-1.7) {d0 穿 1 级，d3 穿 3 级——队尾最慢};
\end{pdfdiagram}
\diagramnote{左：\code{case} 综合成译码器加并行 mux，sel 译出四条互斥选择线，d0 到 d3 到达输出的路径同为一级，地位平等。右：if-else 链综合成级联 2 选 1，req0 控制的那级紧贴输出，d0 只穿一级，链尾的 d3 要穿三级——优先级不是注释，是物理上的路径长度。}

“数字硬件基础”部分的三级报警就是优先级链的天然例子：\code{fire} 一来，\code{fault} 和 \code{info} 的输出被逻辑本身关断——用 if-else 链写，\code{fire} 放第一级，硬件自动获得“最高优先级路径最短”的性质，与“数字硬件基础”部分“越重要的信号越快”的结论吻合。反过来，译码场景写 if-else 链则是自找麻烦：处理器核心相关章节指令译码按 \code{opcode} 选控制信号，十六种编码地位平等，用 \code{case} 既平行又清楚。

SystemVerilog 还提供两个意图声明。\code{unique case} 声明各分支互斥且覆盖完备，综合器可放心按并行译码优化，若实际输入打破声明（分支重叠或遗漏），仿真会报告；\code{priority case} 声明按书写顺序带优先级、允许未列出的组合不覆盖，综合器按链式结构处理。意图声明不是装饰，是让工具替你检查约定的写法；但它只影响检查与优化方向，不改变你写下的逻辑关系本身。

优先级链有一个直接的量化后果：链越长，队尾越慢。四级 if-else 的最后一个分支，数据要穿过三级串联的 2 选 1 mux；而把三级报警写成 if-else 链，\code{fire} 的路径是零级 mux——直接就是输出：

\begin{lstlisting}[language=Verilog]
// 三级报警的 if-else 版：硬件与“数字硬件基础”部分逐门对应
always_comb begin
    if (fire) begin
        fire_lamp  = 1'b1;
        fault_lamp = 1'b0;
        info_lamp  = 1'b0;
    end else if (fault) begin
        fire_lamp  = 1'b0;
        fault_lamp = 1'b1;
        info_lamp  = 1'b0;
    end else if (info) begin
        fire_lamp  = 1'b0;
        fault_lamp = 1'b0;
        info_lamp  = 1'b1;
    end else begin
        fire_lamp  = 1'b0;
        fault_lamp = 1'b0;
        info_lamp  = 1'b0;
    end
end
\end{lstlisting}

综合器把它展开后，\code{fault\_lamp} 只在“无 \code{fire} 且有 \code{fault}”时为 1——\code{fire} 的反相出现在 \code{fault} 的使能项里，与“数字硬件基础”部分“高优先级以反相形式出现在低优先级乘积项中”的结构逐字对应。分支里的常量 \code{1'b0} 不是“关灯动作”，是把对应输出接到固定电平；三个输出在每条路径上都有确定值，这段代码不会推断锁存器，下一专题的默认值纪律它天然满足。

顺带提一下 \code{casez}：它允许分支项里写 \code{?} 作为“这位不比较”的通配，适合译码时只关心部分位的场合——例如“高三位是 \code{101} 就选中，低位任意”。通配位在硬件上是译码矩阵里被省去的比较，写得好能省门；但多个带通配的分支可能同时匹配，\code{casez} 按书写顺序取第一个匹配项，这就把优先级悄悄引了回来，读 \code{casez} 必须像读 if-else 链一样留意次序。入门阶段建议先用朴素的 \code{case} 写全编码，确认需要通配再换 \code{casez}。

\topic{默认赋值与锁存器推断}机制：组合块的输出必须在所有输入组合下都有确定驱动。如果某条执行路径走到底也没给某个输出赋值，综合器面对的规格就成了“这种情况下输出保持原值”——保持需要存储，于是推断出一个电平敏感的锁存器。这是“状态、时钟与时序逻辑”一章讲过的那种靠反馈维持状态的元件，只是这一次不是你画的，是工具替你补的，补完往往只在综合日志里留一行警告。

\begin{lstlisting}[language=Verilog]
// 危险写法：sel 为 2'b11 时 y 没有新值，推断出锁存器
always_comb begin
    case (sel)
        2'b00: y = d0;
        2'b01: y = d1;
        2'b10: y = d2;
    endcase
end
\end{lstlisting}

“没写就不变”在程序里天经地义，在组合块里是危险信号：一段电平敏感的存储悄悄混进了你以为纯组合的路径。锁存器在使能期间对输入直通，毛刺跟着穿过去；它把时序分析从“周期边界”变成“电平窗口”，工具检查困难；测试设备对锁存器的可控性也远差于触发器。“状态、时钟与时序逻辑”一章花了一整节把锁存器升级成边沿触发器，工程 RTL 对意外锁存器的态度就是那一次升级的延续：除非故意，绝不引入。

消除方法有两种，工程上首选第一种。方法一：块首先给所有输出写默认值，再用分支覆盖需要的项——默认值保证任何路径都有确定驱动，锁存器无从产生。方法二：把分支补全，给 \code{case} 加 \code{default}、给 \code{if} 补 \code{else} 覆盖默认路径。两种写法如下：

\begin{lstlisting}[language=Verilog]
// 方法一：默认值先行，分支只写例外（推荐）
always_comb begin
    y = d3;                 // 默认值
    case (sel)
        2'b00: y = d0;
        2'b01: y = d1;
        2'b10: y = d2;
        default: y = d3;
    endcase
end

// 方法二：分支补全，每条路都有明确驱动
always_comb begin
    if (sel == 2'b00)      y = d0;
    else if (sel == 2'b01) y = d1;
    else if (sel == 2'b10) y = d2;
    else                   y = d3;
end
\end{lstlisting}

方法一在长代码里优势明显：输出有十个时，默认值集中写在块首，一眼可查“每个输出的归宿是什么”；方法二把默认赋值散在各分支末尾，输出一多就容易漏一个——而漏一个就是一只锁存器。多输出的组合块几乎总是选方法一。

把坏写法与修复写法分别综合成硬件摆在一起，差别就直观了：前者的输出路径上多了一只你从未画过的存储元件。

\begin{pdfdiagram}
  % 上：缺省分支，综合器补出锁存器
  \node[hw label] at (1.4,5.05) {缺省分支：组合网络之后被补上一只锁存器};
  \node[hw compute, minimum width=2.6cm, minimum height=1.5cm] (cbad) at (0.2,3.5) {组合网络\\（case 译码）};
  \node[hw memory, minimum width=2.2cm, minimum height=1.1cm] (lat) at (3.6,3.5) {电平敏感\\锁存器};
  \node[hw label, anchor=east] at (-2.7,3.95) {d0};
  \node[hw label, anchor=east] at (-2.7,3.5) {d1};
  \node[hw label, anchor=east] at (-2.7,3.05) {d2};
  \node[hw label, anchor=east] at (-2.7,4.7) {sel};
  \draw[hw bus] (-2.55,3.95) -- (-1.1,3.95);
  \draw[hw bus] (-2.55,3.5) -- (-1.1,3.5);
  \draw[hw bus] (-2.55,3.05) -- (-1.1,3.05);
  \draw[hw return, dashed] (-2.55,4.7) -- (0.2,4.7) -- (0.2,4.28);
  \draw[hw bus] (cbad.east) -- (lat.west);
  \draw[hw bus] (lat.east) -- (5.5,3.5);
  \node[hw label, anchor=west] at (5.6,3.5) {y};
  \draw[hw return, dashed] (cbad.south) -- (0.2,2.3) -- (3.6,2.3) -- (3.6,2.92);
  \node[hw label, anchor=west] at (0.3,2.02) {使能：sel=11 时关门，y 保持旧值};
  % 下：完整赋值，纯组合
  \node[hw label] at (1.4,1.66) {默认值先行或分支补全：纯组合网络};
  \node[hw compute, minimum width=2.6cm, minimum height=1.5cm] (cok) at (0.2,0.2) {组合网络\\（覆盖全部分支）};
  \node[hw label, anchor=east] at (-2.7,0.65) {d0};
  \node[hw label, anchor=east] at (-2.7,0.2) {d1};
  \node[hw label, anchor=east] at (-2.7,-0.25) {d2};
  \node[hw label, anchor=east] at (-2.7,1.4) {sel};
  \draw[hw bus] (-2.55,0.65) -- (-1.1,0.65);
  \draw[hw bus] (-2.55,0.2) -- (-1.1,0.2);
  \draw[hw bus] (-2.55,-0.25) -- (-1.1,-0.25);
  \draw[hw return, dashed] (-2.55,1.4) -- (0.2,1.4) -- (0.2,0.98);
  \draw[hw bus] (cok.east) -- (5.5,0.2);
  \node[hw label, anchor=west] at (5.6,0.2) {y};
\end{pdfdiagram}
\diagramnote{同一个 \code{always\_comb} 的两种综合结果。上：\code{sel}=11 时没有任何分支给 \code{y} 新值，综合器补出一只电平敏感锁存器——使能关门期间 \code{y} 由存储维持，透明窗口与毛刺穿透随之混入。下：默认值先行（或分支补全）后，所有输入组合都有确定驱动，输出直接来自组合网络，一个存储元件也没有。}

多输出的例子更能说明问题。下面这个块有三个输出，\code{valid} 分支里漏写了 \code{err}：

\begin{lstlisting}[language=Verilog]
// err 在 valid=1 的路径上没有新值：一只锁存器
always_comb begin
    if (valid) begin
        data_out = data_in;
        ready    = 1'b1;
    end else begin
        data_out = 8'h00;
        ready    = 1'b0;
        err      = 1'b0;
    end
end

// 修好：三个输出的默认值先行，分支只写例外
always_comb begin
    data_out = 8'h00;
    ready    = 1'b0;
    err      = 1'b0;
    if (valid) begin
        data_out = data_in;
        ready    = 1'b1;
    end
end
\end{lstlisting}

坏版本里 \code{err} 在 \code{valid}=1 期间保持旧值——一只锁存器只挂在三个输出之一上，其余两个仍是纯组合。这种“半组合半存储”的混合体最容易蒙混过关：模块在多数测试里表现正常，只有特定输入序列才让 \code{err} 显出记忆。好版本里三个默认值排成一列，每个输出的归宿一目了然，分支只描述“什么情况下偏离默认”。

写完任何 \code{always\_comb}，按这张清单过一遍：

\begin{longtable}{L{0.34\linewidth}L{0.48\linewidth}}
\toprule
检查点 & 预期结果 \\
\midrule
每个输出在块首有默认值，或每条路径都被覆盖 & 任一输入组合下输出都有确定驱动 \\\tablerowrule
\code{case} 有 \code{default} 或分支枚举完备 & 任何输入组合都有分支接住 \\\tablerowrule
综合日志无 inferred latch 字样 & 工具与你对硬件的想象一致 \\\tablerowrule
块内读到的信号都会触发重算 & \code{always\_comb} 敏感表自动覆盖 \\
\bottomrule
\end{longtable}

\begin{warningbox}
常见误区：“综合器帮我补了锁存器，功能仿真也对，就这样用。”仿真对是因为激励恰好没踩中电平窗口里的毛刺。意外锁存器的故障模式是温度、电压、批次相关的间歇性错误，最难排查的一类。综合日志里出现 inferred latch 字样，应当作错误处理，回去补默认值。
\end{warningbox}

\topic{把组合块写成可读的硬件文字}可读性的标准只有一条：每一行都能立刻读回具体硬件。三条习惯服务于这个标准。第一，输入默认、逐项覆盖：块首给默认值，下面的 \code{case} 只写例外项，读者先看“通常是什么”，再看“什么时候不是”。第二，避免嵌套三元：\code{sel ? (s2 ? a : b) : (s3 ? c : d)} 是一团看不清形状的 mux 树，同样的硬件用 \code{case} 写出来，结构一目了然。第三，宽信号先切片命名：把 \code{data[7:4]} 起名为 \code{hi\_nibble} 再使用，读者看到的是有意义的线名，而不是大量缺乏语义的位号。

按这三条习惯，把算术与数据通路相关章节的 ALU 写成 \code{always\_comb} 版本。功能与算术与数据通路相关章节一致：加、减、按位与、按位或、按位异或，外加零标志：

\begin{lstlisting}[language=Verilog]
module alu8 (
    input  wire [7:0] a,
    input  wire [7:0] b,
    input  wire [2:0] op,
    output reg  [7:0] result,
    output wire       zero
);
    always_comb begin
        case (op)
            3'b000:  result = a + b;    // ADD：8 位加法器
            3'b001:  result = a - b;    // SUB：加补器
            3'b010:  result = a & b;    // AND：8 个与门
            3'b011:  result = a | b;    // OR ：8 个或门
            3'b100:  result = a ^ b;    // XOR：8 个异或门
            default: result = 8'h00;    // 默认路径：避免推断锁存器
        endcase
    end

    assign zero = (result == 8'h00);    // 8 个 XNOR 加与树
endmodule
\end{lstlisting}

逐行读回硬件。五种运算各是一群候选电路，全部并排存在、同时在工作：\code{a + b} 一条 8 位进位链，\code{a - b} 沿用算术与数据通路相关章节的加补方案，三个按位运算各 8 个门。\code{case} 对 \code{op} 译码，控制一组 8 比特宽的 5 选 1 mux，从五群候选里挑出 \code{result}——这正是算术与数据通路相关章节“候选结果并排算好再选择”的结构，只是从门级图纸换成了文字。\code{default} 保证完备，这个块不可能推断出锁存器。\code{result} 声明为 \code{reg}，只因为驱动它的是 always 块；它身上没有一位触发器，纯粹是组合网络的输出端。\code{zero} 单独用 \code{assign} 写，因为“一行说得清”，分工习惯在此兑现。

最后说一句意图与实现的关系。算术与数据通路相关章节讲过 ADD 与 SUB 应复用同一条加法链；这里写成 \code{a + b} 和 \code{a - b} 两行，综合器通常能识别并共享加法器，但不保证。需要严格复用时，就得像算术与数据通路相关章节那样显式写出 \code{b\_to\_adder} 和 \code{carry\_in}，把复用变成文字里看得见的结构。这是 RTL 的普遍处境：写法表达意图，综合器决定实现；意图越接近结构，实现越接近预期。“时序逻辑的 Verilog 写法”一节进入时序逻辑，触发器登场之后，这条原则只会更重要。

把“严格复用”写出来看看。下面的版本只用一条加法链，ADD 与 SUB 共享，与算术与数据通路相关章节 74HC283、74HC86、74HC157 三片芯片的实验台同构：

\begin{lstlisting}[language=Verilog]
module alu8_shared (
    input  wire [7:0] a,
    input  wire [7:0] b,
    input  wire [2:0] op,
    output reg  [7:0] result,
    output wire       zero,
    output wire       carry
);
    reg [7:0] b_to_adder;   // B 通路上的可控反相
    reg       carry_in;     // 最低位进位
    wire [8:0] sum9;        // 9 位：和加进位，各得其所

    always_comb begin
        // 默认值先行：加法，B 原样通过，进位 0
        b_to_adder = b;
        carry_in   = 1'b0;
        if (op == 3'b001) begin     // SUB：B 取反加一
            b_to_adder = ~b;
            carry_in   = 1'b1;
        end
    end

    assign sum9 = {1'b0, a} + {1'b0, b_to_adder} + {8'b0, carry_in};

    always_comb begin
        case (op)
            3'b000, 3'b001: result = sum9[7:0];  // ADD/SUB 共用一条链
            3'b010:         result = a & b;
            3'b011:         result = a | b;
            3'b100:         result = a ^ b;
            default:        result = 8'h00;
        endcase
    end

    assign carry = sum9[8];
    assign zero  = (result == 8'h00);
endmodule
\end{lstlisting}

逐行读回硬件。\code{b\_to\_adder} 前面的块是 8 个可控反相器——\code{op} 控制每比特过原值还是反值，正是算术与数据通路相关章节 B 输入那排 XOR 门；\code{carry\_in} 由同一个 \code{op} 派生，对应 \code{sub} 接最低位 \code{Cin}。\code{sum9} 左边主动写成 9 位并用拼接扩位，位宽纪律在此兑现：进位不再被静默丢弃，而是进入有名信号 \code{carry}。\code{case} 把 \code{sum9[7:0]} 的两个分支合并书写，文字层面就声明了共享——综合器不必再猜测。三个块各有分工：B 通路准备、加法链、结果选择，对应算术与数据通路相关章节实验台上的三片芯片。\code{b\_to\_adder} 与 \code{carry\_in} 声明为 \code{reg} 并由 always 块驱动，但它们没有保持任何旧值的路径——默认值先行保证了两端都是纯组合。

这个版本多写了十几行，换来的是面积的确信：一条加法链而不是两条。选择写哪个版本，取决于你在乎什么——在乎写法简短，用前一个；在乎实现确定，用这一个。两种写法都没有对错，区别只在“把多少决定权留给综合器”。本章反复说的话在此又响一遍：Verilog 不是程序，是用文字画电路——你画到哪一层，工具就接手到哪一层。

\section{三、时序逻辑的 Verilog 写法}

前两节写的都是组合逻辑：给定输入，输出立即确定，电路本身没有记忆。“状态、时钟与时序逻辑”一章已经说明，记忆来自触发器——一组 D 触发器共用一个时钟，每个上升沿把各自 D 端的值抄进 Q 端，沿与沿之间 Q 保持不变。时序逻辑的 Verilog 写法因此只描述一件事：存在哪些触发器，每个时钟沿它们的 D 端接的是什么。

\topic{always\_ff 的标准模板}

先讲机制。D 触发器有三个端口：时钟、D、Q。它只在时钟上升沿动作一次，把此刻 D 的值抄进 Q；D 在沿与沿之间怎么变化，Q 一概不理。Verilog 用敏感表声明“什么事件会让这块硬件动作”，于是时钟是敏感表里唯一的内容。一个寄存器的最小模板如下：

\begin{lstlisting}
module dff (
    input  logic clk,
    input  logic d,
    output logic q
);
    // 一只 D 触发器：上升沿把 d 采进 q
    always_ff @(posedge clk) begin
        q <= d;
    end
endmodule
\end{lstlisting}

这三行读回硬件没有歧义：\code{posedge clk} 是触发器的时钟端，\code{q} 是 Q 端，\code{d} 通过非阻塞赋值接到 D 端。用本章开头的三问套一遍：这一行产生了什么硬件——一只触发器；这个信号是组合路径还是触发器——\code{q} 是触发器输出，\code{d} 是组合路径的终点；综合器会推断什么多余的东西——不会，模板里每个符号都有明确的硬件归宿。三问都答得出，模板才算读懂。

一个 \code{always\_ff} 块可以给多个信号赋值，得到的是一排共用同一时钟的触发器：

\begin{lstlisting}
// 8 位寄存器：八只触发器共用一个时钟
logic [7:0] reg8;

always_ff @(posedge clk) begin
    reg8 <= data_in;
end
\end{lstlisting}

“状态、时钟与时序逻辑”一章的 8 位寄存器就是这个模样的八份并排。位宽写在声明里，赋值语句一行不变——文字的简法是 Verilog 的长处，但读的时候要记得展开：这一行是八只触发器，不是一只。

为什么敏感表不写 \code{d}？因为 D 端变化不会让触发器动作，它只是改变“下一次沿到来时将被采样的值”。若把敏感表写成 \code{@(posedge clk or d)}，等于声明“\code{d} 的任何变化也是一次提交事件”——真实世界里不存在这种触发器，综合器要么报错，要么推断出一个并不想要的结构。这里藏着一条读 RTL 的基本功：敏感表不是“块内读到了哪些信号”的清单，而是“哪些边沿构成提交边界”的声明。块内用到却不在敏感表里的信号，正是穿过组合网络送到 D 端的数据输入；区分时钟与数据，先看敏感表，再看赋值方向。

\code{always\_ff} 是 SystemVerilog 的关键字，比 Verilog-2001 的 \code{always @(posedge clk)} 多一层自我声明：它明确告诉综合器“我意图描述触发器”。若块内写法推不出触发器——例如同一个信号在多个块里被赋值——综合器会给出警告，而不是默默猜测。读旧代码遇到 \code{always @(posedge clk)}，硬件含义完全相同，只是少了这层声明带来的检查。

边界只有一条，但必须记住：一个信号只能在一个 \code{always\_ff} 块里被赋值。一只触发器的 D 端只有一个驱动源；同一个信号出现在两个块的 \code{<=} 左端，等于把两个输出短接到同一根导线上，综合器会报多驱动错误。这也是“Verilog 不是程序”的一条硬理由——程序里两个函数可以先后写同一个变量，硬件里两个驱动源不能同时推一根导线。

\topic{非阻塞赋值的语义：先读旧值，同时提交}

机制仍然是“状态、时钟与时序逻辑”一章那句话：时钟沿到来时，所有触发器在同一瞬间翻转，采样的都是沿之前的旧值。非阻塞赋值 \code{<=} 就是这句话的文字形式。一个 \code{always\_ff} 块内所有 \code{<=} 的右端，都用沿到来之前的值计算；全部右端算完之后，左端才同时更新。读 \code{<=} 的正确姿势不是“执行这一句”，而是“给这只触发器的 D 端接线”。

用交换两个寄存器来验证这套语义：

\begin{lstlisting}
always_ff @(posedge clk) begin
    a <= b;   // A 的 D 端接 B 的 Q
    b <= a;   // B 的 D 端接 A 的 Q
end
\end{lstlisting}

两个右端读的都是旧值，于是一个沿之后 \code{a} 与 \code{b} 互换。硬件上这是两只触发器交叉对接：A 的 D 接 B 的 Q，B 的 D 接 A 的 Q。两行代码谁先谁后不影响结果，因为接线没有先后——这与“状态、时钟与时序逻辑”一章“时钟沿同时翻转”的图景严格对应，非阻塞赋值就是把“同时”写进语言的方式。逐拍看更清楚：

\begin{longtable}{L{0.18\linewidth}L{0.16\linewidth}L{0.16\linewidth}L{0.34\linewidth}}
\toprule
时刻 & \code{a} & \code{b} & 说明 \\
\midrule
沿之前 & 1 & 0 & 两个右端各自采样旧值 \\\tablerowrule
沿之后 & 0 & 1 & 两端同时提交，完成互换 \\\tablerowrule
再一个沿后 & 1 & 0 & 再次互换，回到原样 \\
\bottomrule
\end{longtable}

同一场合换成阻塞赋值 \code{=}，语义立刻改变。\code{=} 先更新左端，后面的语句读到的是新值：

\begin{lstlisting}
always_ff @(posedge clk) begin
    a = b;    // a 立刻变成 b 的值
    b = a;    // 读到的是新 a，即 b 的旧值
end
\end{lstlisting}

结果 \code{a} 与 \code{b} 都变成 \code{b} 的旧值，交换失败；硬件上相当于两只触发器的 D 端都接到了 B 的 Q，A 的旧值被丢弃。更麻烦的是，\code{=} 让结果依赖书写顺序——两行一调换，推断出的接线就跟着变。仿真行为随书写顺序漂移，正是把 Verilog 当程序读要付的第一笔学费。

把两种写法画到同一条时间轴上（沿前值均为 a=1、b=0），“同时提交”与“立即生效”就成了两种波形形状：

\begin{waveform}[x=0.9cm,y=0.85cm]
  % 参考虚线：三个时钟沿
  \draw[BookRule, line width=0.45pt, dashed] (2,7.1) -- (2,1.6);
  \draw[BookRule, line width=0.45pt, dashed] (4,7.1) -- (4,1.6);
  \draw[BookRule, line width=0.45pt, dashed] (6,7.1) -- (6,1.6);
  % clk：沿在 x=2,4,6
  \draw[BookAccent, line width=0.9pt] (0.5,6.3) -- (2,6.3) -- (2,6.9) -- (3,6.9) -- (3,6.3) -- (4,6.3) -- (4,6.9) -- (5,6.9) -- (5,6.3) -- (6,6.3) -- (6,6.9) -- (7,6.9) -- (7,6.3) -- (8.6,6.3);
  \node[hw label, anchor=east] at (0.3,6.6) {clk};
  \node[hw label] at (3.0,6.05) {沿 1：两端同时提交，完成互换};
  \node[hw label] at (7.0,6.05) {沿 2：再次互换};
  % 非阻塞：a、b 每拍互换
  \draw[BookAccent, line width=0.9pt] (0.5,5.75) -- (2,5.75) -- (2,5.15) -- (4,5.15) -- (4,5.75) -- (6,5.75) -- (6,5.15) -- (8.6,5.15);
  \node[hw label, anchor=east] at (0.3,5.45) {a（\code{<=}）};
  \draw[BookAccent, line width=0.9pt] (0.5,4.15) -- (2,4.15) -- (2,4.75) -- (4,4.75) -- (4,4.15) -- (6,4.15) -- (6,4.75) -- (8.6,4.75);
  \node[hw label, anchor=east] at (0.3,4.45) {b（\code{<=}）};
  % 阻塞：a=b 先写入，b=a 读到新 a，两值皆 0
  \draw[BookAccent, line width=0.9pt] (0.5,3.45) -- (2,3.45) -- (2,2.85) -- (8.6,2.85);
  \node[hw label, anchor=east] at (0.3,3.15) {a（\code{=}）};
  \draw[BookAccent, line width=0.9pt] (0.5,1.95) -- (8.6,1.95);
  \node[hw label, anchor=east] at (0.3,1.95) {b（\code{=}）};
  \node[hw label, anchor=north] at (4.55,1.5) {\code{a=b} 先写入，\code{b=a} 读到的是新 a——两拍后仍是 0、0，交换失败};
\end{waveform}
\diagramnote{沿前值 a=1、b=0。上两行：\code{<=} 的右端都读沿前旧值，沿 1 之后 a、b 互换，沿 2 再换回——硬件是两只触发器 D 端交叉互连。下两行：\code{=} 立即生效，\code{a=b} 先把 a 写成 0，\code{b=a} 读到的是刚写入的新 a，b 停在 0——等效于两只触发器的 D 端同接 b 的 Q，交换被同化。}

“同时提交”还有一个对全书都重要的推论：\code{<=} 的左端在本沿之后、下一沿之前保持不变，于是组合逻辑在沿与沿之间读到的是稳定电平。“状态、时钟与时序逻辑”一章分析建立时间、保持时间和时钟偏斜时，预设了“采样沿前后各有一段安静期”；在非阻塞赋值的世界里，这个预设自动成立。这就是同步设计把 \code{<=} 定为标准的根本原因，也是本节要记的规则：时序块里只写 \code{<=}。“从 Verilog 读回硬件：综合边界”一节还会用移位寄存器的正反两例，把这条规则在综合层面再算一遍账。

\topic{复位的两种写法}

“状态、时钟与时序逻辑”一章讲复位时分过工：拉入要异步，保证任何时刻都能进入安全状态；释放要同步，让所有触发器在同一拍恢复。Verilog 的两种复位写法正好对应这两个选择，差别只在敏感表。

异步复位把复位信号本身当作一个提交事件，所以它进敏感表：

\begin{lstlisting}
// 异步复位：rst_n 的下降沿不等时钟，立即清零
always_ff @(posedge clk or negedge rst_n) begin
    if (!rst_n)
        q <= '0;   // 复位沿到来，Q 直接归 0
    else
        q <= d;    // 时钟沿到来，正常采样
end
\end{lstlisting}

读回硬件：这是一只带异步清零端的触发器，“状态、时钟与时序逻辑”一章的 74HC74 就有这样一只 \code{/CLR} 引脚。\code{rst\_n} 的下降沿不经过时钟，直接把 Q 拉到 0；\code{rst\_n} 为高时，每个时钟沿正常采样 D。敏感表里出现两个边沿，是“异步复位”在文字上的签名。名字里的 \code{\_n} 后缀与低电平有效的 \code{negedge} 也是一对约定：看到 \code{rst\_n}，就知道复位是拉低生效。

同步复位不进敏感表，它只是 D 端网络的一个输入：

\begin{lstlisting}
// 同步复位：是否清零由时钟沿裁决
always_ff @(posedge clk) begin
    if (!rst_n)
        q <= '0;   // 沿到来且复位有效，采 0
    else
        q <= d;    // 沿到来且复位无效，采 d
end
\end{lstlisting}

代码只差敏感表里的几个字符，硬件却完全不同：\code{rst\_n} 不再接触发器的清零端，而是变成 D 端前面一只二选一 mux 的选择端——复位有效时 mux 选 0，否则选 \code{d}。复位何时生效完全由时钟沿决定，与时钟域外的世界没有直接通道。

两种写法读回的硬件并排画出，差别只在复位信号接到哪里。

\begin{pdfdiagram}[x=0.9cm]
  % 左：异步复位，rst_n 直连 /CLR
  \node[hw label] at (2.0,1.0) {异步复位（带 /CLR 端）};
  \node[hw state, minimum width=2.0cm] (ffa) at (2.0,0) {D 触发器};
  \node[hw label, anchor=east] at (0.0,0.25) {D};
  \draw[hw bus] (0.1,0.25) -- ($(ffa.west)+(0,0.25)$);
  \node[hw label, anchor=east] at (0.0,-0.25) {clk};
  \draw[hw return] (0.1,-0.25) -- ($(ffa.west)+(0,-0.25)$);
  \draw[hw bus] (ffa.east) -- (3.6,0);
  \node[hw label, anchor=west] at (3.7,0) {Q};
  \node[hw label] at (2.0,-1.75) {rst\_n};
  \draw[hw return] (2.0,-1.5) -- (ffa.south);
  \node[hw label, anchor=west] at (2.15,-0.95) {/CLR};
  \node[hw label] at (2.0,-2.3) {敏感表含 negedge rst\_n};
  % 右：同步复位，rst_n 只是 mux 的选择端
  \node[hw label] at (7.9,1.0) {同步复位（D 前一只 mux）};
  \node[hw control, minimum width=1.5cm] (mux) at (6.6,0) {2 选 1\\mux};
  \node[hw label, anchor=east] at (5.0,0.28) {D};
  \draw[hw bus] (5.1,0.28) -- ($(mux.west)+(0,0.28)$);
  \node[hw label, anchor=east] at (5.0,-0.28) {0};
  \draw[hw bus] (5.1,-0.28) -- ($(mux.west)+(0,-0.28)$);
  \node[hw label] at (6.6,-1.75) {rst\_n};
  \draw[hw return] (6.6,-1.5) -- (mux.south);
  \node[hw label, anchor=west] at (6.75,-0.95) {sel};
  \node[hw state, minimum width=1.7cm] (ffb) at (9.3,0) {D 触发器};
  \draw[hw bus] (mux.east) -- (ffb.west);
  \node[hw label, anchor=east] at (7.2,-1.3) {clk};
  \draw[hw return] (7.5,-1.0) -- (9.3,-1.0) -- (ffb.south);
  \draw[hw bus] (ffb.east) -- (10.8,0);
  \node[hw label, anchor=west] at (10.9,0) {Q};
  \node[hw label] at (7.9,-2.3) {敏感表只有 posedge clk};
\end{pdfdiagram}
\diagramnote{两种复位写法读回的硬件。左：异步复位把 rst\_n 接到触发器的异步清零端（74HC74 的 /CLR），下降沿不等时钟、立即清零，敏感表里因此出现 negedge rst\_n。右：同步复位的 rst\_n 只控制 D 端前面的二选一 mux，复位有效时 mux 选 0，是否清零由时钟沿裁决，敏感表里只有 posedge clk。}

什么时候用哪种？对照“状态、时钟与时序逻辑”一章“异步拉入、同步释放”的约定：来自芯片外部的复位按钮是异步事件，入口级用异步写法，保证任何时刻都能把状态机拉进安全状态；释放路径经过两级触发器同步之后，内部各模块看到的已是同步信号，其后的大量寄存器用同步写法即可。“状态、时钟与时序逻辑”一章那只复位同步小电路，写成 Verilog 是两行触发器：

\begin{lstlisting}
// 复位释放同步器：异步拉入、同步释放
always_ff @(posedge clk or negedge rst_n) begin
    if (!rst_n) begin
        rst_ff1 <= 1'b0;
        rst_ff2 <= 1'b0;
    end else begin
        rst_ff1 <= 1'b1;      // 释放时机由时钟裁决
        rst_ff2 <= rst_ff1;
    end
end
// 内部模块使用同步后的 rst_sync_n = rst_ff2
\end{lstlisting}

外部 \code{rst\_n} 拉低时，两级触发器立刻清零，复位“进得去”；外部释放后，\code{rst\_ff1}、\code{rst\_ff2} 在两个连续时钟沿逐级变高，内部所有寄存器看到的是同一个沿上翻来的释放信号，复位“出得来”而且出得整齐。这正是“状态、时钟与时序逻辑”一章“异步拉入、同步释放”的模板化形态。

同步写法还有两个实际好处：综合器可以把复位逻辑并进数据路径，省掉专用的全局复位网络；复位路径与普通数据路径一样参与建立时间检查，时序分析更简单。FPGA 和 ASIC 的标准单元库通常两种复位端的触发器都备有，选型由系统复位策略决定，不是语法偏好。下表把两种写法并排对照：

\begin{longtable}{L{0.13\linewidth}L{0.33\linewidth}L{0.22\linewidth}L{0.20\linewidth}}
\toprule
复位写法 & 读回的硬件 & 生效时刻 & 典型场合 \\
\midrule
异步复位 & 触发器的异步清零端（如 74HC74 的 \code{/CLR}） & \code{rst\_n} 边沿立即生效，不等时钟 & 上电、急停、外部复位入口 \\\tablerowrule
同步复位 & D 端前面的二选一 mux & 下一个时钟沿 & 内部已同步的复位域 \\
\bottomrule
\end{longtable}

最后一个细节：不是每个寄存器都要复位。“状态、时钟与时序逻辑”一章的原则是“控制状态机必须复位，数据通路视情况而定”。写进 Verilog 就是少写一个 \code{if}：不复位的寄存器，模板里干脆不出现 \code{rst\_n}，综合器就少布一根复位线，触发器也可能换成无复位端的更小单元。复位是硬件需求，不是代码洁癖。

\topic{使能端与计数器、移位寄存器模板}

“状态、时钟与时序逻辑”一章把计数器定义为“寄存器加‘加一组合逻辑’构成的反馈”。使能端是插在反馈路径上的一只 mux：使能有效时 D 接 \code{q+1}，无效时 D 接 \code{q} 自身。写成 Verilog 只有三行：

\begin{lstlisting}
// 标准计数器：en 有效才加一，否则保持
always_ff @(posedge clk) begin
    if (en)
        q <= q + 1'b1;
end
\end{lstlisting}

注意这里没有 \code{else} 分支。时序块里“没写”不等于“没变”，而是“D 端接回 Q 的旧值”——触发器天然会保持，保持不需要任何代码。这是时序块与组合块最关键的差别；“从 Verilog 读回硬件：综合边界”一节讲锁存器推断时，同一句“没写”在组合块里会惹出完全不同的麻烦。使能端也是同步设计里“减速”的正道：想让计数器每秒才加一，就给它一个每秒出现一次的使能脉冲，而不是去动时钟——“状态、时钟与时序逻辑”一章交通灯的 1Hz 分频使能就是这个思路；时钟本身应保持完整，“电源、时钟与信号完整性”一章还会从时钟树的角度重讲这条纪律。

工程里的计数器通常有三个控制：同步清零、装载、使能。三者的优先级写在 \code{if} 链里，越靠前优先级越高：

\begin{lstlisting}
// 带清零、装载、使能的计数器：优先级 clr > load > en
always_ff @(posedge clk) begin
    if (clr)
        q <= '0;              // 清零最优先
    else if (load)
        q <= load_val;        // 其次装载
    else if (en)
        q <= q + 1'b1;        // 再次计数
    // 三者都无效：保持
end
\end{lstlisting}

读回硬件：D 端是一条三级优先 mux 链，\code{clr} 控制的那一级离触发器最近、优先级最高。\code{if} 链的书写顺序就是 mux 的物理顺序——调整条件顺序，等于在电路里挪动 mux 的位置。“状态、时钟与时序逻辑”一章交通灯控制器里那个“装载 30/5/2 秒、减到 0 发 \code{timer\_done}”的定时器，把加一换成减一、再配一个比较器，就是这个模板的直接变形：

\begin{lstlisting}
// 交通灯定时器：装载秒数，1Hz 使能下递减，到 0 发脉冲
always_ff @(posedge clk) begin
    if (load)
        timer <= load_val;
    else if (tick_1hz && (timer != 0))
        timer <= timer - 1'b1;
end
assign timer_done = (timer == 0);
\end{lstlisting}

移位寄存器同样四行。以“状态、时钟与时序逻辑”一章反复出现的 74HC595 式串入并出为例：

\begin{lstlisting}
// 74HC595 式移位寄存器：en 有效时左移一格
always_ff @(posedge clk) begin
    if (en)
        q <= {q[6:0], sin};
end
\end{lstlisting}

拼接表达式读作连线：第 $i$ 只触发器的 D 接第 $i-1$ 只的 Q，第 0 只的 D 接 \code{sin}。八只触发器排成一列，每个沿所有 bit 同时挪一格——“同时”依然来自 \code{<=} 的旧值语义，第 $i$ 只采到的是第 $i-1$ 只的旧值，串行输入一个沿只能前进一级。\code{<=} 与 \code{=} 的差别在这条单句里还显不出来，多句级联的写法才有“一个沿里输入 bit 直达末尾”的坑，“从 Verilog 读回硬件：综合边界”一节开头用三级例子细算。595 芯片还有第二级存储寄存器——把移位结果一次性送到输出引脚的那一排——写成 Verilog 是另一个时钟下的整块复制：

\begin{lstlisting}
// 595 的完整模型：移位级 + 存储级，两个时钟各管一排
always_ff @(posedge srclk) begin
    if (sr_en)
        shift <= {shift[6:0], sin};
end

always_ff @(posedge rclk) begin
    q_out <= shift;   // 存储级：移位结果整体亮相
end
\end{lstlisting}

两个 \code{always\_ff} 块、两个时钟、两排触发器，与芯片手册里 SRCLK 和 RCLK 两只引脚一一对应。再读 595 的数据手册时看到这两只时钟，就能直接想起这两段代码。

\topic{两段式状态机模板}

“状态、时钟与时序逻辑”一章给状态机画的结构图是三块：状态寄存器、下一状态组合逻辑、输出组合逻辑。两段式模板把这个结构原样搬进 Verilog：一段 \code{always\_ff} 只做一件事——沿到来时 \code{state <= next\_state}；一段 \code{always\_comb} 根据当前状态和输入算出 \code{next\_state} 与全部输出。分工固定之后，读任何状态机都有了固定路线：\code{case} 语句对照转移表，输出赋值对照输出表，而 \code{always\_ff} 段永远只做一件事，没有值得多看的内容。模板的价值，就在于让无聊的段落真正无聊。

现在把“状态、时钟与时序逻辑”一章的交通灯 FSM 完整写成模板。需求复述一遍：东西绿 30 秒、东西黄 5 秒、全红 2 秒、南北绿 30 秒、南北黄 5 秒、全红 2 秒，循环往复；计时交给一个独立的带装载递减计数器，它减到 0 时给出单周期脉冲 \code{timer\_done}，这是状态机唯一的转移条件；灯色只由当前状态决定，是 Moore 型输出；复位入口选全红。先声明状态编码与灯色，并写状态寄存器段：

\begin{lstlisting}
typedef enum logic [2:0] {
    EW_GREEN, EW_YELLOW, ALL_RED_EW,
    NS_GREEN, NS_YELLOW, ALL_RED_NS
} state_t;

typedef enum logic [1:0] {RED, YELLOW, GREEN} light_t;

state_t state, next_state;
light_t ew_light, ns_light;

// 第一段：状态寄存器。沿到来时提交下一状态
always_ff @(posedge clk or negedge rst_n) begin
    if (!rst_n)
        state <= ALL_RED_EW;   // 上电先全红
    else
        state <= next_state;
end
\end{lstlisting}

第二段是下一状态逻辑，它就是“状态、时钟与时序逻辑”一章转移表的逐行翻译：

\begin{lstlisting}
// 第二段（上）：下一状态逻辑，纯组合
always_comb begin
    next_state = state;        // 默认保持：timer_done=0 时不动
    case (state)
        EW_GREEN:   if (timer_done) next_state = EW_YELLOW;
        EW_YELLOW:  if (timer_done) next_state = ALL_RED_EW;
        ALL_RED_EW: if (timer_done) next_state = NS_GREEN;
        NS_GREEN:   if (timer_done) next_state = NS_YELLOW;
        NS_YELLOW:  if (timer_done) next_state = ALL_RED_NS;
        ALL_RED_NS: if (timer_done) next_state = EW_GREEN;
        default:    next_state = ALL_RED_EW;   // 非法编码回全红
    endcase
end
\end{lstlisting}

输出段对应“状态、时钟与时序逻辑”一章的输出表。它可以并入上一段，这里分开写，以强调它回答的是另一个问题——“这个状态点亮哪盏灯”：

\begin{lstlisting}
// 第二段（下）：输出逻辑，Moore 型，只由状态决定
always_comb begin
    ew_light = RED;            // 先给默认：双向全红
    ns_light = RED;
    case (state)
        EW_GREEN:   ew_light = GREEN;
        EW_YELLOW:  ew_light = YELLOW;
        NS_GREEN:   ns_light = GREEN;
        NS_YELLOW:  ns_light = YELLOW;
        default: ;             // 两个全红态保持默认
    endcase
end
\end{lstlisting}

逐段读回硬件。\code{state} 是 3 只触发器——六状态二进制编码，与“状态、时钟与时序逻辑”一章的编码表一致；若改用独热方案，把 \code{typedef} 换成 6 位向量加六个单 bit 常量，就是 6 只触发器。\code{next\_state} 的组合网络是一只以 \code{state} 为选择端的 8 输入 mux：注意每个分支若不满足 \code{timer\_done}，mux 的输出就是 \code{state} 自身——“保持”是 mux 把旧值接回 D 端，不是免费的静止。\code{timer\_done} 只出现在 \code{if} 条件里，它是 mux 链上的辅助选择信号。输出段是两个译码器：输入 \code{state}，输出两组灯色；进 \code{case} 前先给默认值、再按需覆盖，保证每种状态下每根输出线都有确定电平，不给锁存器留任何机会（“从 Verilog 读回硬件：综合边界”一节细讲这条纪律）。

三段代码与硬件的对应可以直接画成结构图——它与“状态、时钟与时序逻辑”一章那张三框结构图是同一张图：

\begin{pdfdiagram}
  \node[hw compute, minimum width=2.1cm, minimum height=1.0cm] (nsl) at (0,0) {下一状态逻辑\\always\_comb};
  \node[hw state, minimum width=2.6cm, minimum height=1.0cm] (srg) at (3.4,0) {状态寄存器 state\\always\_ff};
  \node[hw compute, minimum width=2.1cm, minimum height=1.0cm] (oln) at (6.8,0) {输出逻辑\\always\_comb};
  \draw[hw bus] (nsl.east) -- (srg.west);
  \node[hw label] at (1.6,0.68) {next\_state};
  \draw[hw bus] (srg.east) -- (oln.west);
  \node[hw label] at (5.1,0.68) {state};
  \draw[hw bus] (srg.south) -- (3.4,-1.0) -- (0,-1.0) -- (nsl.south);
  \node[hw label, below] at (1.7,-1.0) {state 送回（当前状态）};
  \draw[hw return, dashed] (-2.1,1.35) -- (0,1.35) -- (0,0.53);
  \node[hw label, above] at (-1.05,1.35) {timer\_done};
  \draw[hw return, dashed] (4.1,-1.6) -- (4.1,-0.53);
  \node[hw label, anchor=west] at (4.2,-1.55) {clk};
  \draw[hw bus] (oln.east) -- (9.2,0);
  \node[hw label, below] at (8.6,-0.08) {灯色输出};
\end{pdfdiagram}
\diagramnote{两段式模板的硬件形状：左框（\code{always\_comb}）按 \code{state} 与 \code{timer\_done} 算出 \code{next\_state}，中框（\code{always\_ff}）在 \code{clk} 沿把 \code{next\_state} 提交进状态寄存器，右框（\code{always\_comb}）把 \code{state} 译成 \code{ew\_light}、\code{ns\_light} 两组灯色。\code{state} 经下方通道送回左框——“保持”就是 mux 把旧值接回 D 端。\code{timer\_done} 是这台状态机唯一的转移条件，计数器本体在模块之外。}

用一条逐拍轨迹把三段连起来看。假设 \code{state} 当前是 \code{EW\_GREEN}，定时器还在计数，\code{timer\_done=0}：组合段的默认行生效，算出 \code{next\_state = state}，时钟沿把 \code{state} 原样写回——状态机原地不动，灯色不变。若干秒后定时器减到 0，\code{timer\_done=1} 的那一拍，组合段算出 \code{next\_state = EW\_YELLOW}，沿到来时状态寄存器提交，下一拍起输出译码器看到新状态，东西向由绿转黄。每个周期状态机恰好前进一步或原地不动，这正是“状态、时钟与时序逻辑”一章“状态只在时钟边界提交”的代码形态：

\begin{longtable}{L{0.14\linewidth}L{0.22\linewidth}L{0.22\linewidth}L{0.24\linewidth}}
\toprule
时钟沿 & \code{timer\_done}（沿前） & \code{state}（沿后） & 灯色变化 \\
\midrule
第 $n$ 沿 & 0 & \code{EW\_GREEN}（保持） & 不变 \\\tablerowrule
第 $n+1$ 沿 & 1 & \code{EW\_YELLOW}（提交） & 沿后东西转黄 \\\tablerowrule
第 $n+2$ 沿 & 0 & \code{EW\_YELLOW}（保持） & 不变 \\
\bottomrule
\end{longtable}

还有两个边界要点，都是“状态、时钟与时序逻辑”一章强调过的工程细节在代码里的形状。其一，复位状态是 \code{ALL\_RED\_EW} 而不是某个绿灯态——上电后所有方向先看到红灯。其二，\code{default} 分支把 3 位编码里两个未使用的编码导向全红；删掉这两处，代码更短，路口更危险，而综合器不会替你保留安全余量。最后注意定时器不在状态机里：“状态、时钟与时序逻辑”一章的分工是状态机只回答“现在在哪个阶段、下一个阶段是什么”，计数交给带装载的递减计数器，两者靠 \code{timer\_done}、装载值、启动三条线相连。这个边界在 Verilog 里体现为两个模块各写各的——状态机模板里一个计数语句都不出现。

两段式不是唯一写法。把状态寄存器与下一状态逻辑挤在同一段里的“一段式”也能综合，但状态跳转和输出译码混在一起，读起来要自己在脑子里重新分块；在两段之外再加一段输出寄存器的“三段式”，用一拍延迟换输出无毛刺，适合驱动外部危险负载。对初学者，两段式与“状态、时钟与时序逻辑”一章结构图的对应最整齐：一段文字对应图上一个方框，读代码就是读图。

五个模板与它们读回的硬件，汇总成一张对照表，作为本节的收尾：

\begin{longtable}{L{0.22\linewidth}L{0.32\linewidth}L{0.30\linewidth}}
\toprule
模板 & 读回的硬件 & 敏感表的形状 \\
\midrule
最小寄存器 & 一只 D 触发器 & 只有 \code{posedge clk} \\\tablerowrule
异步复位 & 带 \code{/CLR} 端的触发器（74HC74 式） & 时钟沿加复位沿 \\\tablerowrule
同步复位寄存器 & 触发器加 D 端 mux & 只有 \code{posedge clk} \\\tablerowrule
计数器 & 触发器排加加一器加使能 mux & 只有 \code{posedge clk} \\\tablerowrule
移位寄存器 & 触发器排首尾相接 & 只有 \code{posedge clk} \\\tablerowrule
两段式状态机 & 状态寄存器加两片组合网络 & 只有 \code{posedge clk}（或加复位沿） \\
\bottomrule
\end{longtable}

右列值得多看一眼：除了异步复位，所有模板的敏感表都只有 \code{posedge clk}。同步设计的文字签名就是这么朴素——一个时钟沿统领全部提交。敏感表里多出任何别的东西，都要像对待异步复位那样追问一句：这个边沿对应哪只真实引脚？

本节五个模板共用同一条读法：敏感表指出提交边界，\code{<=} 的两端指出触发器与 D 端网络，\code{if} 链指出 mux 的优先级。把这三样读出来，任何时序 RTL 都能翻回“状态、时钟与时序逻辑”一章的触发器图。Verilog 不是程序，是用文字画电路；时序部分画的，是“沿”与“沿之间”。

\section{四、从 Verilog 读回硬件：综合边界}

前三节解决“怎么写”，本节解决“写完之后综合器怎么想”。综合器不是编译器：编译器把语句翻译成顺序执行的指令流，综合器把文字翻译成门与触发器的网表，而且只接受语言的一个子集。读工程 RTL 必须知道这条边界的走向——哪些写法读回确定的硬件，哪些写法会被推断成意料之外的元件，哪些写法根本不属于 RTL。

在清点之前，先把综合器本身的工作流程摆出来。从 RTL 文字到门级网表，它固定走四步：解析、推断、优化、映射；本节要清点的全部边界问题，都发生在“推断”这一步。

\begin{pdfdiagram}[x=0.85cm]
  \node[hw state, minimum width=1.8cm, minimum height=1.0cm, align=center] (rtl) at (0,0) {RTL 源码};
  \node[hw compute, minimum width=1.7cm, minimum height=1.0cm, align=center] (parse) at (2.2,0) {解析};
  \node[hw compute, minimum width=1.7cm, minimum height=1.0cm, align=center] (infer) at (4.4,0) {推断};
  \node[hw compute, minimum width=1.7cm, minimum height=1.0cm, align=center] (opt) at (6.6,0) {优化};
  \node[hw compute, minimum width=1.7cm, minimum height=1.0cm, align=center] (map) at (8.8,0) {映射到\\目标库};
  \node[hw memory, minimum width=1.8cm, minimum height=1.0cm, align=center] (net) at (11.0,0) {门级\\网表};
  \draw[hw bus] (rtl.east) -- (parse.west);
  \draw[hw bus] (parse.east) -- (infer.west);
  \draw[hw bus] (infer.east) -- (opt.west);
  \draw[hw bus] (opt.east) -- (map.west);
  \draw[hw bus] (map.east) -- (net.west);
  \node[hw label] at (1.1,0.75) {语法树};
  \node[hw label] at (3.3,0.75) {门与触发器};
  \node[hw label] at (5.5,0.75) {化简、共享};
  \node[hw label] at (7.7,0.75) {库单元};
  % 各阶段注记（虚线下挂）
  \draw[hw return, dashed] (rtl.south) -- (0,-0.75);
  \node[hw label, align=center, anchor=north] at (0,-0.85) {只接受\\可综合子集};
  \draw[hw return, dashed] (infer.south) -- (4.4,-0.75);
  \node[hw label, align=center, anchor=north] at (4.4,-0.85) {意外的锁存器\\在此混入};
  \draw[hw return, dashed] (opt.south) -- (6.6,-0.75);
  \node[hw label, align=center, anchor=north] at (6.6,-0.85) {常数折叠\\资源共享};
  \draw[hw return, dashed] (map.south) -- (8.8,-0.75);
  \node[hw label, align=center, anchor=north] at (8.8,-0.85) {工艺库决定\\门数与延迟};
  \node[hw label] at (5.5,-2.2) {边界在输入侧：不能读回门与触发器的写法，走不进这条链};
\end{pdfdiagram}
\diagramnote{综合器的四步流程：解析把源码读成语法树，推断把写法翻成门与触发器（分支没写全时，意外补出的锁存器也在这一步混入），优化做常数折叠与资源共享，映射把逻辑绑定到目标工艺库的单元，最终产出门级网表。本节五类问题里，赋值符号、锁存器推断、位宽与符号约定都落在“推断”一步；\code{\#} 延时、\code{initial}、\codetext{\$display} 这类仿真专属写法不属于可综合子集，走不进这条链。}

\topic{阻塞与非阻塞的正反两例}

“时序逻辑的 Verilog 写法”一节立过规矩：时序块里只写 \code{<=}。本专题用一对正反例把这条规矩算到综合层面，先看它为什么对，再看违反它损失什么。

正例：三级移位寄存器，非阻塞赋值。

\begin{lstlisting}
// 正例：三级移位寄存器，每个 <= 都采旧值
always_ff @(posedge clk) begin
    s1 <= in;
    s2 <= s1;
    s3 <= s2;
end
\end{lstlisting}

读回硬件：三只触发器首尾相接，\code{in} 进第一级的 D，\code{s1} 的 Q 进第二级的 D，\code{s2} 的 Q 进第三级的 D。每个沿所有右端同时采旧值，输入 bit 每拍恰好前进一级。设三级的初值都是 0，逐拍推演：

\begin{longtable}{L{0.16\linewidth}L{0.14\linewidth}L{0.14\linewidth}L{0.14\linewidth}L{0.14\linewidth}}
\toprule
沿序号 & \code{in}（沿前） & \code{s1}（沿后） & \code{s2}（沿后） & \code{s3}（沿后） \\
\midrule
第 1 沿 & 1 & 1 & 0 & 0 \\\tablerowrule
第 2 沿 & 0 & 0 & 1 & 0 \\\tablerowrule
第 3 沿 & 1 & 1 & 0 & 1 \\
\bottomrule
\end{longtable}

“每拍前进一级”不是文字游戏，是硬件结构决定的：第 $i$ 级的 D 端物理上接着第 $i-1$ 级的 Q，沿到来时第 $i-1$ 级的新值还没来得及穿过时钟到输出的延迟，第 $i$ 级采到的必然是旧值。“状态、时钟与时序逻辑”一章讲时钟到输出延迟与保持时间时画的沿口竞态，正是这套语义的物理底板。仿真器兑现“旧值语义”的方式也值得一提：它先把块内所有 \code{<=} 的右端按当前值算好、暂存，再把左端一起更新——计算有先后，生效同时刻，与硬件的沿口提交一一对应。

反例：同样三行，换成阻塞赋值。

\begin{lstlisting}
// 反例：= 立即生效，新值沿书写顺序连跳三级
always_ff @(posedge clk) begin
    s1 = in;
    s2 = s1;
    s3 = s2;
end
\end{lstlisting}

逐句追踪：第一句 \code{s1 = in} 立即生效，\code{s1} 当场变成 \code{in} 的新值；第二句读到的 \code{s1} 已是新值，\code{s2} 跟着变成 \code{in}；第三句同理——同一个沿里，\code{in} 的值沿书写顺序连跳三级，直达 \code{s3}。硬件读法：三只触发器的 D 端全接到了 \code{in}，三级移位退化成三只并联的单级采样器。书写顺序在这里成了硬件的一部分：三行换个排列，推断出的接线就换一副面孔。

两种写法的接线差别，画成图就是两副完全不同的面孔：

\begin{pdfdiagram}
  % 正例：<= 三级移位链
  \node[hw label] at (2.4,4.55) {正例：三行 \code{<=}，三只触发器首尾相接};
  \node[hw state, minimum width=1.1cm, minimum height=0.8cm] (g1) at (0,3.6) {s1};
  \node[hw state, minimum width=1.1cm, minimum height=0.8cm] (g2) at (2.4,3.6) {s2};
  \node[hw state, minimum width=1.1cm, minimum height=0.8cm] (g3) at (4.8,3.6) {s3};
  \node[hw label, anchor=east] at (-1.7,3.6) {in};
  \draw[hw bus] (-1.6,3.6) -- (g1.west);
  \draw[hw bus] (g1.east) -- (g2.west);
  \draw[hw bus] (g2.east) -- (g3.west);
  \draw[hw bus] (g3.east) -- (6.6,3.6);
  \node[hw label, anchor=west] at (6.7,3.6) {三拍后到达};
  \draw[BookTeal, dashed, line width=0.58pt] (-1.6,2.95) -- (5.4,2.95);
  \draw[hw return, dashed] (0,2.95) -- (0,3.17);
  \draw[hw return, dashed] (2.4,2.95) -- (2.4,3.17);
  \draw[hw return, dashed] (4.8,2.95) -- (4.8,3.17);
  \node[hw label, anchor=east] at (-1.7,2.95) {clk};
  % 反例：= D 端并联接 in
  \node[hw label] at (2.4,2.1) {反例：三行 \code{=}，D 端并联接 in，同拍直达};
  \node[hw state, minimum width=1.1cm, minimum height=0.8cm] (h1) at (0,0.4) {s1};
  \node[hw state, minimum width=1.1cm, minimum height=0.8cm] (h2) at (2.4,0.4) {s2};
  \node[hw state, minimum width=1.1cm, minimum height=0.8cm] (h3) at (4.8,0.4) {s3};
  \node[hw label, anchor=east] at (-1.7,0.4) {in};
  \draw[BookMuted, line width=0.62pt] (-1.6,0.4) -- (-1.0,0.4) -- (-1.0,1.4) -- (4.8,1.4);
  \draw[hw bus] (0,1.4) -- (0,0.82);
  \draw[hw bus] (2.4,1.4) -- (2.4,0.82);
  \draw[hw bus] (4.8,1.4) -- (4.8,0.82);
  \fill[BookMuted] (0,1.4) circle (0.04);
  \fill[BookMuted] (2.4,1.4) circle (0.04);
  \node[hw label, above] at (1.9,1.4) {in 同拍到达三只 D};
  \draw[hw bus] (h3.east) -- (6.6,0.4);
  \node[hw label, anchor=west] at (6.7,0.4) {同拍直达};
  \draw[BookTeal, dashed, line width=0.58pt] (-1.6,-0.35) -- (5.4,-0.35);
  \draw[hw return, dashed] (0,-0.35) -- (0,-0.02);
  \draw[hw return, dashed] (2.4,-0.35) -- (2.4,-0.02);
  \draw[hw return, dashed] (4.8,-0.35) -- (4.8,-0.02);
  \node[hw label, anchor=east] at (-1.7,-0.35) {clk};
\end{pdfdiagram}
\diagramnote{同一个 \code{always\_ff} 块的两种硬件读法。上：三行 \code{<=} 都采沿前旧值，综合出首尾相接的三级移位链，输入 bit 每拍恰好前进一级，三拍后到达 \code{s3}。下：三行 \code{=} 立即生效，等效于三只触发器的 D 端并联接 \code{in}——同拍直达 \code{s3}，三级移位退化成一级采样。写法差三个字符，硬件差两级触发器。}

这种错误最阴险的地方在于，它可能通过粗略的功能检查。若测试只看“输入最终出现在输出”，反例一样能做到——错的只是拍数：本该三拍后到达，实际一拍直达。而拍数恰恰是硬件语义的一部分：“状态、时钟与时序逻辑”一章给状态机画对齐波形时逐周期核对控制线，核对的就是这种东西。写法差一个字符，硬件差两级触发器；仿真里差一个字符，行为差两级流水。

那么组合块为什么反其道而行，用 \code{=}？因为 \code{always\_comb} 描述的是导线网络，电平必须立即传播才能正确建模级联的门。“组合逻辑的 Verilog 写法”一节讲过组合块可以引入临时变量分步计算——先算中间量、再算输出——与“与门输出接或门输入”一一对应；若改用 \code{<=}，等于声明“中间量要等某个沿才生效”，既不符合导线的事实，也会让分支完备性检查错乱。组合块里本来就没有“沿之前的旧值”这个概念，\code{<=} 的旧值语义无从谈起。

时序块里确实需要中间量时，工程写法是把两类赋值分开用：临时变量用 \code{=} 在块内当场算清，真正进触发器的信号用 \code{<=} 提交。临时变量综合后是 D 端组合网络的一部分，不进任何触发器——只在块内当场消费的变量没有独立的硬件对应物，它只是表达式的一种写法。

记忆规则一句话：时序用 \code{<=}，组合用 \code{=}。物理依据也是一句话：\code{<=} 对应“触发器在沿处同时提交”，\code{=} 对应“导线上的电平立即传播”。两个赋值符号不是两种执行速度，而是两类硬件。

组合块内部对 \code{=} 的用法也有一对正反：先算后用，中间量是一根跳线；先读后写，就得为“上一次求值的旧值”多造一只锁存器。并排画出：

\begin{pdfdiagram}[x=0.85cm]
  % 上：正确用法，先算后用
  \node[hw label] at (1.0,4.6) {正确：先算后用，\code{t} 是中间跳线};
  \node[anchor=west, font=\scriptsize\ttfamily, text=BookInk] at (-5.4,3.8) {t = a \& b;};
  \node[anchor=west, font=\scriptsize\ttfamily, text=BookInk] at (-5.4,3.1) {y = t | c;};
  \node[hw compute, minimum width=1.2cm, minimum height=0.9cm] (g1) at (0.0,3.4) {AND};
  \node[hw compute, minimum width=1.2cm, minimum height=0.9cm] (g2) at (2.3,3.4) {OR};
  \node[hw label, anchor=east] at (-1.45,3.62) {a};
  \node[hw label, anchor=east] at (-1.45,3.18) {b};
  \draw[hw bus] (-1.3,3.62) -- ($(g1.west)+(0,0.22)$);
  \draw[hw bus] (-1.3,3.18) -- ($(g1.west)+(0,-0.22)$);
  \draw[hw bus] (g1.east) -- (g2.west);
  \node[hw label, above] at (1.15,3.42) {t};
  \node[hw label] at (2.3,2.6) {c};
  \draw[hw bus] (2.3,2.8) -- (g2.south);
  \draw[hw bus] (g2.east) -- (4.1,3.4);
  \node[hw label, anchor=west] at (4.2,3.4) {y};
  \node[hw label, anchor=north] at (1.0,2.25) {综合结果：两只门，无存储};
  % 下：误用，先读后写
  \node[hw label] at (1.4,1.3) {误用：先读后写，\code{t} 要保持上一次求值的旧值};
  \node[anchor=west, font=\scriptsize\ttfamily, text=BookInk] at (-5.4,0.4) {y = t | c;   // t 是旧值};
  \node[anchor=west, font=\scriptsize\ttfamily, text=BookInk] at (-5.4,-0.3) {t = a \& b;   // 下一遍可见};
  \node[hw compute, minimum width=1.2cm, minimum height=0.9cm] (h1) at (0.0,0.0) {AND};
  \node[hw memory, minimum width=1.9cm, minimum height=0.9cm, align=center] (lt) at (2.1,0.0) {被推断的\\锁存器};
  \node[hw compute, minimum width=1.2cm, minimum height=0.9cm] (h2) at (4.3,0.0) {OR};
  \node[hw label, anchor=east] at (-1.45,0.22) {a};
  \node[hw label, anchor=east] at (-1.45,-0.22) {b};
  \draw[hw bus] (-1.3,0.22) -- ($(h1.west)+(0,0.22)$);
  \draw[hw bus] (-1.3,-0.22) -- ($(h1.west)+(0,-0.22)$);
  \draw[hw bus] (h1.east) -- (lt.west);
  \draw[hw bus] (lt.east) -- (h2.west);
  \node[hw label, above] at (3.2,0.02) {t};
  \node[hw label] at (4.3,-0.85) {c};
  \draw[hw bus] (4.3,-0.65) -- (h2.south);
  \draw[hw bus] (h2.east) -- (6.1,0.0);
  \node[hw label, anchor=west] at (6.2,0.0) {y};
  \node[hw label, anchor=north] at (2.1,-1.35) {语句集合相同、次序颠倒，硬件多出一只电平敏感存储};
\end{pdfdiagram}
\diagramnote{组合块里 \code{=} 的正反两例。上：先写 \code{t} 再用它，\code{t} 综合成 AND 与 OR 之间的一根中间跳线，块内没有任何存储。下：先读 \code{t} 后写 \code{t}，读到的只能是“上一遍求值留下的旧值”——保持旧值需要存储，综合器为 \code{t} 补出一只电平敏感锁存器，透明窗口与毛刺穿透随之而来。同样的语句换个次序，硬件多出一种元件；组合块里被读的信号，永远先写后读。}

反例也值得画一张逐拍表。同样的初值、同样的输入序列，阻塞版本的三级在每个沿后全部等于当时的 \code{in}：

\begin{longtable}{L{0.16\linewidth}L{0.14\linewidth}L{0.14\linewidth}L{0.14\linewidth}L{0.14\linewidth}}
\toprule
沿序号 & \code{in}（沿前） & \code{s1}（沿后） & \code{s2}（沿后） & \code{s3}（沿后） \\
\midrule
第 1 沿 & 1 & 1 & 1 & 1 \\\tablerowrule
第 2 沿 & 0 & 0 & 0 & 0 \\\tablerowrule
第 3 沿 & 1 & 1 & 1 & 1 \\
\bottomrule
\end{longtable}

两张表放在一起，正反例的差别一眼可见：正例的 \code{s3} 是 \code{in} 延迟三拍的副本，反例的 \code{s3} 就是 \code{in} 本身——一个是移位寄存器，一个是一根导线带三只并联触发器。写代码时二者只差三个字符，读代码时若扫一眼赋值符号就翻页，这个差别就漏过去了。所以读时序块的第一眼，永远先落在赋值符号上。

还有一种更隐蔽的混用：同一个时序块里，一部分信号用 \code{<=}、另一部分信号用 \code{=}。仿真器按各自规则照单处理，结果往往“大体像对的”——错位只出现在混用处附近的拍序上。更麻烦的是仿真与综合对书写顺序的敏感度不完全一致：仿真严格按顺序执行阻塞赋值，而综合器推断组合前级时并不承诺维持这份顺序语义，于是仿真通过、硬件出错，成为最难追查的一类缺陷。工程评审把“时序块里出现 \code{=}”列为必查项，查的不是风格，是这块硬件里可能多了一根直通导线、少了一级流水。

早年资料里还有一种说法——“非阻塞赋值是并行执行”。它并不准确：仿真器照样一句一句地算，只是左端的生效时刻被推迟到整个块算完之后。把 \code{<=} 理解成“延迟到沿尾生效的赋值”也不算错，但都不如直接读成“D 端接线”来得干净：接线没有执行，只有连接。

\topic{锁存器推断的产生与消除}

“组合逻辑的 Verilog 写法”一节已经看过推断的产生：组合块里某条路径没给输出赋值，综合器面对“保持原值”的规格，只好搬出“状态、时钟与时序逻辑”一章那只电平敏感的锁存器来填空。本专题把视角挪到综合边界上，回答三个相关的问题：时序块为什么天然免疫，综合警告怎么读，除了默认值和补全分支还有什么手段。

先回答免疫问题。时序块里“没写”是常态：“时序逻辑的 Verilog 写法”一节的计数器只在 \code{en} 有效时加一，其余分支一字不写，语义是“保持旧值”——而保持是触发器的本能，D 端经一只 mux 接回 Q 即可，不需要任何新元件。同一个“没写”，在时序块里是免费的本能，在组合块里却要多造一只锁存器，差别全在信号背后有没有记忆元件。对照如下：

\begin{longtable}{L{0.15\linewidth}L{0.31\linewidth}L{0.22\linewidth}L{0.20\linewidth}}
\toprule
块类型 & 某分支没给信号赋值 & 综合结果 & 处置 \\
\midrule
\code{always\_ff} & 保持旧值，D 端经 mux 接回 Q & 触发器加使能，通常正是本意 & 无需处理 \\\tablerowrule
\code{always\_comb} & 语义要求保持，组合网络却无记忆 & 推断锁存器，几乎从非本意 & 默认值、补全、\code{unique} 三选一 \\
\bottomrule
\end{longtable}

所以锁存器推断有一个可靠的判别式：它只可能发生在 \code{always\_comb}（或电平敏感的 \code{always}）里，永远不可能发生在 \code{always\_ff @(posedge clk)} 里。读综合报告时若看到锁存器警告指向一个自认是时序逻辑的块，那一定是敏感表写成了电平敏感，或者块内混进了组合馈通——先查敏感表，再查赋值符号。

综合警告怎么读？典型的锁存器警告会给出信号名与位置，含义是“综合器为这个信号推断了一只锁存器”。读它抓三个要素：哪个信号、哪个块、哪条路径漏了赋值。这类警告不是可以直接忽略的噪音——每一条都对应真实硬件里多出来的一只透明锁存器，带着“状态、时钟与时序逻辑”一章分析过的全部麻烦：使能窗口内输入直通、毛刺穿透、时序分析从周期边界变成电平窗口。工程上的纪律是：锁存器警告数必须为零，或者每一处都有写下来的理由。

消除手段，“组合逻辑的 Verilog 写法”一节给了两种——块首写默认值、分支补全——这里补上第三种：用 \code{unique} 把分支性质变成工具可检查的声明。

\begin{lstlisting}
// unique：声明分支互斥且覆盖所有取值
always_comb begin
    unique case (sel)
        2'b00: y = d0;
        2'b01: y = d1;
        2'b10: y = d2;
        2'b11: y = d3;
    endcase
end
\end{lstlisting}

\code{unique case} 向综合器声明两件事：各分支互斥，且覆盖所有可能取值。声明属实，综合器按并行译码放手优化，遇到遗漏会警告而不是静默推断锁存器；声明不属实——例如实际出现了没列出的编码——仿真器会在运行期报错。它把“我已经检查过分支完备性”这句口头承诺，变成仿真与综合都能核对的约定。若确实需要优先级语义、分支可能重叠，对应的关键字是 \code{priority}：它声明“按书写顺序取第一个匹配”，同样禁止锁存器推断，但保留 if-else 链的次序含义。

最后留一道边界题：有没有故意要锁存器的场合？有——某些超低功耗或特殊时序设计会显式使用电平敏感存储，但那属于高级技法，会用专门的写法与评审流程把它标记出来。对本书读者，规则不变：组合块的输出必须在每条路径上有确定值，锁存器警告数为零。

分支不全也不止 \code{case} 缺 \code{default} 一种形状，\code{if} 缺 \code{else} 同样漏风：

\begin{lstlisting}
// if 缺 else：sel=0 时 y 保持，推断锁存器
always_comb begin
    if (sel)
        y = a;
end
\end{lstlisting}

读法与 \code{case} 缺口的读法一致：顺着每条执行路径问一句“\code{y} 在这条路上被写了吗”，只要有一条路的答案是否定，锁存器就进门。形状越朴素越容易被漏看——三行的块在评审里往往比三十行的块更危险。

把推断出的锁存器与触发器并排看，差别回到“状态、时钟与时序逻辑”一章的透明窗口：触发器只在沿上采样，沿与沿之间输出是安静的；锁存器在整个使能电平期间透明，输入毛刺会原样穿出。时序分析也随之分裂：触发器路径按周期边界算，锁存器路径要按电平窗口算，而工具的默认检查流程大多按前者设计——这就是“意外锁存器”常常同时带来功能风险与验证盲区的原因。

顺带划清一个名词边界：本节说的是“推断”——工具替你补出来的锁存器。若设计真的需要电平敏感存储，应当显式实例化并写明意图，让它在评审与报告里以本来面目出现，而不是作为某段组合代码的副产品悄悄混入。另外，不同综合器的警告措辞并不统一，有的写推断锁存器，有的写某信号可能被保持；读陌生工具的报告时，先把它的锁存器措辞认下来，再谈清零纪律。

\topic{不可综合写法}

Verilog 语言先于综合工具而存在，它的全集是写给仿真器的；RTL 只是全集中“能读回门与触发器”的子集。判断一个写法是否在子集里，只需问一句：它对应哪种真实存在的硬件结构？以下几类常见写法通不过这一问，逐一看原因。

延时语句 \code{\#10}。它的含义是“仿真时间前进 10 个单位”。真实电路里没有“等待 10 个时间单位”的元件；延迟由门延迟与布线寄生决定——那是“电源、时钟与信号完整性”一章的话题——不由代码决定。混进 RTL 的延时语句，综合器有的报错、有的悄悄忽略。忽略比报错更危险：仿真里存在的延迟在硬件里凭空消失，仿真与实物就此分叉，而工具一言不发。

初始化块 \code{initial}。它描述“仿真开始那一瞬间的值”。ASIC 的触发器上电取值不定，\code{initial} 没有硬件对应；多数 FPGA 的触发器带配置上电初值机制，这些平台的综合器因此接受 \code{initial}，把它变成上电初值——但这是平台特性，不是语言承诺。要跨平台移植的代码把初始化交给复位（“时序逻辑的 Verilog 写法”一节的两种写法），不依赖 \code{initial}。读 FPGA 专用代码时见到 \code{initial}，要明白它读回的是“配置完成时刻的触发器初值”，而不是任何运转中的硬件。

打印系统任务。\codetext{\$display} 往仿真终端写一行文本，\codetext{\$monitor} 跟踪信号变化并自动打印，\codetext{\$fatal} 报错并终止仿真。硬件里没有打印机，也没有终端，这些调用没有任何门级对应物。“testbench 与仿真阅读”一节会看到它们在 testbench 里大显身手——那里才是它们的合法住所。

实数变量 \code{real}。浮点变量没有位宽，综合器不知道给它分配几只触发器。硬件里的“小数”是定点数：算术与数据通路相关章节的 Q 格式，用一个整数加上约定的二进制小数点位置来表示，所有运算仍是整数运算。\code{real} 只可能出现在仿真侧，或编译期常量折叠的表达式里。

那么，为什么 testbench 可以尽情使用它们？因为 testbench 没有综合义务。它是仿真世界里的虚拟实验台：产生时钟、驱动激励、观察响应、打印报告，这一切只发生在仿真器内部，永远不需要变成硅片或查找表。同一语言、两个读者——RTL 写给综合器，testbench 写给仿真器。工程目录通常把 \filepath{rtl/} 与 \filepath{tb/} 分开放置，读代码时先看文件在哪个目录，就知道该换上哪种读法。反过来的判别同样好用：在声称要综合的文件里看到 \codetext{\$display}，基本可以断定作者没打算把它做成硬件。

还要区分第三类文件：仿真模型。存储器厂商提供的模型、晶振的行为级描述，里面满是延时与 \code{initial}——它们同样不综合，用途是在仿真里扮演真实器件，好让你的 RTL 有逼真的对手。它们介于 RTL 与 testbench 之间：不比 testbench 更接近硅片，也不比 RTL 更远离。下表汇总这几类写法的边界：

\begin{longtable}{L{0.17\linewidth}L{0.29\linewidth}L{0.24\linewidth}L{0.18\linewidth}}
\toprule
写法 & 硬件里不存在的东西 & 综合器反应 & 合法出现场合 \\
\midrule
\code{\#10} 延时 & “等待 10 个时间单位”的元件 & 报错或悄悄忽略 & testbench \\\tablerowrule
\code{initial} & ASIC 触发器的确定上电值 & 多数报错；FPGA 流程常收作上电初值 & testbench、FPGA 专用 RTL \\\tablerowrule
\codetext{\$display} 一族 & 打印机与终端 & 忽略或报错 & testbench \\\tablerowrule
\code{real} & 无位宽的浮点存储 & 报错 & testbench、常量表达式 \\
\bottomrule
\end{longtable}

把这一专题收拢成一句话：仿真器执行的是语言，综合器接受的是结构。凡是不能在脑子里展开成门、触发器与连线的写法，都不属于 RTL——它要么属于 testbench，要么属于仿真模型，要么属于还没人发明出来的硬件。

阅读顺序也因此确定：拿到一份陌生代码，先做分拣——哪些是 RTL，哪些是测试与模型——再对 RTL 部分动用本章的读法。把 testbench 的 \code{initial} 当成设计缺陷去改，与把 RTL 的锁存器警告当成风格问题去忍，是同一种错位的两个方向。

\topic{signed 与位宽陷阱}

算术与数据通路相关章节讲过：同一组 bit，按无符号读和按补码读是两个数。“组合逻辑的 Verilog 写法”一节的运算符表默认了“读数方式已经明确”，本专题处理这张表底下的暗礁。Verilog 的运算符对两种读法共用一套符号，歧义全靠声明消解——综合出来的是补零移位器还是补符号移位器、比较器按哪种数轴比大小，全看声明。三个要点分开说。

第一，\code{>>} 与 \code{>>>}。逻辑右移 \code{>>} 高位补 0；算术右移 \code{>>>} 高位补符号位——但 \code{>>>} 只在操作数声明为 \code{signed} 时才真正是算术移位，对无符号操作数写 \code{>>>} 仍然补 0。移位器补什么由声明决定，不由符号的长相决定：

\begin{lstlisting}
logic        [7:0] u = 8'b1100_0000;   // 无符号：192
logic signed [7:0] s = 8'b1100_0000;   // 补码：-64

u >>  2    // 0011_0000：高位补 0，得 48
u >>> 2    // 0011_0000：无符号，仍补 0
s >>> 2    // 1111_0000：补符号位，得 -16
\end{lstlisting}

第二，\codetext{\$signed} 与比较时的扩展规则。\codetext{\$signed(x)} 强制把 \code{x} 按补码解释。比较运算的规则是：两个操作数都是 \code{signed}，先符号扩展到同宽，再按补码比较；只要有一个是 \code{unsigned}，两个都按无符号处理。最后这条是著名的反直觉规则，给出一个容易出错的示例：

\begin{lstlisting}
logic        [3:0] a = 4'b1100;   // 无符号：读作 12
logic signed [3:0] b = 4'b0111;   // signed：读作 +7

if (a < b)          // a 无符号，b 被当成无符号
                    // 实际比较 12 < 7，不成立

if ($signed(a) < b) // 两边同按补码：-4 < 7，成立
\end{lstlisting}

逐行算：\code{a} 是 \code{4'b1100}，按无符号读作 12；\code{b} 声明了 \code{signed}，\code{4'b0111} 读作 +7。第一个比较里 \code{a} 是无符号，于是 \code{b} 也被当成无符号，实际比较是 $12<7$，不成立。若设计本意是 \code{a} 存放补码 $-4$，必须写成 \codetext{\$signed(a) < b}，两边同按补码，$-4<7$ 成立。少写一个 \codetext{\$signed}，比较结果整个翻转——而仿真不会报任何错，综合器也不会，因为对工具来说两种读法都“合法”。符号与位宽是阅读 RTL 时必须亲手核对的隐性约定，工具默认你清楚自己在写什么。

第三，位宽与字面量。表达式的运算位宽按参与者的最大宽度进行，赋值时再截断或扩展到左端宽度。“时序逻辑的 Verilog 写法”一节计数器模板里写 \code{q + 1'b1} 而不写 \code{q + 1}，原因就在此：不带位宽的十进制字面量是 32 位有符号数，\code{q + 1} 会先把 \code{q} 扩展到 32 位再相加，赋回 4 位的 \code{q} 时虽然会截断，但中间过程一旦接进更宽的表达式，差异就会露头：

\begin{lstlisting}
logic [3:0] q = 4'd15;
logic [3:0] n;
logic [7:0] wide;

n    = q + 1'b1;   // 语境宽 4 位：和为 1_0000，进位被截断，n = 0
wide = q + 1'b1;   // 语境宽 8 位：q 先零扩展，wide = 16
\end{lstlisting}

同一笔 \code{q + 1'b1}，赋给 4 位的 \code{n} 得 0，赋给 8 位的 \code{wide} 得 16，差别全在左端的接收宽度。语境宽度规则要记全：表达式的运算位宽按语境内所有参与者的最大宽度取定，赋值左端也算参与者；算完再按左端宽度截断或扩展。所谓“进位丢弃”，不是加法器没算出第 5 位，而是接收端只有 4 位、进位在赋回时被截掉——想留住进位，接收端必须够宽。Harris \& Harris 的纪律值得照抄：字面量永远带位宽与基数（\code{4'b1100}、\code{8'd30}），算术表达式每一段的宽度都要心里有数。三个陷阱的共同根源是一个：Verilog 不会替你做“这组 bit 是什么数、有多宽”的决定，它把决定权连同出错的自由一起交给写法。读别人的 RTL，见到算术运算先核对每个操作数的 \code{signed} 声明与位宽，再谈逻辑对错。

符号扩展与零扩展的方向也要分清。把窄信号赋给宽信号时，\code{signed} 窄信号做符号扩展——高位复制符号位；\code{unsigned} 窄信号做零扩展。一次赋值里，读数方式与扩展方式都由声明牵着走：

\begin{lstlisting}
logic signed [3:0] s4 = 4'b1100;   // 补码：-4
logic        [3:0] u4 = 4'b1100;   // 无符号：12
logic        [7:0] w1, w2;

w1 = s4;   // 符号扩展：1111_1100，按无符号读是 252
w2 = u4;   // 零扩展：0000_1100，仍是 12
\end{lstlisting}

\code{w1} 的 252 往往让第一次遇到的人愣住：赋进去的是 $-4$，读出来怎么成了 252？答案是一切都按规则发生——\code{w1} 是无符号的 8 位向量，$-4$ 的补码按无符号读正是 252。位没有丢，读数方式变了。若接收端也声明 \code{signed}，读回 $-4$，这次扩展就“透明”了。可见扩展本身从不出错，出错的是两端的读数方式不一致。

比较器的硬件也因此分两路。无符号比较沿“组合逻辑的 Verilog 写法”一节说的减法路径看最高位借位；有符号比较同样做减法，判据却换成“差的符号位连同溢出一起判定”——算术与数据通路相关章节补码减法判正负的那套逻辑。两种比较器面积相近，但接错声明，等于把数轴的原点挪了位置：无符号的 \code{4'b1100} 在数轴右端（12），有符号的它紧挨原点左侧（$-4$）。同一个数轴、两个原点，是比较陷阱的几何图像。

还有一处与位宽有关的阅读细节：拼接与切片永远指名道姓。拼接结果的宽度是各段宽度之和，切片抽出的宽度是下标区间——它们不随语境伸缩，因此是最安全的位宽操作。真正需要提防的是隐式部分：运算扩展、赋值截断、比较对齐，这三处都按规则默默进行；规则本身不复杂，复杂的是读代码的人要记得它们存在。

工程代码里常见的防御写法也值得认得：比较前把两个操作数显式调成同一宽度同一符号，或都用 \codetext{\$signed} 包起来；计数比较能用 \code{==} 就不用 \code{>=}——“组合逻辑的 Verilog 写法”一节算过账，等值比较比大小比较便宜；字面量一律带位宽。这些写法不改变硬件，改变的是“读者需要猜的东西”的数量。

\topic{generate 与参数化}

处理器核心相关章节最小 CPU 的数据通路、教学计算机项目中的面包板机的总线，都是“同一结构复制 N 份”的产物：8 只触发器排成寄存器，8 个全加器排成并行加法器。“从电路图到 RTL：两种描述同一个硬件”一节讲过实例化是把画好的图放进另一张图——那是手动复制；本专题的 generate 是自动复制。先把最重要的话说在前面：generate 是编译期的复制，不是运行时的循环。它在综合开始之前就被展开成 N 份独立硬件，电路上没有任何“循环”在运行。

\code{parameter} 与 \code{localparam} 把尺寸做成参数：

\begin{lstlisting}
module shift_reg #(
    parameter int WIDTH = 8        // 实例化时可覆盖
) (
    input  logic             clk,
    input  logic             en,
    input  logic             sin,
    output logic [WIDTH-1:0] q
);
    localparam int LAST = WIDTH - 1;   // 派生尺寸，内部常量

    always_ff @(posedge clk) begin
        if (en)
            q <= {q[LAST-1:0], sin};
    end
endmodule
\end{lstlisting}

\code{parameter} 在实例化时可被覆盖，同一个模块描述因此能生成 8 位、16 位、32 位的移位寄存器；\code{localparam} 是模块内部常量，用来给派生尺寸命名，外部改不到。算术与数据通路相关章节的加法器位宽、处理器核心相关章节的数据通路宽度、经典芯片与平台案例各芯片的地址线数，都是 \code{parameter} 的天然用例——同一个设计，改一个数字就换一种规模，硬件随文字参数重新生成。

状态机的编码也可以用同一组机制收回纸面。“时序逻辑的 Verilog 写法”一节的交通灯模板用 \code{typedef enum}，把具体编码交给综合器选择；若设计要求编码确定——例如按“状态、时钟与时序逻辑”一章的格雷方案减少翻转时的多位同时变化——就用 \code{localparam logic [2:0] EW\_GREEN = 3'b000;} 这样的常量逐个指定。许多综合器还提供状态编码属性（写在声明旁的专用属性标注），作用相同：把“用哪种编码”从综合器的自由裁量变成设计者的显式决定。属性本身不改变仿真语义，只是递给综合器的便签；读代码时把它当作有工具效力的注释即可。

以交通灯的六个状态为例，三种常用编码的位模式并排对照，位宽与翻转特性的差别一眼可读：

\begin{pdfdiagram}
  \node[hw label] at (1.23,3.85) {二进制（3 位）};
  \node[hw label] at (3.53,3.85) {格雷（3 位）};
  \node[hw label] at (5.66,3.85) {独热（6 位）};
  % S0：bin 000，gray 000，onehot 000001
  \node[hw label, anchor=east] at (0.28,3.2) {S0};
  \foreach \x/\b in {0.6/0, 1.02/0, 1.44/0, 2.6/0, 3.02/0, 3.44/0, 4.4/0, 4.82/0, 5.24/0, 5.66/0, 6.08/0, 6.5/1} {
    \draw[BookTeal, line width=0.6pt, fill=BookCodeBg] (\x,2.95) rectangle ({\x+0.42},3.45);
    \node at ({\x+0.21},3.2) {\b};
  }
  % S1：bin 001，gray 001，onehot 000010
  \node[hw label, anchor=east] at (0.28,2.58) {S1};
  \foreach \x/\b in {0.6/0, 1.02/0, 1.44/1, 2.6/0, 3.02/0, 3.44/1, 4.4/0, 4.82/0, 5.24/0, 5.66/0, 6.08/1, 6.5/0} {
    \draw[BookTeal, line width=0.6pt, fill=BookCodeBg] (\x,2.33) rectangle ({\x+0.42},2.83);
    \node at ({\x+0.21},2.58) {\b};
  }
  % S2：bin 010，gray 011，onehot 000100
  \node[hw label, anchor=east] at (0.28,1.96) {S2};
  \foreach \x/\b in {0.6/0, 1.02/1, 1.44/0, 2.6/0, 3.02/1, 3.44/1, 4.4/0, 4.82/0, 5.24/0, 5.66/1, 6.08/0, 6.5/0} {
    \draw[BookTeal, line width=0.6pt, fill=BookCodeBg] (\x,1.71) rectangle ({\x+0.42},2.21);
    \node at ({\x+0.21},1.96) {\b};
  }
  % S3：bin 011，gray 010，onehot 001000
  \node[hw label, anchor=east] at (0.28,1.34) {S3};
  \foreach \x/\b in {0.6/0, 1.02/1, 1.44/1, 2.6/0, 3.02/1, 3.44/0, 4.4/0, 4.82/0, 5.24/1, 5.66/0, 6.08/0, 6.5/0} {
    \draw[BookTeal, line width=0.6pt, fill=BookCodeBg] (\x,1.09) rectangle ({\x+0.42},1.59);
    \node at ({\x+0.21},1.34) {\b};
  }
  % S4：bin 100，gray 110，onehot 010000
  \node[hw label, anchor=east] at (0.28,0.72) {S4};
  \foreach \x/\b in {0.6/1, 1.02/0, 1.44/0, 2.6/1, 3.02/1, 3.44/0, 4.4/0, 4.82/1, 5.24/0, 5.66/0, 6.08/0, 6.5/0} {
    \draw[BookTeal, line width=0.6pt, fill=BookCodeBg] (\x,0.47) rectangle ({\x+0.42},0.97);
    \node at ({\x+0.21},0.72) {\b};
  }
  % S5：bin 101，gray 111，onehot 100000
  \node[hw label, anchor=east] at (0.28,0.10) {S5};
  \foreach \x/\b in {0.6/1, 1.02/0, 1.44/1, 2.6/1, 3.02/1, 3.44/1, 4.4/1, 4.82/0, 5.24/0, 5.66/0, 6.08/0, 6.5/0} {
    \draw[BookTeal, line width=0.6pt, fill=BookCodeBg] (\x,-0.15) rectangle ({\x+0.42},0.35);
    \node at ({\x+0.21},0.10) {\b};
  }
  % 组注记
  \node[hw label, anchor=north, align=center] at (1.23,-0.55) {3 只触发器\\相邻转移可多位同变};
  \node[hw label, anchor=north, align=center] at (3.53,-0.55) {3 只触发器\\相邻态只变 1 位};
  \node[hw label, anchor=north, align=center] at (5.66,-0.55) {6 只触发器\\译码最简、数量最多};
\end{pdfdiagram}
\diagramnote{六状态机三种编码的位模式对照（S0 到 S5 依次对应 \code{EW\_GREEN}、\code{EW\_YELLOW}、\code{ALL\_RED\_EW}、\code{NS\_GREEN}、\code{NS\_YELLOW}、\code{ALL\_RED\_NS}）。二进制 3 位最省触发器，但相邻状态的转移可能多位同时翻转；格雷同样 3 位，相邻状态只差 1 位，翻转少、译码毛刺少；独热每个状态一只触发器共 6 只，下一状态与输出译码最简，代价是触发器数量翻倍。“时序逻辑的 Verilog 写法”一节模板用 \code{typedef enum} 时选择权在综合器；用 \code{localparam} 常量或编码属性，可以把它收回纸面。}

\code{generate for} 复制结构。以 N 组独立使能的寄存器堆为例：

\begin{lstlisting}
genvar i;
generate
    for (i = 0; i < DEPTH; i = i + 1) begin : gen_regs
        // 每一份复制都是一组独立的触发器
        always_ff @(posedge clk) begin
            if (we[i])
                regs[i] <= wdata;
        end
    end
endgenerate
\end{lstlisting}

读回硬件：综合后得到 \code{DEPTH} 组完全相同的触发器排，每组一个使能端，与把单组代码手工抄 \code{DEPTH} 遍一字不差。\code{begin} 后的标号 \code{gen\_regs} 不是装饰：它给每份复制命名，层次路径 \code{gen\_regs[3]} 会出现在仿真波形和综合报告里，用来定位第 3 份——没有标号，工具只能给自动生成一串无意义的编号。

把寄存器堆按 8 份展开画出来，“编译期复制”的含义就成了图上的并排八份——每份都是独立的一组触发器，层次路径各不相同。

\begin{pdfdiagram}[x=0.85cm]
  \node[hw control, minimum width=3.4cm, minimum height=2.4cm, align=center] (gen) at (0,0.6) {generate for\\i = 0 .. 7\\标号 gen\_regs};
  \node[hw label, align=center, anchor=north] at (0,-0.95) {编译期复制：\\综合开始前展开};
  \draw[hw bus] (gen.east) -- (3.1,0.6);
  \node[hw label, above] at (2.4,0.65) {展开 8 份};
  \node[hw state, minimum width=4.6cm, minimum height=0.72cm] (r0) at (5.7,2.55) {gen\_regs[0]：regs[0] 触发器组};
  \node[hw state, minimum width=4.6cm, minimum height=0.72cm] (r1) at (5.7,1.7) {gen\_regs[1]：regs[1] 触发器组};
  \node[hw state, minimum width=4.6cm, minimum height=0.72cm] (r2) at (5.7,0.85) {gen\_regs[2]：regs[2] 触发器组};
  \node at (5.7,0.2) {$\vdots$};
  \node[hw state, minimum width=4.6cm, minimum height=0.72cm] (r7) at (5.7,-0.6) {gen\_regs[7]：regs[7] 触发器组};
  \node[hw label, anchor=north, align=center] at (5.7,-1.2) {每份独立使能 \code{we[i]}，共享 \code{clk} 与 \code{wdata}\\与手工抄写 8 遍一字不差};
\end{pdfdiagram}
\diagramnote{\code{generate for} 在综合开始前把寄存器组复制成 8 份独立硬件：每份一组触发器、各自一个使能端 \code{we[i]}，时钟与写数据共享。标号 \code{gen\_regs} 给每份命名，\code{gen\_regs[3]} 这样的层次路径会出现在波形与综合报告里，用来定位具体某一份。电路上没有任何“循环”在运行——它只是原理图复制粘贴的文字形式。}

流水线与脉动阵列这类“N 级相同结构串联”的形态同理，只是级间连线改用数组下标表达——第 $i$ 级的输出接第 $i+1$ 级的输入：

\begin{lstlisting}
genvar k;
generate
    for (k = 0; k < STAGES; k = k + 1) begin : gen_stage
        always_ff @(posedge clk) begin
            if (k == 0)                // k 是编译期常量
                pipe[k] <= sample_in;
            else
                pipe[k] <= pipe[k-1];  // 第 k 级接第 k-1 级
        end
    end
endgenerate
\end{lstlisting}

展开之后是 \code{STAGES} 级触发器排首尾相接，与“时序逻辑的 Verilog 写法”一节的移位寄存器同构，区别只在每级的内容从单 bit 换成一个 8 位样本。generate 不创造新硬件，它只是把重复了一遍又一遍的硬件收拢成一段参数化文字。同族的 \code{generate if} 按参数裁剪结构，本专题末尾会给出完整例子。

边界同样由“编译期”三个字决定：generate 的循环界与条件必须是编译期常量——\code{parameter}、\code{localparam}、\code{genvar}——不能是信号。硬件的份数在上板之前就定死了，运行时不能“多生成一份”。运行时要做的是“选择”而不是“生成”：从 N 份里选一份用 mux，那是“组合逻辑的 Verilog 写法”一节组合逻辑的工作，不是 generate 的工作。把这两件事分清，读参数化模块就不会把 \code{for} 想成程序循环——它只是原理图复制粘贴的文字形式。

参数化的最后一环在实例化一侧。“从电路图到 RTL：两种描述同一个硬件”一节说过实例化是放元件；带参数的元件在放置时给出本次的规格——宽度、深度、级数——同一模块名在不同位置可以放出不同规模的硬件。读层次化设计时，模块定义处的 \code{parameter} 默认值只是“不传参数时的规格”，真实尺寸以实例化处为准；追踪一个信号的位宽，要一路追到最上层的参数传递，而不是停在模块声明那一行。

把按参数裁剪写成代码，就是 \code{generate if}：

\begin{lstlisting}
generate
    if (WIDTH > 1) begin : gen_wide
        // 多位版本：进位逐级串联成链
        assign c_out = c[WIDTH];
    end else begin : gen_bit
        // 单位版本：没有进位链，直接相连
        assign c_out = c_in;
    end
endgenerate
\end{lstlisting}

条件里只能出现编译期常量，于是裁剪在综合开始前一锤定音：被选中的分支展开成真实电路，被裁掉的分支连一个门都不会剩下——留下的电路里找不到任何“判断”的痕迹。读带 \code{generate if} 的模块，顺序是先确定本次实例的参数值，再决定哪一份电路真正存在；另一份只是文字，不是硬件。这与运行时选择形成最后一组对照：\code{if} 语句选的是数据走向，\code{generate if} 选的是电路有无，一个发生在每一个时钟周期，一个发生在按下综合按钮的那一刻。

\section{五、testbench 与仿真阅读}

仿真器里没有电压与电流，只有一个按时间推进的事件队列；被测的 RTL 在队列里被逐事件推演，而推动这一过程的环境叫 testbench。testbench 本身不是硬件，也永远不会变成硬件：它是一台“用文字搭出来的测试台”，负责产生时钟与复位、施加激励、观察输出并报告结果。分清“哪些代码要变成电路、哪些代码只在仿真器里工作”，是读懂一切工程仿真的前提。

\topic{testbench 是激励环境，不是被综合的硬件}先看机制。仿真器启动时需要一个顶层模块，这个模块没有端口——它对外的接口就是仿真器本身。testbench 正是这个无端口 module：内部用 \code{logic} 与 \code{wire} 声明一组信号，实例化被测模块（工程上常称 DUT），把这些信号接到 DUT 的端口上；然后用 \code{initial} 与 \code{always} 块按时间驱动 DUT 的输入，用系统任务观察 DUT 的输出：\codetext{\$display} 在调用时刻打印一行，\codetext{\$monitor} 挂上之后每当所列信号变化就自动打印，\codetext{\$fatal} 打印错误并以失败姿态终止仿真（本节“激励与自检的写法”一段细讲）。整个结构可以读成面包板上那套动作的翻版：实例化 DUT 是把芯片插上面包板，驱动输入是拨动开关，观察输出是看 LED。

写法上，一个最小 testbench 骨架只有五个零件：无端口 module、信号声明、DUT 实例化、时钟激励、测试流程。

\begin{lstlisting}
module tb_xxx;                    // 无端口：仿真器从这里启动
  logic        clk, rst_n;        // 激励：将要接到 DUT 输入端
  wire  [7:0]  y;                 // 观察：接 DUT 输出端

  xxx dut (                       // DUT：被测模块，可综合的 RTL
    .clk(clk), .rst_n(rst_n), .y(y)
  );

  always #5 clk = ~clk;           // 时钟激励，下一小节细讲

  initial begin                   // 测试流程：复位、激励、自检
    rst_n = 1'b0;
    #20  rst_n = 1'b1;
    // ... 施加激励、与预期比对 ...
    $finish;
  end
endmodule
\end{lstlisting}

边界只有一条判据：这段代码要不要变成芯片上的电路。DUT 以及被它实例化的一切，必须遵守前四节的全部综合纪律，每一行都要能读回触发器、组合门与连线；testbench 里的一切则不受约束——\code{initial}、\code{\#} 延时、\codetext{\$display}、文件读写全部合法，因为它们只活在仿真器里，综合器永远看不到。工程上 testbench 还有一项特权：层次化引用，例如 \code{dut.u\_ram.mem[0]} 可以直接伸进 DUT 内部读写某个存储单元。这种写法不产生任何硬件，只是仿真器提供的观察与干预通道，本节后面的案例会用它给整机预置程序。

把这台虚拟实验台画成挂接图，四个部件的分工就是 testbench 的全部结构：

\begin{pdfdiagram}
  % testbench 边界：全部部件不综合
  \draw[draw=BookNavy, dashed, rounded corners=3pt, line width=0.5pt] (-0.9,-2.6) rectangle (10.3,2.2);
  \node[hw label] at (6.55,1.92) {tb\_xxx：仿真器内部，不综合};
  \node[hw control, minimum width=2.5cm, minimum height=1.3cm] (stim) at (0.9,0.3) {激励\\时钟/复位/输入序列};
  \node[hw compute, minimum width=2.0cm, minimum height=1.3cm] (dut) at (4.6,0.3) {DUT\\被测 RTL};
  \node[hw state, minimum width=2.5cm, minimum height=1.3cm] (chk) at (8.3,0.3) {自检\\\code{!==} 比对预期};
  \node[hw control, minimum width=2.6cm, minimum height=0.9cm] (wd) at (4.6,-1.7) {看门狗：超时 \codetext{\$fatal}};
  \draw[hw bus] (stim.east) -- (dut.west);
  \node[hw label, above] at (2.85,0.32) {驱动输入};
  \draw[hw bus] (dut.east) -- (chk.west);
  \node[hw label, above] at (7.05,0.32) {观察输出};
  \draw[hw return, dashed] (wd.north) -- (dut.south);
  \node[hw label, anchor=west] at (4.75,-0.85) {监视停滞};
  \node[hw label, anchor=north] at (8.3,-0.55) {\codetext{\$display} / \codetext{\$fatal} 报告结果};
\end{pdfdiagram}
\diagramnote{testbench 的四个部件：激励发生器产生时钟与复位、按时间驱动输入序列；DUT 是边界内唯一要综合的部分，遵守前四节全部纪律；自检把输出与预期逐项比对（\code{!==} 让 X 也判失败）；看门狗在程序陷入死循环时强制作废本次仿真。四个部件连同虚线边界，只存在于仿真器里。}

\topic{时钟与复位的产生}机制：仿真时间是离散的。骨架里那一行 \code{always \#5} 让 \code{clk} 每过 5 个时间单位取反一次：0 时刻为 0，第 5 单位变 1，第 10 单位变 0，再第 15 单位变 1——周期因而是 10 个时间单位，这是全仿真界通行的写法：\code{\#} 后面的数字是半周期。文件开头的编译指令 \code{`timescale 1ns/1ns} 把“单位”定为纳秒，这条时钟就是周期 10 ns、频率 100 MHz 的方波。这个 always 块在硬件里的对应物是本书末章的教学计算机项目的 555 振荡器：一个不依赖别人、自己永远走下去的激励源。

复位的写法同样对应面包板上的复位按钮：上电即按下，按住若干拍，再松开。

\begin{lstlisting}
  initial begin
    rst_n = 1'b0;                 // 上电即复位
    repeat (5) @(negedge clk);    // 按住 5 拍
    rst_n = 1'b1;                 // 在时钟低电平期间释放
  end
\end{lstlisting}

拉低若干拍再释放，是为了让 DUT 里每个触发器都至少经历一次复位、进入已知初态——这正是本书末章的教学计算机项目“复位覆盖全部时序状态”那一课的仿真版：不复位，寄存器初值全是 X，波形从第一拍起就不可读。选在低电平期间释放是个小习惯：让释放动作远离时钟沿，避免与 DUT 的时序块争抢同一时刻。

边界：\code{\#} 延时与 \code{repeat} 等待只在 testbench 里合法。RTL 中不存在“等 5 ns”的电路——“从 Verilog 读回硬件：综合边界”一节讲不可综合写法时已经说过，综合器面对延时语句要么报错要么悄悄忽略，两者都危险。真实硬件的时钟来自晶振或 555，复位来自 RC 电路或按钮；仿真里的这两三行是它们的替身，而不是对它们的描述。

把 \code{`timescale}、时钟、复位释放与自检检查点画到同一条时间轴上，这几行代码的时间含义就全部可见：

\begin{waveform}[x=0.9cm,y=0.85cm]
  % 顶部说明
  \node[hw label] at (4.3,7.55) {\code{`timescale 1ns/1ns}：1 单位 = 1 ns；\code{always \#5} 给半周期 5 ns、周期 10 ns};
  % 参考虚线：复位释放（50 ns）与两个检查点（60、70 ns）
  \draw[BookRule, line width=0.45pt, dashed] (5.5,7.1) -- (5.5,3.3);
  \draw[BookRule, line width=0.45pt, dashed] (6.6,4.2) -- (6.6,3.3);
  \draw[BookRule, line width=0.45pt, dashed] (7.7,4.2) -- (7.7,3.3);
  \node[hw label, anchor=west] at (5.62,5.85) {复位释放（50 ns）};
  \node[hw label, anchor=west] at (5.62,4.6) {检查点};
  % clk：沿在 0.55 的倍数处（5 ns 一格）
  \draw[BookAccent, line width=0.9pt] (0,6.3) -- (0.55,6.3) -- (0.55,6.9) -- (1.1,6.9) -- (1.1,6.3) -- (1.65,6.3) -- (1.65,6.9) -- (2.2,6.9) -- (2.2,6.3) -- (2.75,6.3) -- (2.75,6.9) -- (3.3,6.9) -- (3.3,6.3) -- (3.85,6.3) -- (3.85,6.9) -- (4.4,6.9) -- (4.4,6.3) -- (4.95,6.3) -- (4.95,6.9) -- (5.5,6.9) -- (5.5,6.3) -- (6.05,6.3) -- (6.05,6.9) -- (6.6,6.9) -- (6.6,6.3) -- (7.15,6.3) -- (7.15,6.9) -- (7.7,6.9) -- (7.7,6.3) -- (8.5,6.3);
  \node[hw label, anchor=east] at (-0.2,6.6) {clk};
  % rst_n：按住 5 拍，第 5 个 negedge（50 ns）释放
  \draw[BookAccent, line width=0.9pt] (0,5.0) -- (5.5,5.0) -- (5.5,5.6) -- (8.5,5.6);
  \node[hw label, anchor=east] at (-0.2,5.3) {rst\_n};
  \node[hw label] at (2.75,4.68) {repeat (5) @(negedge clk)：按住 5 拍};
  % 自检行：每个 negedge 触发比对
  \draw[BookMuted, line width=0.6pt] (0,3.9) -- (8.5,3.9);
  \foreach \x in {1.1, 2.2, 3.3, 4.4, 5.5, 6.6, 7.7} {
    \draw[BookMuted, line width=0.6pt] (\x,3.75) -- (\x,4.05);
  }
  \node[hw label, anchor=east] at (-0.2,3.9) {自检触发};
  \node[hw label, anchor=east] at (5.35,3.45) {每个 negedge 比对 \code{out} 与预期，避开沿上竞争};
  % 时间轴
  \draw[BookMuted, line width=0.6pt, ->] (0,2.9) -- (8.8,2.9);
  \foreach \x/\t in {0/0, 1.1/10, 2.2/20, 3.3/30, 4.4/40, 5.5/50, 6.6/60, 7.7/70} {
    \draw[BookMuted, line width=0.5pt] (\x,2.96) -- (\x,2.84);
    \node[hw label, anchor=north] at (\x,2.8) {\t};
  }
  \node[hw label, anchor=west] at (8.9,2.9) {ns};
\end{waveform}
\diagramnote{仿真时间轴与三段代码的对应：\code{`timescale 1ns/1ns} 把 1 个时间单位定为 1 ns，\code{always \#5} 产生周期 10 ns 的时钟；复位拉低后经 \code{repeat (5) @(negedge clk)} 按住 5 拍，在第 5 个下降沿（50 ns）释放——刻意选在低电平期间，远离采样沿；自检块在每个下降沿比对 \code{out} 与预期，同样避开沿上的读写竞争。RTL 仿真里时间是离散的，这条轴上的每一格都是事件队列里的一个时刻。}

\topic{激励与自检的写法}机制：激励是按时间施加到 DUT 输入端的值序列；自检是让 testbench 自己把实际输出与预期值逐项比对，只在最后报告结论。人眼读波形只能算抽查——波形越长，漏看越多；能自动比对的测试才谈得上可重复。三个系统任务的分工前面已经交代：\codetext{\$display} 报进展，\codetext{\$monitor} 追变化，\codetext{\$fatal} 报错退出，脚本凭退出状态判失败。

写法：把激励与预期都预置成数组，逐项驱动、逐项比对。以测一个 8 位加法器（端口 \code{a}、\code{b}、\code{y}）为例：

\begin{lstlisting}
  logic [15:0] stim [0:3];        // {a, b} 激励对
  logic [7:0]  gold [0:3];        // 预期和
  integer i, errs;
  initial begin
    stim[0] = {8'd3,   8'd4};   gold[0] = 8'd7;
    stim[1] = {8'd9,   8'd9};   gold[1] = 8'd18;
    stim[2] = {8'd200, 8'd100}; gold[2] = 8'd44;  // 溢出回绕
    stim[3] = {8'd0,   8'd0};   gold[3] = 8'd0;
    errs = 0;
    for (i = 0; i < 4; i = i + 1) begin
      {a, b} = stim[i];  #10;   // 驱动后等组合路径稳定
      if (y !== gold[i]) begin
        $display("FAIL: i=%0d y=%0d expect=%0d", i, y, gold[i]);
        errs = errs + 1;
      end
    end
    if (errs == 0) $display("ALL PASS");
    else           $fatal(1, "%0d errors", errs);
    $finish;
  end
\end{lstlisting}

注意比对用 \code{!==}（区分 X 的严格不等）：输出哪怕是 X，也会被正确地判成失败。若用普通的 \code{!=}，含 X 的比较结果是“不确定”，错误就从指缝里漏过去了。

规模再大一些的工程用 golden 参考：同一模块写两份，一份是行为级的“会说话的规格”——怎么好懂怎么写，一份是“从电路图到 RTL：两种描述同一个硬件”一节到“从 Verilog 读回硬件：综合边界”一节风格的 RTL——怎么可综合怎么写；同一份激励同时输入两者，逐拍比对输出。此时预期值不再由人手工计算，而是由行为模型现场给出，RTL 每改一版就自动重测一轮。边界要记牢：自检的质量不会超过预期值的质量——预期算错，全部 PASS 也是假象；预期表的第一读者永远是人。

\topic{波形怎么读}波形窗口是三样东西的并排：纵轴是信号名，横轴是仿真时间，中间是每个信号的取值轨道。单比特信号画高低线，多比特总线画成首尾尖起的块、块内标注当前值。读波形先认四种特殊显示：

\begin{longtable}{L{0.24\linewidth}L{0.3\linewidth}L{0.34\linewidth}}
\toprule
显示 & 含义 & 常见根源 \\
\midrule
X（通常红色） & 未知态：仿真器不知道该值是 0 还是 1 & 寄存器没复位、存储器没初始化、多驱动冲突 \\\tablerowrule
Z（通常蓝色细线） & 高阻：此刻没有任何驱动者 & 三态门全部关闭，如本书末章的教学计算机项目空拍时的总线 \\\tablerowrule
总线值突然变 X & 两个驱动者给出相反电平 & 总线争用：同一拍开了两个“出”使能 \\\tablerowrule
时钟沿上的瞬时细线 & 零延迟模型中的事件竞争 & 同一时刻多个块读写同一信号 \\
\bottomrule
\end{longtable}

最有用的排障手法是顺着 X 往回走：找到第一个变 X 的信号，看它的驱动源；驱动源若也是 X，就再往上追一层。链条的尽头几乎总是同一个地方——一个没接复位的触发器，或一块没写初值的存储器。这与本书末章的教学计算机项目上电检查的第一条互为镜像：真机上“复位后全部寄存器应为已知值”，仿真里“复位释放后波形不应再出现 X”。Z 的读法则反过来：看到 Z 先别慌，问“此刻谁该驱动”。答案若本来就是“没人”——例如本书末章的教学计算机项目空拍时的总线——Z 是正常；答案若该有人驱动，那就是使能链断在了某一级。

把“顺着 X 往回走”画成波形：q1 没接复位，它的 X 在一个沿之后传染给 q2，再顺着组合路径染到输出 y：

\begin{waveform}[x=0.9cm,y=0.85cm]
  % 参考虚线：沿 1（复位释放）与沿 2（X 被采进 q2）
  \draw[BookRule, line width=0.45pt, dashed] (2,7.3) -- (2,1.5);
  \draw[BookRule, line width=0.45pt, dashed] (4,7.3) -- (4,1.5);
  \node[hw label] at (4.55,7.5) {沿 1：复位释放；沿 2：q2 采进 q1 的 X};
  % clk
  \draw[BookAccent, line width=0.9pt] (0.5,6.4) -- (2,6.4) -- (2,7.0) -- (3,7.0) -- (3,6.4) -- (4,6.4) -- (4,7.0) -- (5,7.0) -- (5,6.4) -- (6,6.4) -- (6,7.0) -- (7,7.0) -- (7,6.4) -- (8.6,6.4);
  \node[hw label, anchor=east] at (0.3,6.7) {clk};
  % rst_n：低 5 拍后释放（沿 1）
  \draw[BookAccent, line width=0.9pt] (0.5,5.3) -- (2,5.3) -- (2,5.9) -- (8.6,5.9);
  \node[hw label, anchor=east] at (0.3,5.6) {rst\_n};
  % q1：从未复位，恒为 X
  \draw[BookRed, line width=0.8pt, fill=BookRed!12] (0.5,4.2) rectangle (8.6,5.0);
  \node at (4.55,4.6) {X：从未复位};
  \node[hw label, anchor=east] at (0.3,4.6) {q1};
  % q2：有复位，沿 2 采进 q1 的 X
  \draw[BookTeal, line width=0.9pt, fill=BookCodeBg] (0.5,2.7) rectangle (4,3.5);
  \node at (2.25,3.1) {0};
  \draw[BookRed, line width=0.8pt, fill=BookRed!12] (4,2.7) rectangle (8.6,3.5);
  \node at (6.3,3.1) {X};
  \node[hw label, anchor=east] at (0.3,3.1) {q2};
  % y = q2 + 1（组合）：q2 变 X 后立即跟着变 X
  \draw[BookTeal, line width=0.9pt, fill=BookCodeBg] (0.5,1.6) rectangle (4,2.4);
  \node at (2.25,2.0) {1};
  \draw[BookRed, line width=0.8pt, fill=BookRed!12] (4,1.6) rectangle (8.6,2.4);
  \node at (6.3,2.0) {X};
  \node[hw label, anchor=east] at (0.3,2.0) {y = q2+1};
  % 排障顺序
  \node[hw label, anchor=north] at (4.55,1.3) {排障反向走：y 的 X $\leftarrow$ q2 的 X $\leftarrow$ q1（源头：未复位）};
\end{waveform}
\diagramnote{X 沿数据路径的传播链：q1 从未复位，恒为 X；沿 2 这个 X 被采进 q2，组合输出 y 随即跟着变 X。排障沿链反向走——先看到 y 变 X，追到 q2，再追到源头 q1：一只没接复位的触发器。rst\_n 对 q2 有效、对 q1 缺席，正是本图假设的接线；给 q1 补上复位，整条链在第一拍之后就不再有 X。}

\topic{案例：把本书末章的教学计算机项目教学计算机写进仿真}现在把全部零件合起来。目标：用“从电路图到 RTL：两种描述同一个硬件”一节到“从 Verilog 读回硬件：综合边界”一节的写法，把本书末章的教学计算机项目那台 8-bit 总线机器写成 RTL；再用本节的写法给它配 testbench——把本书末章的教学计算机项目程序一（从 1 加到 4）预置进 RAM，让仿真自己跑到 HLT，并自动核对 OUT 依次显示 1、3、6、10。

机器的核心是控制单元，而本书末章的教学计算机项目“从 Verilog 读回硬件：综合边界”一节已经给出了它的全部内容：一张以 \{opcode, T\} 为地址的微码表。RTL 化几乎是逐行誊写：每个 T 状态一个分支，每根有效控制线一个赋值。数据通路沿用本书末章的教学计算机项目的总线结构——任一时刻至多一个驱动者，空拍时总线高阻。先写存储器：

\begin{lstlisting}
module sap1_ram (                 // 16 字节：程序与数据同住
    input  logic       clk,
    input  logic       we,        // 即控制线 RI
    input  logic [3:0] addr,      // 即 MAR 的输出
    input  logic [7:0] din,
    output logic [7:0] dout
);
  logic [7:0] mem [0:15];

  always_ff @(posedge clk) begin  // 写：时钟沿、RI 有效时
    if (we) mem[addr] <= din;
  end

  always_comb begin               // 读：组合读出，相当于常开的 RO 三态门
    dout = mem[addr];
  end
endmodule
\end{lstlisting}

注意 RAM 没有复位端：复位清的是机器状态，不是程序——面包板上拨进去的字节，也不会被复位按钮抹掉。整机如下，控制线与本书末章的教学计算机项目“从 Verilog 读回硬件：综合边界”一节一一对应，只有三处形式差异：J 在这里叫 \code{pi}，$\Sigma$O 写作 \code{so}，SU 不再单列为一根控制线，改由 \code{op} 译码在 ALU 内部取代：

\begin{lstlisting}
module sap1_machine (
    input  logic       clk,
    input  logic       rst_n,
    output logic [7:0] out,       // OUT 寄存器，接显示排
    output logic       halted     // HLT 后时钟冻结
);
  logic [7:0] a, b, ir, bus;
  logic [3:0] pc, mar;
  logic [2:0] t;                  // T 状态计数器：T0..T7 循环
  logic       zf, cf;
  logic [7:0] ram_dout;
  logic [3:0] op;
  logic [8:0] alu9;               // 第 9 位就是进位

  // “出”使能（谁驱动总线）与“入”使能（谁在这一沿装载）
  logic co, ro, ao, io, so;
  logic mi, ii, ce, ai, bi, oi, ri, pi, fi, hlt;

  assign op = ir[7:4];

  // 总线：任一时刻至多一个驱动者；无人驱动时为高阻（本书末章的教学计算机项目的空拍）
  assign bus = co ? {4'h0, pc}      :
               ro ? ram_dout        :
               ao ? a               :
               io ? {4'h0, ir[3:0]} :
               so ? alu9[7:0]       : 8'hzz;

  // ALU：减法是取反加一（算术与数据通路相关章节）
  always_comb begin
    if (op == 4'h3) alu9 = {1'b0, a} + {1'b0, ~b} + 9'd1;
    else            alu9 = {1'b0, a} + {1'b0, b};
  end

  // 微码译码：本书末章的教学计算机项目微码 ROM 的 RTL 抄本，行地址是 {op, t}
  always_comb begin
    co = 0; ro = 0; ao = 0; io = 0; so = 0;
    mi = 0; ii = 0; ce = 0; ai = 0; bi = 0;
    oi = 0; ri = 0; pi = 0; fi = 0; hlt = 0;
    case (t)
      3'd0: begin co = 1; mi = 1; end        // T0: MAR <- PC
      3'd1: begin ro = 1; ii = 1; ce = 1; end// T1: IR <- RAM[MAR], PC <- PC+1
      3'd2: case (op)                        // T2: 地址场 / 立即数 / 跳转 / 输出 / 停机
        4'h1, 4'h2, 4'h3, 4'h4: begin io = 1; mi = 1; end
        4'h5: begin io = 1; ai = 1; end      // LDI: A <- 立即数
        4'h6: begin io = 1; pi = 1; end      // JMP: PC <- 地址场
        4'h7: begin io = 1; pi = cf; end     // JC:  cf=1 才跳
        4'h8: begin io = 1; pi = zf; end     // JZ:  zf=1 才跳
        4'hE: begin ao = 1; oi = 1; end      // OUT: OUT <- A
        4'hF: hlt = 1;                       // HLT: 冻结
        default: ;
      endcase
      3'd3: case (op)                        // T3: 访存与装载
        4'h1:       begin ro = 1; ai = 1; end// LDA: A <- RAM[MAR]
        4'h2, 4'h3: begin ro = 1; bi = 1; end// ADD/SUB: B <- RAM[MAR]
        4'h4:       begin ao = 1; ri = 1; end// STA: RAM[MAR] <- A
        default: ;
      endcase
      3'd4: case (op)                        // T4: 写回 A 与标志
        4'h2, 4'h3: begin so = 1; ai = 1; fi = 1; end
        default: ;
      endcase
      default: ;                             // T5..T7: 空拍，总线高阻
    endcase
  end

  // 全部寄存器共用一个时钟沿；“入”使能决定谁在这一沿装载
  always_ff @(posedge clk or negedge rst_n) begin
    if (!rst_n) begin
      pc <= 0; mar <= 0; ir <= 0; a <= 0; b <= 0;
      out <= 0; t <= 0; zf <= 0; cf <= 0; halted <= 0;
    end else if (!halted) begin
      t <= t + 3'd1;
      if (mi)  mar <= bus[3:0];
      if (ii)  ir  <= bus;
      if (ce)  pc  <= pc + 4'd1;
      if (pi)  pc  <= bus[3:0];
      if (ai)  a   <= bus;
      if (bi)  b   <= bus;
      if (oi)  out <= bus;
      if (fi)  begin zf <= (alu9[7:0] == 8'h00); cf <= alu9[8]; end
      if (hlt) halted <= 1'b1;
    end
  end

  sap1_ram u_ram (
    .clk(clk), .we(ri), .addr(mar), .din(bus), .dout(ram_dout)
  );
endmodule
\end{lstlisting}

对照本书末章的教学计算机项目的逐拍表读这段代码：T0 的 \code{co+mi} 就是 \code{CO+MI}，T1 的 \code{ro+ii+ce} 就是 \code{RO+II+CE}，ADD 的 T3、T4 正是 \code{RO+BI} 与 $\Sigma$O+\code{AI+FI}。微码 ROM 在面包板上是一片 EEPROM，在这里是一个组合逻辑块——同一张真值表的两种硬件实现，行为完全一致。\code{ce} 与 \code{pi} 永远不同时有效，所以 PC 的两个装载条件互斥；这份互斥不是语法保证的，是微码表保证的，与本书末章的教学计算机项目“出使能互斥”完全同构。

testbench 负责剩下的一切：时钟、复位、预置程序（对应编程模式拨码）、自检，外加一只看门狗。

\begin{lstlisting}
`timescale 1ns/1ns
module tb_sap1;
  logic       clk   = 1'b0;
  logic       rst_n = 1'b0;
  wire  [7:0] out;
  wire        halted;

  sap1_machine dut (              // 被测整机
    .clk(clk), .rst_n(rst_n), .out(out), .halted(halted)
  );

  always #5 clk = ~clk;           // 时钟：周期 10 ns

  initial begin                   // 复位：拉低 5 拍后释放
    repeat (5) @(negedge clk);
    rst_n = 1'b1;
  end

  initial begin                   // 预置程序：本书末章的教学计算机项目程序一逐字节抄录
    dut.u_ram.mem[ 0] = 8'h1F;    // LDA 15    A <- sum
    dut.u_ram.mem[ 1] = 8'h2E;    // ADD 14    A <- sum + i
    dut.u_ram.mem[ 2] = 8'h4F;    // STA 15    sum 写回
    dut.u_ram.mem[ 3] = 8'hE0;    // OUT       显示部分和
    dut.u_ram.mem[ 4] = 8'h1E;    // LDA 14    A <- i
    dut.u_ram.mem[ 5] = 8'h3D;    // SUB 13    i - 4，判顶
    dut.u_ram.mem[ 6] = 8'h8B;    // JZ 11     到顶即停机
    dut.u_ram.mem[ 7] = 8'h1E;    // LDA 14
    dut.u_ram.mem[ 8] = 8'h2C;    // ADD 12    i + 1
    dut.u_ram.mem[ 9] = 8'h4E;    // STA 14    i 写回
    dut.u_ram.mem[10] = 8'h60;    // JMP 0     回循环头
    dut.u_ram.mem[11] = 8'hF0;    // HLT
    dut.u_ram.mem[12] = 8'h01;    // 常数 1
    dut.u_ram.mem[13] = 8'h04;    // 上限 4
    dut.u_ram.mem[14] = 8'h01;    // i   初值 1
    dut.u_ram.mem[15] = 8'h00;    // sum 初值 0
  end

  // 自检：OUT 每装入一个新值，就与预期序列对一项
  logic [7:0] expect [0:3];
  logic [7:0] out_d;
  integer     seen;
  initial begin
    expect[0] = 8'd1; expect[1] = 8'd3;
    expect[2] = 8'd6; expect[3] = 8'd10;
    out_d = 8'h00;                // 复位后 OUT 的已知初值
    seen  = 0;
  end

  always @(negedge clk) begin     // 在时钟低电平期间检查，避开沿上竞争
    if (rst_n && (out !== out_d)) begin
      out_d = out;
      if (seen > 3)
        $fatal(1, "FAIL: extra OUT write %0d", out);
      else if (out !== expect[seen])
        $fatal(1, "FAIL: step %0d expect %0d got %0d", seen, expect[seen], out);
      else
        $display("[%0t ns] OUT = %0d  PASS step %0d", $time, out, seen);
      seen = seen + 1;
    end
  end

  initial begin                   // 收尾：HLT 后核对输出次数
    wait (halted);
    repeat (2) @(negedge clk);
    if (seen == 4) $display("TEST PASS: OUT sequence 1 3 6 10, %0d checks", seen);
    else           $fatal(1, "FAIL: halted after %0d outputs", seen);
    $finish;
  end

  initial begin                   // 看门狗：程序死循环时强制作废本次仿真
    #20000;
    $fatal(1, "FAIL: timeout, machine never halted");
  end
endmodule
\end{lstlisting}

这份 testbench 把前面几个写法一次用全：时钟与复位的产生、层次化预置（\code{dut.u\_ram.mem} 就是编程模式的拨码）、数组化的预期值、\codetext{\$fatal} 与 \codetext{\$display} 自检；看门狗保证微码若有错、程序陷入死循环时，仿真不会无限空转。在 1 ns 时标、复位拉低 5 拍的这份激励下，仿真输出为：

\begin{lstlisting}
[320 ns] OUT = 1  PASS step 0
[1200 ns] OUT = 3  PASS step 1
[2080 ns] OUT = 6  PASS step 2
[2960 ns] OUT = 10  PASS step 3
TEST PASS: OUT sequence 1 3 6 10, 4 checks
\end{lstlisting}

四次 OUT 与本书末章的教学计算机项目的逐圈表逐项吻合；机器执行 41 条指令后在 HLT 的 T2 冻结，从复位释放到冻结共 323 拍。面包板上要按三百余次单步才能看完的过程，仿真器里只是一瞬间——但两者跑的是同一张微码表、同一份程序、同一个输出序列。仿真不是另一台机器，是同一台机器的理想化影子。

这份 testbench 里有两条层次化通道值得单独画出：预置程序写进 \code{dut.u\_ram.mem} 的路径，和从 \code{out} 到自检的观测路径。它们都不产生任何硬件。

\begin{pdfdiagram}[x=0.85cm]
  % testbench 边界：全部部件不综合
  \draw[draw=BookNavy, dashed, rounded corners=3pt, line width=0.5pt] (-1.0,-3.3) rectangle (10.4,2.9);
  \node[hw label] at (6.8,2.62) {tb\_sap1：仿真器内部，全部不综合};
  \node[hw control, minimum width=1.9cm, minimum height=1.2cm, align=center] (stim) at (0.35,1.3) {激励\\clk / rst\_n};
  % DUT 边界：唯一可综合的部分
  \draw[draw=BookMuted, dashed, rounded corners=3pt, line width=0.5pt] (1.9,-0.9) rectangle (5.5,2.3);
  \node[hw label] at (3.7,2.03) {dut：sap1\_machine（可综合 RTL）};
  \node[hw label] at (3.7,1.35) {控制单元与数据通路（从略）};
  \node[hw memory, minimum width=2.7cm, minimum height=1.1cm, align=center] (uram) at (3.9,0.1) {u\_ram：sap1\_ram\\mem[0:15]};
  \draw[hw bus] (stim.east) -- (1.9,1.3);
  % out 观测路径（实线：真实端口）
  \node[hw state, minimum width=2.3cm, minimum height=1.2cm, align=center] (chk) at (8.4,1.3) {自检\\expect[0:3]};
  \draw[hw bus] (5.5,1.3) -- (chk.west);
  \node[hw label, above] at (6.7,1.32) {out};
  \node[hw label, anchor=east, align=center] at (8.2,0.1) {每个 negedge 比对\\out 与 expect[seen]};
  % 预置路径（虚线：层次化写）
  \node[hw control, minimum width=2.9cm, minimum height=0.95cm, align=center] (pre) at (0.35,-1.9) {initial：预置程序};
  \draw[hw return, dashed] (pre.east) -- (3.9,-1.9) -- (uram.south);
  \node[hw label, anchor=north] at (1.8,-2.1) {\code{dut.u\_ram.mem[0..15]} $\leftarrow$ 程序字节};
  % 层次化观测（虚线：直读内部信号）
  \draw[hw return, dashed] (chk.south) -- (8.4,-2.9) -- (5.0,-2.9) -- (5.0,-0.9);
  \node[hw label, anchor=east] at (4.85,-2.62) {探针也可直读 \code{dut.pc}、\code{dut.t}（测点）};
\end{pdfdiagram}
\diagramnote{两条层次化通道。写：预置 initial 用 \code{dut.u\_ram.mem[...]} 穿过两层实例边界，直接把程序字节写进 RAM——对应面包板编程模式的拨码。读：自检沿真实端口 \code{out} 逐项比对预期序列；需要更深的位置时，层次路径还能直读 \code{dut.pc}、\code{dut.t} 作波形测点。两条通道只存在于仿真器里，综合器永远看不到。}

\topic{仿真与真实硬件的差别}仿真通过之后，最容易产生的错觉是“硬件也就这样了”。三处差别必须记牢。第一是零延迟的理想时序：RTL 仿真里组合逻辑在零时间内完成，所有触发器在同一个时钟沿同时更新；“电源、时钟与信号完整性”一章讲的传播延迟、建立与保持时间、时钟偏斜，在这幅画面里根本不存在——它们要换成标注了真实延时的门级仿真才看得见，或者上板以后由示波器告诉你。第二是电气问题不可见：仿真世界里信号只有 0、1、X、Z 四种取值，没有电压与电流；电源跌落、串扰、短路发热、毛刺携带的能量，统统无法在逻辑仿真里现身。第三是 X 的语义边界：X 的意思是“仿真器不知道”，而不是某种物理电压。真实芯片上电后，每一位总是确定的 0 或 1（或正在翻转的亚稳态），只是你无法预知；仿真却可能因若干巧合，让这份“不知道”从未传播到输出——错误被 X 的乐观藏了起来，仿真通过而硬件照样出错。

\begin{longtable}{L{0.3\linewidth}L{0.32\linewidth}L{0.22\linewidth}}
\toprule
仿真里的世界 & 真实硬件 & 补课章节 \\
\midrule
组合逻辑零延迟，沿上同时更新 & 传播延迟、建立与保持时间、时钟偏斜 & “电源、时钟与信号完整性”一章 \\\tablerowrule
信号只有 0/1/X/Z & 电压、电流、串扰、发热 & “数字硬件基础”部分、“电源、时钟与信号完整性”一章 \\\tablerowrule
X 表示“不知道”，可以不传播 & 每位总是确定的 0 或 1（或亚稳态） & “状态、时钟与时序逻辑”一章、“电源、时钟与信号完整性”一章 \\\tablerowrule
时钟与复位几行代码就有 & 晶振、复位电路要单独设计 & 本书末章的教学计算机项目、“电源、时钟与信号完整性”一章 \\
\bottomrule
\end{longtable}

所以仿真通过的正确定位是：“可以上板”的必要条件，而不是充分条件。它保证的是逻辑意图被正确地翻译成了 RTL；时序是否收敛、电气是否健康，是另外两关，分别在时序分析与“电源、时钟与信号完整性”一章的测量里。把这三关的顺序搞反——先上板、出了怪事再回头怀疑逻辑——是调试时间最大的浪费。




























\section{六、从 RTL 到可实现约束}

RTL 仿真回答“在给定激励下，描述的状态怎样变化”，但芯片或 FPGA 还必须回答“这些变化能否在真实时钟和引脚边界内发生”。工具首先进行展开与静态检查：解析参数和 generate 分支，确定端口宽度与驱动关系，再报告未连接端口、位宽截断、组合环、多重驱动、未完全赋值以及不可达状态。lint 不是语法检查的别名；它要在仿真尚未表现异常时指出可能形成锁存器、X 扩散或硬件歧义的结构。

时序约束把外部契约写给静态时序分析器。主时钟约束给出周期与波形；派生时钟说明分频、倍频或门控后的相位关系；输入延迟描述外部发送器相对本芯片采样沿何时给出数据；输出延迟描述板级接收器要求本芯片何时稳定输出。没有输入/输出延迟的设计即使内部路径全部通过，也没有证明芯片边界满足接口时序。

\begin{longtable}{L{0.18\linewidth}L{0.34\linewidth}L{0.38\linewidth}}
\toprule
约束对象 & 它表达的硬件事实 & 缺失或误用的后果 \\
\midrule
主时钟 & 寄存器采样边沿和周期预算 & 工具不知道组合路径必须多快到达 \\\tablerowrule
派生时钟 & 两个相关时钟的频率、相位和来源 & 把相关路径误当异步，或使用错误采样沿 \\\tablerowrule
输入延迟 & 板外器件已经消耗的周期预算 & 内部输入路径通过，系统级建立时间仍可能失败 \\\tablerowrule
输出延迟 & 板外接收器要求保留的建立/保持预算 & 输出到达过晚或变化过早 \\\tablerowrule
时钟不确定度 & 抖动、偏斜和建模余量 & 纸面周期没有给真实分配网络留出裕量 \\\tablerowrule
例外路径 & 功能上不按普通单周期采样的路径 & 错误例外会隐藏真实失败，缺少例外会产生伪告警 \\
\bottomrule
\end{longtable}

false path 只能用于功能上永远不会被同步采样的路径，例如已经由异步协议隔离的跨域关系；multicycle path 表示接收端确实允许多个周期后采样，因此还要同步修正保持检查。两者都不是“让报告变绿”的选项。每条例外都应能回到电路模式、状态条件和接口协议，证明被排除的采样事件不会发生。

\section{七、静态时序分析把路径拆成预算}

静态时序分析（Static Timing Analysis, STA）不依赖某一组输入向量，而是沿时序图检查所有可达寄存器路径。建立时间检查比较数据最晚到达时刻与接收寄存器的要求时刻；保持时间检查比较数据最早到达时刻与旧值必须维持到的时刻。两者使用不同延迟角：建立失败通常来自数据路径过慢，保持失败通常来自数据路径过快或时钟偏斜不利。

设发射寄存器时钟到输出为 $t_{cq}$，组合路径最大延迟为 $t_{comb,max}$，接收寄存器建立时间为 $t_{setup}$，时钟偏斜与不确定度预算合计为 $t_u$，则单周期路径至少要求

\[
T_{clk} \ge t_{cq}+t_{comb,max}+t_{setup}+t_u.
\]

若 $T_{clk}=2.0$ ns，四项分别为 $0.12$、$1.62$、$0.08$ 和 $0.10$ ns，要求时刻与到达时刻之差只有 $0.08$ ns。这个正裕量不是“还有很多空间”，因为布线拥塞、工艺角和串扰模型更新都可能消耗它。报告阅读应先确认起点、终点、时钟关系和例外是否正确，再看组合单元、扇出、布线和逻辑级数；直接优化报告中最长的一串单元，可能只是在修一条约束错误的伪路径。

建立失败可以通过减少逻辑级数、复制高扇出驱动、调整布局或增加流水级修复。保持失败不能靠降低时钟频率解决，因为同一采样沿附近的数据仍然到得过早；常见修复是在数据路径增加受控延迟、调整时钟树或修正不合理的跨域约束。任何自动修复后都要重新检查所有角和所有模式，避免一个角的改动在另一个角制造新失败。

\section{八、CDC、复位与形式性质}

跨时钟域（Clock Domain Crossing, CDC）检查关注的不是组合功能，而是数据怎样跨越无固定相位关系的采样边界。单 bit 电平可以经两级同步器降低亚稳态传播概率；脉冲必须保证目标域能够看到，或改成 toggle/握手协议；多 bit 数据要用稳定快照握手、异步 FIFO 或带编码保证的专用结构。把多位总线的每一位分别接两级同步器，仍可能在目标域拼出源域从未存在过的混合值。

复位域跨越（Reset Domain Crossing, RDC）还要检查不同模块解除复位的先后关系。异步置位可以快速把输出拉到安全状态，但解除通常要在本地时钟域同步完成。若生产者已经开始发送而消费者仍处于复位，valid、credit 或指针状态会失配。设计应规定复位后谁先声明就绪，以及旧事务如何作废。

形式验证用“所有满足前提的状态和输入”检查性质，而不是只运行若干测试向量。适合 RTL 的性质通常分三类：安全性说明坏事永远不发生，例如 FIFO 不在空时读取；活性说明在公平前提下请求最终得到响应；等价性说明综合或优化前后的可观察行为一致。断言的前提必须描述环境契约，否则工具可能用现实中不可能的输入找到反例，也可能因为前提过强而把错误排除在搜索空间之外。

\begin{longtable}{L{0.18\linewidth}L{0.34\linewidth}L{0.38\linewidth}}
\toprule
验证层 & 主要输入 & 通过时真正证明了什么 \\
\midrule
lint/展开 & RTL 结构、参数和连接 & 未发现选定规则覆盖的静态结构缺陷 \\\tablerowrule
仿真 & 测试激励、参考模型、覆盖点 & 给定场景下行为符合预期，未覆盖状态仍未知 \\\tablerowrule
形式性质 & 假设、断言和可达状态空间 & 在假设成立的全部可达状态内性质成立 \\\tablerowrule
CDC/RDC & 时钟域、复位域和跨域结构 & 跨边界协议符合已识别的安全模板与时序要求 \\\tablerowrule
STA & 网表、延迟、时钟和接口约束 & 所有被分析路径在指定模式与工艺角满足时序 \\\tablerowrule
实现后检查 & 布局布线、功耗和物理规则 & 逻辑已映射到可制造或可配置资源并满足物理边界 \\
\bottomrule
\end{longtable}

\section{九、用证据闭合一次 RTL 修改}

一次看似局部的控制器修改可能同时改变状态可达性、组合路径、复位行为和接口背压。完整闭环应从需求中提取不变量与完成条件，先更新断言和测试，再运行 lint、仿真、CDC/RDC、综合与 STA。若修改增加一级流水寄存器，还要同步更新 valid、错误位、事务 ID 和输出延迟；只让数据多走一拍而控制仍按旧拍到达，会形成比原始时序失败更隐蔽的功能错误。

验收记录应保留工具版本、顶层参数、约束集合、未豁免告警、覆盖结果、最差建立/保持裕量以及关键资源变化。这样下一次频率下降、面积上升或跨域告警出现时，才能判断它来自功能改动、约束改动还是实现随机性。第 3 章建立的采样边界在这里成为可检查的时序路径；本章形成的网表与时序证据，则交给后续处理器章节作为“电路确实能按所述拍序运行”的实现基础。




























\section{十、从门级网表到可制造芯片}

综合完成时，设计仍然只是一张带时序约束的门级网表：其中已经出现具体的与非门、触发器、时钟门控单元和存储宏，却还没有决定这些单元位于硅片的什么位置，也没有形成可供制造的几何图形。把网表变成芯片，需要后端实现、签核、制造、封装和首次上电五个连续阶段。每个阶段都产生一种新的事实；后一个阶段不能用自己的通过结果替代前一个阶段的证明。

\begin{pdfdiagram}[x=1cm,y=1cm]
  \node[hw state, minimum width=2.5cm] (rtl) at (1.4,2.4) {RTL 与约束};
  \node[hw compute, minimum width=2.5cm] (syn) at (4.6,2.4) {综合与映射};
  \node[hw state, minimum width=2.5cm] (net) at (7.8,2.4) {门级网表};
  \node[hw compute, minimum width=2.5cm] (pnr) at (10.9,2.4) {布局布线};
  \draw[hw bus] (rtl) -- node[hw label,above]{逻辑/时序约束} (syn);
  \draw[hw bus] (syn) -- node[hw label,above]{库单元实例} (net);
  \draw[hw bus] (net) -- node[hw label,above]{位置与连线} (pnr);

  \node[hw control, minimum width=2.5cm] (sign) at (10.9,0.2) {签核与版图};
  \node[hw block, minimum width=2.5cm] (fab) at (7.8,0.2) {掩膜与晶圆制造};
  \node[hw block, minimum width=2.5cm] (pkg) at (4.6,0.2) {封装与量产测试};
  \node[hw state, minimum width=2.5cm] (bring) at (1.4,0.2) {上电与硅后验证};
  \draw[hw bus] (pnr) -- (10.9,1.3) -- (sign);
  \draw[hw bus] (sign) -- node[hw label,above]{GDS/掩膜数据} (fab);
  \draw[hw bus] (fab) -- node[hw label,above]{合格裸片} (pkg);
  \draw[hw bus] (pkg) -- node[hw label,above]{可供电器件} (bring);
\end{pdfdiagram}
\diagramnote{RTL 到可运行芯片的两行实现链。上行把行为约束落实为单元位置和物理连线；下行把通过签核的版图变成晶圆、封装器件和可上电样片。箭头上的产物是阶段间的交接契约，不能把“仿真通过”直接解释为“芯片可制造”。}

\topic{标准单元库把逻辑功能连接到器件与版图}

\term{标准单元}{standard cell}是预先完成晶体管级设计、版图和时序表征的基本块。一个二输入与非门通常同时提供多种驱动强度和阈值电压版本：高驱动版本能推动更大负载，但面积与动态功耗更高；低阈值版本切换更快，但漏电更大。综合器先选择逻辑功能，后端工具再根据扇出、线长、时序和功耗不断调整具体版本。存储器通常不由标准单元拼接，而以 SRAM、ROM 或寄存器文件宏的形式进入平面规划，因为宏的端口、宽高比、电源引脚和内部时序已经固定。

\term{平面规划}{floorplan}先确定芯片边界、I/O 区域、电源网格和大宏位置。宏若挡住主要数据通路，后续布线会产生绕行；宏若离使用者过远，访问延迟和动态功耗会上升；电源网格过弱则会在同时切换时产生 IR drop。平面规划不是美观问题，而是把面积、连线、电源和热密度提前放进同一张约束图。

\topic{布局、时钟树与布线把网表变成有延迟的几何结构}

\term{布局}{placement}为每个标准单元选择合法位置。工具在满足行高、阻塞区和拥塞约束的同时，尽量缩短关键网线并控制局部功耗密度。布局之后进行\term{时钟树综合}{Clock Tree Synthesis, CTS}：从时钟入口插入分级缓冲和必要的门控单元，使各寄存器看到的时钟沿满足偏斜、转换时间和负载限制。时钟树完成后，保持时间可能比布局阶段更紧，因为真实时钟到达差已经出现。

\term{布线}{routing}为电源、时钟和普通信号分配金属层、线宽、间距与过孔。局部低层金属连接相邻单元，高层宽线承担长距离、时钟和电源。布线产生的电阻与电容会反过来改变门延迟，因此工具必须抽取寄生参数，再用更新后的延迟重新执行静态时序分析。若一条路径失败，修复可能是换单元、插缓冲、移动逻辑、调整时钟或增加流水级；任何修复都要重新检查其他模式和工艺角。

\topic{签核回答“可制造且在规定边界内工作”}

\begin{longtable}{L{0.18\linewidth}L{0.34\linewidth}L{0.38\linewidth}}
\toprule
签核项目 & 检查的物理事实 & 失败意味着什么 \\
\midrule
STA & 所有时钟模式、PVT 角与例外路径 & 数据可能在采样边界前未稳定，或旧值保持时间不足 \\\tablerowrule
DRC & 线宽、间距、包围、密度等制造规则 & 几何图形超出工艺可稳定制造的窗口 \\\tablerowrule
LVS & 版图抽取网表与设计网表是否等价 & 版图中存在漏连、短接或器件参数不一致 \\\tablerowrule
功耗与 IR drop & 电源网格上的电流、压降和热点 & 单元在负载瞬态下可能得到不足电压 \\\tablerowrule
电迁移 & 金属与过孔的长期电流密度 & 连线可能在寿命期内因原子迁移失效 \\\tablerowrule
信号完整性 & 串扰、转换时间与噪声窗口 & 相邻活动可能改变关键网的延迟或逻辑判决 \\
\bottomrule
\end{longtable}

签核使用的是特定工艺设计套件、单元库、寄生模型和工作条件。通过签核表示“在这些模型和角落内没有发现违反约束”，不表示所有样片必然完全相同。制造波动、封装寄生、板级供电和未建模使用条件仍要由量产测试与硅后验证继续约束。

\topic{从掩膜到封装，逻辑状态变成可供电的器件}

签核版图经数据准备生成各层掩膜。晶圆制造反复执行薄膜沉积、光刻、刻蚀、离子注入和平坦化，在硅上形成晶体管与多层互连。制造完成后先做晶圆级探针测试，把明显失效的裸片排除；合格裸片切割、贴装并通过键合线、倒装焊凸点或硅中介层连接到封装基板。封装同时承担电源分配、信号逃逸、机械保护和热传导，封装寄生电感与热阻会进入真实工作边界。

首次上电应从最小可信条件开始：确认电源轨顺序与静态电流，确认复位与参考时钟，再尝试 JTAG 身份读取、片内 SRAM 测试、Boot ROM 取指和外部接口训练。若一上电就运行复杂软件，任何异常都可能同时来自电源、时钟、制造缺陷、封装连接、DFT 配置或逻辑设计，证据难以分层。

\section{十一、FPGA：位流怎样把通用阵列变成专用电路}

FPGA 不是“可以运行 Verilog 的处理器”。它是一张预制的可配置硬件网格，配置位决定查找表实现哪条真值关系、交换矩阵接通哪些导线、触发器是否启用、I/O 使用何种电平标准。RTL 综合与布局布线完成后生成的\term{位流}{bitstream}写入配置存储器，配置完成时，整张通用网格同时成为目标电路。

\begin{pdfdiagram}[x=1cm,y=1cm]
  \node[hw control, minimum width=2.4cm] (cfg) at (1.4,3.2) {配置端口\\位流装载};
  \node[hw state, minimum width=2.6cm] (bits) at (4.5,3.2) {配置 SRAM\\连接与功能位};
  \draw[hw bus] (cfg) -- (bits);

  \node[hw compute, minimum width=2.5cm, minimum height=1.2cm] (clb) at (2.0,1.0) {逻辑块\\LUT + FF + carry};
  \node[hw compute, minimum width=2.5cm, minimum height=1.2cm] (dsp) at (5.1,1.0) {DSP 阵列\\乘加与累加};
  \node[hw memory, minimum width=2.5cm, minimum height=1.2cm] (bram) at (8.2,1.0) {BRAM/URAM\\片上存储};
  \node[hw block, minimum width=2.5cm, minimum height=1.2cm] (io) at (11.3,1.0) {I/O + SerDes\\PLL/MMCM};
  \draw[hw bus] (clb) -- node[hw label,above]{可编程互连} (dsp);
  \draw[hw bus] (dsp) -- node[hw label,above]{可编程互连} (bram);
  \draw[hw bus] (bram) -- node[hw label,above]{可编程互连} (io);
  \draw[hw return] (bits.south) -- (4.5,2.25) -- (6.6,2.25);
  \node[hw label,anchor=west] at (6.7,2.25) {控制 LUT 内容、开关矩阵、I/O 与时钟资源};
\end{pdfdiagram}
\diagramnote{FPGA 的两层结构。上层配置存储器保存位流状态，下层是被这些状态重新连接的逻辑、DSP、存储、I/O 和时钟资源。位流不是程序指令流；它在配置边界上一次性决定后续每个时钟周期中并行存在的电路。}

\topic{LUT、触发器和进位链构成基本逻辑块}

$k$ 输入 LUT 本质上是一块 $2^k$ bit 的小存储器：输入作为地址，存储位作为输出，因此可以表达任意 $k$ 输入布尔函数。LUT 后接触发器，允许同一逻辑块输出组合结果或寄存结果。专用 carry chain 把相邻逻辑块的进位直接串联，避免宽加法器经过普通交换矩阵；这就是 FPGA 上加法器通常比任意同深度 LUT 网络更快的原因。

普通可编程互连需要配置开关和长线段，延迟常常高于 LUT 本身。FPGA 时序优化因而不只减少逻辑级数，还要减少远距离跨区、控制高扇出并使用专用时钟网络。BRAM、DSP、SerDes 和 PLL 是硬化资源：把常用结构直接做成固定电路，单位面积、速度和功耗显著优于用 LUT 拼接，但端口形状与数量受到器件限制。

\topic{配置启动是一条独立于用户逻辑的状态机}

上电后，FPGA 先保持用户 I/O 在安全状态，清除或初始化配置存储器，再从 SPI flash、并行 flash、JTAG 或主机接口接收位流。配置逻辑逐帧校验并写入内部配置 SRAM，全部完成后才释放全局启动信号，让用户时钟、触发器和 I/O 进入运行态。若位流校验失败、供电不稳定或配置时钟停止，用户逻辑不会获得一个“部分可用”的正常状态；平台应保持输出安全并报告配置错误。

位流还包含器件型号、资源位置和布线开关选择，通常不能跨器件直接复用。RTL 是相对可移植的功能描述，约束与位流则属于具体板卡和具体器件。调试时必须区分三类失败：RTL 行为错误、实现时序失败、配置与板级边界失败。

\section{十二、DFT 与量产测试：怎样观察芯片内部状态}

芯片正常工作时，数百万个触发器只通过少量外部引脚间接影响结果。制造测试若仍从功能接口施加输入，许多内部节点需要运行很长序列才能被激活和观察。\term{可测性设计}{Design for Test, DFT}增加专用测试路径，把“难以控制、难以观察”的内部状态转换成可移入和可移出的结构。

\begin{pdfdiagram}[x=1cm,y=1cm]
  \node[hw block, minimum width=2.4cm] (si) at (1.4,2.6) {scan in};
  \node[hw state, minimum width=2.2cm] (ff0) at (4.0,2.6) {扫描触发器 0};
  \node[hw state, minimum width=2.2cm] (ff1) at (6.8,2.6) {扫描触发器 1};
  \node[hw state, minimum width=2.2cm] (ff2) at (9.6,2.6) {扫描触发器 2};
  \node[hw block, minimum width=2.4cm] (so) at (12.2,2.6) {scan out};
  \draw[hw bus] (si) -- (ff0);
  \draw[hw bus] (ff0) -- (ff1);
  \draw[hw bus] (ff1) -- (ff2);
  \draw[hw bus] (ff2) -- (so);

  \node[hw compute, minimum width=2.6cm] (comb0) at (5.4,0.4) {功能组合逻辑 A};
  \node[hw compute, minimum width=2.6cm] (comb1) at (8.2,0.4) {功能组合逻辑 B};
  \draw[hw return] (ff0.south) -- (4.0,1.45) -- (comb0.west);
  \draw[hw return] (comb0.east) -- (6.8,1.45) -- (ff1.south);
  \draw[hw return] (ff1.south east) -- (7.1,1.2) -- (comb1.west);
  \draw[hw return] (comb1.east) -- (9.6,1.2) -- (ff2.south);
  \node[hw label] at (6.8,3.55) {scan\_enable=1：移位装载/读出；scan\_enable=0：按功能路径采样};
\end{pdfdiagram}
\diagramnote{同一组三只扫描触发器的两种连接视图。测试模式把触发器串成可移位链，先装入指定内部状态，再切回功能模式捕获一拍组合逻辑响应，最后切回扫描模式读出。两种模式共享存储元件，但具有不同的数据路径和完成边界。}

\topic{scan 与 ATPG 把内部组合逻辑变成可直接施测的对象}

扫描触发器在功能模式下与普通触发器相同；测试模式下，多只触发器首尾相连形成 scan chain。测试机先移入一组内部状态，切回功能时钟捕获一拍响应，再把结果移出。\term{自动测试向量生成}{Automatic Test Pattern Generation, ATPG}根据网表和故障模型寻找能激活并观察目标故障的向量。最常见的 stuck-at 模型假设节点固定为 0 或 1；transition 模型检查节点能否在规定时间内完成跳变。

大型芯片会并行使用多条扫描链并配合压缩器，缩短移位时间和外部引脚需求。扫描模式会改变时钟、功耗和时序条件，因此测试控制器必须规定移位频率、捕获脉冲、复位状态与模拟宏隔离，不能把测试模式视为普通功能模式的另一组输入。

\topic{存储器与高速接口需要专用自测}

SRAM 宏内部节点无法由普通 scan 直接覆盖，通常使用\term{存储器内建自测}{Memory Built-In Self-Test, MBIST}控制器执行 March 类序列：按正向和反向地址顺序进行写 0、读 0、写 1、读 1，暴露单元、耦合、译码和读写路径故障。发现坏行或坏列后，冗余分析逻辑选择 spare row/column，并把修复映射写入熔丝或一次性存储。

\term{逻辑内建自测}{Logic BIST, LBIST}常以 LFSR 产生伪随机激励，以多输入签名寄存器压缩响应。它适合上电自检和现场诊断，但压缩签名只能说明结果是否匹配，不能单独定位每个内部故障。SerDes、PLL、ADC 等模拟或混合信号宏还需要环回、眼图监测、校准码和专用模拟测试入口。

\topic{JTAG 边界扫描把芯片内部测试延伸到电路板}

JTAG TAP 由 TCK、TMS、TDI、TDO 和可选 TRST 构成状态机。边界扫描寄存器位于核心逻辑与封装引脚之间：正常模式透明传递信号，测试模式则可控制输出引脚并采样输入引脚。这样不运行 CPU、不依赖 DRAM，也能检查 BGA 封装下不可接触的焊点、相邻引脚短路和板级开路。JTAG 还可选择芯片 ID、旁路寄存器和片内调试模块，但量产产品必须用认证、生命周期状态或熔丝策略限制调试访问。

\topic{量产测试区分结构缺陷、速度分档与现场可靠性}

\begin{longtable}{L{0.18\linewidth}L{0.28\linewidth}L{0.42\linewidth}}
\toprule
阶段 & 主要对象 & 通过边界 \\
\midrule
晶圆探针测试 & 裸片电源、scan、MBIST、基础模拟宏 & 排除明显缺陷并记录可修复单元 \\\tablerowrule
封装后测试 & 封装连接、全速功能、温度与接口 & 确认封装寄生和装配未破坏功能 \\\tablerowrule
速度/功耗分档 & 最大频率、漏电、工作电压 & 把同一设计按可稳定工作范围分成产品等级 \\\tablerowrule
老化与抽检 & 高温高压、热循环和长期应力 & 估计早期失效率并验证寿命模型 \\\tablerowrule
系统上电自检 & 现场 SRAM、链路、时钟与传感器 & 确认当前机器满足继续启动的最低条件 \\
\bottomrule
\end{longtable}

量产测试通过不表示芯片在任意板卡和任意固件下都可运行；它证明的是器件在规定测试条件内满足结构、速度和功耗界限。系统启动仍要重新验证供电、时钟、存储训练和外部链路，因为这些边界直到芯片被装入具体机器后才形成。

\section*{章末练习}

\begin{chapterexercise}{综合练习：module 读图至位宽陷阱}
\begin{exercisesubproblem}{（a）module 读图}
阅读 \code{module m (input logic clk, en, rst, input logic [7:0] d, output logic [7:0] q); always\_ff @(posedge clk) begin if (rst) q <= 8'h00; else if (en) q <= d; end endmodule}。在纸上画出它对应的框图：标明每个元件的种类、位宽与全部连线，并指出哪部分电路对应 \code{if (rst)}、哪部分对应 \code{if (en)}
\exercisecheck{学习目标}{module 逐行读回硬件}
\end{exercisesubproblem}
\begin{exercisesubproblem}{（b）位宽陷阱}
\code{logic [3:0] a, b; logic c; assign c = (a + b) > 4'd15;} 想检测 4 位加法的进位。说明 \code{c} 的实际行为与原因，并给出两种正确写法
\exercisecheck{学习目标}{位宽由语境决定}
\end{exercisesubproblem}
\end{chapterexercise}

\begin{chapterexercise}{综合练习：锁存器消除至非阻塞交换}
\begin{exercisesubproblem}{（a）锁存器消除}
\code{always\_comb begin if (sel) y = a; end} 会推断出什么硬件？为什么会这样？给出两种消除写法，并说明两者综合出的电路
\exercisecheck{学习目标}{组合块每条路径都赋值}
\end{exercisesubproblem}
\begin{exercisesubproblem}{（b）非阻塞交换}
复位后 \code{a = 8'd1}、\code{b = 8'd2}。块 \code{always\_ff @(posedge clk) begin a <= b; b <= a; end} 运行两拍后 \code{a}、\code{b} 各是多少？若把 \code{<=} 全换成 \code{=} 呢？两种写法各自对应什么电路
\exercisecheck{学习目标}{非阻塞=先读旧值再同时提交}
\end{exercisesubproblem}
\end{chapterexercise}

\begin{chapterexercise}{两段式 FSM 改写}
“状态、时钟与时序逻辑”一章交通灯 FSM 曾把状态转移与灯输出全写在一个 \code{always\_ff} 里（一段式）。参照本章“时序逻辑的 Verilog 写法”一节的两段式模板，写出改写后两个块的骨架（注明每块的类型与内容），并说明灯输出现在属于哪类电路
\exercisecheck{学习目标}{状态寄存器与次态逻辑分家}
\end{chapterexercise}

\begin{chapterexercise}{testbench 自检写法}
为本章“testbench 与仿真阅读”一节的 \code{sap1\_ram} 写自检 testbench 的核心部分：向地址 7 写入 \code{8'hA5} 再读回，比对不符时用 \codetext{\$fatal} 报告，要求不依赖人眼读波形。列出生成时钟、驱动写读、自检三段各自的写法
\exercisecheck{学习目标}{预期+比对+报错才算自检}
\end{chapterexercise}

\begin{chapterexercise}{约束边界}
解释为什么只有主时钟约束仍不能证明芯片接口时序
\exercisecheck{检查要点}{输入/输出延迟属于板级契约}
\end{chapterexercise}

\begin{chapterexercise}{例外审查}
比较 false path 与 multicycle path 的功能含义
\exercisecheck{检查要点}{例外必须由真实采样关系证明}
\end{chapterexercise}

\begin{chapterexercise}{建立计算}
用正文给出的四项延迟计算 2.0 ns 周期下的裕量
\exercisecheck{检查要点}{到达与要求时刻使用同一单位}
\end{chapterexercise}

\begin{chapterexercise}{保持修复}
说明为什么降低频率不能修复普通保持失败
\exercisecheck{检查要点}{保持检查围绕同一采样沿}
\end{chapterexercise}

\begin{chapterexercise}{跨域选择}
为单 bit 电平、窄脉冲和连续多 bit 数据分别选择跨域结构
\exercisecheck{检查要点}{不逐位同步多 bit payload}
\end{chapterexercise}

\begin{chapterexercise}{验证闭环}
为控制器增加一级流水后列出必须同步更新和重新检查的对象
\exercisecheck{检查要点}{数据、控制、错误与事务身份同行}
\end{chapterexercise}

\section*{参考解答与要点}

\begin{chapteranswer}{综合练习：module 读图至位宽陷阱}
\begin{answersubproblem}{（a）module 读图}
框图：8 个 D 触发器并联成 8 位寄存器，时钟端共接 \code{clk}；每个触发器的 D 输入前串两级 2 选 1 选择器——最靠近 D 端、数据最后经过的一级由先写的 \code{rst} 控制，\code{rst=1} 时选常数 0；靠外的一级由 \code{en} 控制，\code{en=0} 时把 \code{q} 接回自身（保持），\code{en=1} 时放行 \code{d}。全部动作同步于时钟沿，无锁存器、无异步逻辑；端口 4 入 1 出，无隐藏状态
\answercheck{必须掌握的要点}{\code{always\_ff} 里每层 \code{if} 都是 D 端前的一级 mux，先写的条件优先级最高，控制最靠近数据输入端（最后选择）的那一级}
\end{answersubproblem}
\begin{answersubproblem}{（b）位宽陷阱}
\code{c} 恒为 0。比较 \code{>} 的两端位宽取语境最大值：左端 \code{a+b} 为 4 位、右端 \code{4'd15} 为 4 位，于是加法按 4 位计算，第 5 位进位在比较前已被丢弃，和最大 15，永远不大于 15。改正一：\code{assign \{c, s\} = a + b;}，5 位拼接作接收端，加法按 5 位算，进位自然进 \code{c}；改正二：先把加数显式扩展再比，\code{(\{1'b0, a\} + \{1'b0, b\}) > 5'd15}
\answercheck{必须掌握的要点}{表达式位宽取语境最大；进位只在接收端够宽时才存在}
\end{answersubproblem}
\end{chapteranswer}

\begin{chapteranswer}{综合练习：锁存器消除至非阻塞交换}
\begin{answersubproblem}{（a）锁存器消除}
\code{sel=0} 的路径上 \code{y} 没有被赋值，“保持旧值”在组合逻辑里没有免费实现，综合器只好推断出一级锁存器：\code{sel} 成了它的使能端。消除法一：块首加默认赋值 \code{y = 1'b0;}，让后面的 \code{if} 覆盖它；法二：补全分支 \code{else y = 1'b0;}。两种写法综合出的电路相同：一个由 \code{sel} 选择的 2 选 1 mux
\answercheck{必须掌握的要点}{锁存器来自“某条路径没赋值”，与 \code{always\_comb} 本身无关}
\end{answersubproblem}
\begin{answersubproblem}{（b）非阻塞交换}
非阻塞版：右端全部读旧值、时钟沿同时提交——一拍后 \code{a=2, b=1}，两拍后 \code{a=1, b=2}，每拍交换一次；对应电路是两只触发器 D 端交叉互连。阻塞版 \code{a = b; b = a;} 顺序生效：\code{a} 先得 2，\code{b} 再取到新的 \code{a} 也是 2——一拍后 \code{a=b=2}，两拍后仍是 2、2；对应电路是两只触发器的 D 端同接 \code{b}，交换被同化
\answercheck{必须掌握的要点}{\code{<=} 的语义就是触发器的并行装载；时序块里的交换必须用非阻塞}
\end{answersubproblem}
\end{chapteranswer}

\begin{chapteranswer}{两段式 FSM 改写}
骨架：\code{always\_ff @(posedge clk or negedge rst\_n)} 里只写 \code{if (!rst\_n) state <= S0; else state <= next\_state;}；\code{always\_comb} 里块首先给 \code{next\_state} 与各路灯输出默认赋值，再 \code{case (state)} 逐状态计算次态与输出。灯输出现在是纯组合电路（对状态的译码），与状态寄存器分家、不再占用触发器；一段式里它们是被时钟沿寄存的
\answercheck{必须掌握的要点}{时序块只装状态寄存器；次态与输出全部组合}
\end{chapteranswer}

\begin{chapteranswer}{testbench 自检写法}
时钟段沿用“testbench 与仿真阅读”一节骨架：\code{always \#5 clk = \~{}clk;}。驱动段写在 \code{initial} 里：\code{addr = 4'd7; din = 8'hA5; we = 1'b1;}，等一个上升沿后 \code{we = 1'b0;}。自检段：再隔一拍读回，\codetext{if (dout !== 8'hA5) \$fatal(1, "RAM FAIL"); else \$display("PASS"); \$finish;}。激励按拍推进，预期值 \code{8'hA5} 写进比对语句，严格不等 \code{!==} 让 X 也判失败
\answercheck{必须掌握的要点}{预期值+逐项比对+\codetext{\$fatal} 三件套；眼看波形不算自检}
\end{chapteranswer}

\begin{chapteranswer}{约束边界}
主时钟只给出内部采样周期；外部发送器的到达时间和接收器的要求时间还要由输入、输出延迟表达。
\answercheck{必须保留的检查点}{内部通过不等于板级接口通过}
\end{chapteranswer}

\begin{chapteranswer}{例外审查}
false path 表示功能上不发生同步采样，multicycle 表示允许若干周期后采样且需配套保持关系。
\answercheck{必须保留的检查点}{不能用例外隐藏普通慢路径}
\end{chapteranswer}

\begin{chapteranswer}{建立计算}
$2.0-(0.12+1.62+0.08+0.10)=0.08$ ns。裕量为正但很小，仍要检查角、模式和建模变化。
\answercheck{必须保留的检查点}{不遗漏不确定度预算}
\end{chapteranswer}

\begin{chapteranswer}{保持修复}
保持约束要求新数据不要在同一采样沿后过早到达，改变下一采样沿的周期并不改变这一局部关系。
\answercheck{必须保留的检查点}{用数据延迟或时钟树修复}
\end{chapteranswer}

\begin{chapteranswer}{跨域选择}
电平用同步器；窄脉冲用展宽、toggle 或握手；连续多 bit 流用异步 FIFO。
\answercheck{必须保留的检查点}{协议要定义接受与复位边界}
\end{chapteranswer}

\begin{chapteranswer}{验证闭环}
更新流水 valid、异常/错误位、ID、接口延迟与断言，重跑仿真、CDC、综合和 STA，并比较覆盖、资源与裕量。
\answercheck{必须保留的检查点}{延迟变化必须在接口契约中可见}
\end{chapteranswer}
